\medterm{Rabeprazole} A medicine that greatly reduces acid made by the stomach. It treats reflux and ulcers and may be combined with antibiotics for Helicobacter pylori infection; possible effects include headache, diarrhea, and rarely low magnesium or bowel infection.

\medterm{Rabies} A viral infection of the brain and nervous system caused by rabies virus, a virus that infects mammals. It spreads mainly when saliva from an infected animal enters a bite, scratch, broken skin, or the eyes, mouth, or another mucous membrane. After a variable symptom-free period, early symptoms may include fever, weakness, headache, or pain, itching, or unusual sensations near the exposure site. Later symptoms may include confusion, agitation, hallucinations, fear of water, difficulty swallowing, paralysis, coma, and death. Once symptoms begin, rabies is nearly always fatal, but vaccination before exposure or post-exposure treatment can prevent disease before symptoms start.

\medterm{Rabies Vaccination} Immunization used before or after exposure to reduce the risk of rabies. Post-exposure management
also includes wound care and, when indicated, rabies immunoglobulin.

\medterm{Rabies Virus} The virus that causes rabies, a usually fatal viral infection of the central nervous system transmitted most often through the saliva of an infected animal.

\medterm{Rachitic Rosary} A physical examination finding in pediatric rickets (vitamin D deficiency) characterized by prominent, beadlike, symmetrical nodular expansion of the costochondral junctions along the anterolateral rib cage resulting from disordered osteoid mineralization and chondrocyte hypertrophy.
\synonyms
Beaded Ribs; Costochondral Rosary

\medterm{Radial} Relating to the radius bone or the thumb side of the forearm and hand. The word can also describe structures arranged like spokes spreading from a central point.

\medterm{Radial Artery} An artery that runs along the thumb side of the forearm to the wrist and hand. Its pulse can be felt at the wrist, and it supplies blood to the forearm and parts of the hand.

\medterm{Radial Artery Thrombosis} A blood clot blocking the artery on the thumb side of the forearm, sometimes after an arterial catheter or injury. It may cause a painful, pale, cold, numb, or weak hand and can threaten hand circulation.

\medterm{Radial Club Hand} A birth-related difference in which the thumb-side forearm bone is small or absent and the hand bends toward the thumb side. The thumb may also be underdeveloped, and treatment depends on movement, function, and associated conditions.

\medterm{Radial Head Subluxation} A common elbow injury in young children, often called nursemaid’s elbow, caused when a sudden pull partly slips a forearm bone out of its supporting band. The child may stop using the arm but usually has little swelling. A trained clinician can put the bone back in place.

\synonyms
Nursemaid's Elbow; Pulled Elbow; Annular Ligament Displacement

\medterm{Radial Nerve} A major nerve of the arm that helps straighten the elbow, wrist, and fingers and carries feeling from part of the back of the hand. Injury can cause wrist drop, finger weakness, or numbness.

\medterm{Radial Nerve Dysfunction} Impaired function of the radial nerve causing weakness of wrist or finger extension and sensory changes over part of the posterior arm or hand; causes include compression, trauma, and systemic neuropathy.

\medterm{Radial Nerve Palsy} Dysfunction of the radial nerve causing weakness of wrist and finger extension, often producing wrist drop, with sensory loss over part of the back of the hand in some cases.

\medterm{Radial Tunnel Syndrome} Irritation or pressure on the radial nerve near the elbow, causing aching outer-forearm pain and sometimes weakness. Symptoms may worsen with repeated wrist or forearm movement and require evaluation to exclude other causes.

\medterm{Radial Vein} Veins running alongside the radial artery on the thumb side of the forearm. They drain blood from the forearm and join other deep veins near the elbow.

\medterm{Radiation} Energy traveling as particles or electromagnetic waves, including X-rays, gamma rays, ultraviolet light, and radio waves. Depending on its type and amount, radiation can diagnose disease, treat cancer, or damage living tissue.

\medterm{Radiation Biology} The study of how radiation affects cells, tissues, organisms, and populations. It examines injury, repair, mutations, cancer risk, treatment effects, and protective measures after exposure to ionizing or nonionizing radiation.

\medterm{Radiation Contamination} Radioactive material present on or inside a person, object, surface, or body site when it should not be. It can spread or continue exposing tissue and may require isolation, monitoring, removal of the material, and specialist advice.

\medterm{Radiation Dermatitis Skin Care} Care of skin affected by radiation treatment. The skin may become red, dry, sore, peeling, or moist, and care focuses on protecting the skin and treating irritation or infection when present.

\synonyms
Radiation Skin Care; Radiation Dermatitis Management; Post-Radiotherapy Skin Care

\medterm{Radiation Dosage} The amount of radiation given during an imaging test or radiation treatment. The dose is planned to obtain the needed benefit while limiting unnecessary exposure.

\medterm{Radiation Dose} The amount of ionizing radiation absorbed or received by a person or tissue. Dose varies with the
examination, equipment, protocol, and patient size; imaging should be justified and optimized
to obtain the needed information with the lowest reasonable exposure.

\medterm{Radiation Emergency} An event involving actual or suspected harmful radiation exposure or contamination that requires urgent protective and medical
measures.

\medterm{Radiation Enteritis} Inflammation or injury of the small or large intestine after radiation treatment. Early symptoms may include diarrhea, cramps, nausea, or urgency; months or years later, scarring can cause narrowing, poor absorption, bleeding, or blockage.

\medterm{Radiation Exposure} Contact with ionizing or non-ionizing radiation from medical imaging, treatment, sunlight, workplaces, or other sources. The possible harm depends mainly on the type, dose, body area, and duration of exposure.

\medterm{Radiation Fibrosis} Long-term thickening and stiffening of tissue caused by radiation injury, sometimes appearing months or years after treatment. It can affect skin, lungs, bowel, or other organs and may cause pain, reduced movement, or shortness of breath.

\medterm{Radiation Field} The area and distribution of radiation delivered during imaging or radiation treatment. In therapy, its size, shape, direction, and dose are planned to cover the target while limiting exposure to nearby healthy tissue.

\medterm{Radiation-Induced Angiosarcoma} A rare, fast-growing cancer of blood-vessel lining cells that develops years after radiation in previously treated tissue. It may appear as a bruise, purple patch, swelling, or lump and requires prompt specialist assessment.

\medterm{Radiation Leukemia} Radiation leukemia is leukemia that develops after significant exposure to ionizing radiation; it is a late treatment- or exposure-related complication.

\medterm{Radiation Monitoring} Measurement of radiation levels or individual exposure using devices such as personal dosimeters and area monitors.

\medterm{Radiation Necrosis} Permanent death of tissue caused by radiation injury, often developing months or years after treatment. In the brain it may cause headache, seizures, swelling, weakness, or speech and vision problems and can resemble returning cancer on scans.

\medterm{Radiation Oncologist} A doctor who plans and oversees radiation treatment for cancer and selected noncancerous conditions. They choose treatment areas and doses, coordinate care, monitor side effects, and adjust therapy during treatment.

\medterm{Radiation Oncology} A medical specialty that uses ionizing radiation to treat cancer. Treatment may aim to
destroy or control tumors, prevent cancer from returning, or relieve symptoms.

\medterm{Radiation Physics} The study of how radiation is produced, travels, interacts with tissues, and is measured. It supports medical imaging, radiation treatment, dose calculation, and protection from unnecessary exposure.

\medterm{Radiation Pneumonitis} Inflammation of lung tissue occurring weeks or months after radiation to the chest. It may cause dry cough, fever, breathlessness, or low oxygen and can later leave permanent scarring; new breathing symptoms need medical assessment.

\medterm{Radiation Poisoning} Illness caused by substantial exposure to ionizing radiation, potentially producing skin injury,
bone-marrow suppression, gastrointestinal symptoms, or organ failure.

\medterm{Radiation Pregnancy Screening} Assessment of pregnancy status before procedures involving ionizing radiation or
radioactive materials when pregnancy is possible. The decision to proceed requires
clinical justification, risk optimization, and consideration of gestational age and
fetal dose.

\medterm{Radiation Proctitis} Inflammation and injury of the rectum after pelvic radiation therapy, causing rectal pain, urgency, diarrhea, mucus, or bleeding.

\medterm{Radiation Protection} The practice of keeping radiation exposure as low as reasonably achievable while
obtaining the required diagnostic or therapeutic result. It relies on justification,
optimization, time reduction, distance, shielding, and appropriate contamination
control.

\medterm{Radiation Recall} Inflammation that returns in an area previously treated with radiation after a later trigger, often a medicine. The skin or lining may become red, sore, swollen, peel, or blister, and the treating team should review the suspected trigger.

\medterm{Radiation-Recall Dermatitis} An inflammatory skin reaction that occurs in a previously irradiated area after exposure to a
triggering substance, often months or years later. The area may become red, swollen, painful,
or blistered and requires clinical evaluation.

\medterm{Radiation Safety} Measures that limit unnecessary exposure to ionizing radiation while allowing needed imaging or treatment. They include using the shortest practical exposure time, keeping as much distance as possible, using shielding, preventing contamination, and following pregnancy-related precautions when relevant.

\medterm{Radiation Sickness} Acute radiation syndrome caused by high-dose whole-body or substantial partial-body
exposure to ionizing radiation. It may produce nausea, vomiting, marrow failure,
infection, bleeding, gastrointestinal injury, neurovascular collapse, and death
depending on dose.

\medterm{Radiation Therapy} Treatment that uses carefully aimed ionizing radiation to damage cancer cells so they stop dividing or die. It may be given from outside the body or from a source placed near or inside a tumor. Side effects depend on the body area and dose and may include tiredness or skin changes.

\medterm{Radiation Tolerance} The amount of radiation that a tissue or organ can receive before injury becomes likely. It varies with the tissue, dose, treatment schedule, and previous radiation.

\medterm{Radical Cystectomy} Surgery to remove the urinary bladder, usually for muscle-invasive or high-risk bladder cancer. Because urine can no longer leave through the bladder, another route for urine drainage is created.

\medterm{Radical Cystoprostatectomy} Radical cystoprostatectomy removes the urinary bladder, prostate, and nearby structures, usually to treat invasive bladder cancer in men.

\medterm{Radical Hysterectomy} An operation that removes the uterus and cervix together with surrounding parametrial tissues and the upper vagina, usually with pelvic lymph-node assessment or removal. It is used mainly for selected invasive cervical cancers and is more extensive than a simple or total hysterectomy.

\medterm{Radical Lymphadenectomy} Extensive surgery to remove lymph nodes and nearby tissue to look for or treat cancer spread. It may provide important staging information but can cause pain, nerve or vessel injury, and long-term swelling from lymph-fluid blockage.

\medterm{Radical Mastectomy} Radical mastectomy removes the breast, chest muscles, and many underarm lymph nodes; it is now rarely used because less extensive cancer surgery is often effective.

\medterm{Radical Nephrectomy} Radical nephrectomy removes an entire kidney and sometimes nearby tissues, usually to treat a kidney tumor.

\medterm{Radical Prostatectomy} Surgery to remove the whole prostate gland and usually the nearby seminal vesicles, most often to treat prostate cancer that has not spread widely. The urinary passage is rejoined to the bladder; possible effects include urine leakage, erection problems, bleeding, and infection.

\medterm{Radicular Cyst} An inflammatory odontogenic cyst arising at the apex of a nonvital tooth, usually from chronic pulpal infection and apical periodontitis.

\medterm{Radiculitis} Irritation or inflammation of a nerve root where it leaves the spine. It can cause pain spreading along an arm or leg, tingling, numbness, weakness, or changed reflexes because of a disk problem, narrowing, infection, or other disease.

\medterm{Radiculopathy} Dysfunction of a spinal nerve root, usually from disc herniation, narrowing, or inflammation. It may cause
shooting pain, tingling, numbness, weakness, or altered reflexes along the affected nerve pathway.

\medterm{Radioactive Iodine} A radioactive form of iodine taken up by thyroid cells. It is used to image the thyroid or destroy overactive thyroid tissue and selected thyroid cancer cells; radiation-safety precautions are needed after some treatments.

\medterm{Radioactive Iodine Ablation} Treatment using radioactive iodine to destroy remaining thyroid tissue or selected thyroid cancer cells after surgery. Thyroid cells absorb the iodine; the decision depends on cancer type, spread, recurrence risk, and the person’s overall health.

\medterm{Radioactive Iodine Uptake} A test measuring the percentage of administered radioactive iodine taken up by the thyroid over time, used to evaluate hyperthyroidism and selected thyroid disorders.

\medterm{Radioactive Iodine Uptake Test} A nuclear medicine test that measures how much radioactive iodine the thyroid gland absorbs. It helps evaluate hyperthyroidism and some thyroid disorders.

\medterm{Radioactivity} Energy released when unstable atoms change into more stable forms. The energy may be particles or rays and can be used in medical imaging or treatment, but enough exposure can damage living tissue.

\medterm{Radioallergosorbent Test} An in vitro radioimmunoassay that detects and quantifies allergen-specific IgE antibodies in a patient's blood serum, used to identify allergic sensitization to environmental aeroallergens, foods, medications, or venoms without risking in vivo anaphylaxis.

\synonyms
RAST

\medterm{Radiobiology} The study of how radiation affects living cells and tissues, including DNA damage, cell death, repair, mutations, treatment effects, and the development of cancer.

\medterm{Radiocarpal Joint} The main wrist joint between the forearm bone on the thumb side and several wrist bones. It allows the hand to bend forward, backward, and from side to side and may be injured by falls or arthritis.

\medterm{Radiochemical Purity} The proportion of a radioactive medicine present in the intended chemical form. Unwanted forms can make a scan less accurate, send radiation to the wrong tissues, or reduce treatment effectiveness.

\medterm{Radiochemistry} The chemistry of radioactive substances and compounds. In medicine it includes making and labeling radiopharmaceuticals, studying radioactive decay, separating products, and checking radiochemical purity and safety.

\medterm{Radiocontrast} A substance given during an imaging test to make selected tissues, blood vessels, or body spaces easier to see. Depending on the agent, allergic reactions, kidney injury, or leakage-related damage may occur.

\medterm{Radiodermatitis} Skin injury caused by ionizing radiation, such as radiation treatment. It may cause redness, dryness, itching, hair loss, peeling, sores, or long-term skin changes; severity depends on the dose, treatment area, and repeated exposure.

\medterm{Radioembolisation} Alternate name; see \textit{Radioembolization}.

\medterm{Radioembolization} A treatment that delivers tiny radioactive beads through an artery to selected liver tumors. The beads reduce the tumor’s blood supply and release radiation nearby; suitability depends on liver function, tumor extent, and blood flow.

\medterm{Radio-Femoral Delay} A clinically palpable lag or latency in the arrival of the femoral pulse compared to the simultaneous radial pulse, often accompanied by reduced femoral pulse amplitude. It is a classic clinical sign of coarctation of the aorta and extensive aortoiliac occlusive disease.

\medterm{Radiofrequency} A type of electromagnetic energy, between some radio waves and microwaves, used in communication and in selected medical tests and treatments. Medical uses include heating tissue or delivering energy during procedures.

\synonyms
Radiofrequency Energy; RF Radiation

\medterm{Radiofrequency Ablation} A procedure that uses high-frequency electrical energy to heat and destroy selected tissue, often under imaging guidance.

\medterm{Radiofrequency Ablation For Chronic Pain} A procedure that applies radiofrequency energy to selected sensory nerves or other targets to interrupt pain signaling for a limited period in carefully selected chronic pain conditions.

\medterm{Radiograph} An image made with X-rays. Dense structures such as bone appear light, while air appears dark and soft tissues appear in shades of gray. Radiographs are used to assess bones, the chest, teeth, and selected abdominal or foreign-body problems; they use a small amount of radiation.

\synonyms
X-Ray; Roentgenogram; Radiographic Image

\medterm{Radiographer} A trained professional who performs diagnostic imaging examinations or delivers radiotherapy, depending on specialty and jurisdiction.

\medterm{Radiography} Creation of images using X-rays. It is useful for bones, chest disease, dental structures, and many other conditions; radiation exposure is kept as low as reasonably achievable.

\medterm{Radioimmunoassay} A laboratory test that uses an antibody and a small radioactive label to measure a hormone, medicine, or protein in a sample. It can be very sensitive but is often replaced by newer nonradioactive tests.

\medterm{Radioimmunoguided Surgery} Surgery in which a radioactive tracer attached to an antibody helps the surgeon locate tissue carrying a particular marker, often tumor tissue. Special equipment detects the tracer during the operation.

\medterm{Radioimmunotherapy} Treatment that attaches a radioactive substance to an antibody so the antibody carries radiation toward cells with a particular marker, often cancer cells. Radiation can still affect nearby normal tissue and cause side effects.

\medterm{Radioiodine} A radioactive form of iodine used in nuclear medicine to image or treat thyroid tissue. It is taken up by thyroid cells; precautions are needed to limit radiation exposure to others.

\medterm{Radioiodine Therapy} Treatment using radioactive iodine, which is taken up mainly by thyroid cells to destroy overactive thyroid tissue or selected thyroid cancer cells. It requires individualized dosing and temporary radiation-safety precautions; pregnancy and breastfeeding are contraindications.

\medterm{Radioisotope} An unstable isotope that emits radiation as it decays; radioisotopes are used in medical imaging, diagnosis, and selected treatments.

\synonyms
Radioactive Isotope; Radiotracer

\medterm{Radioligand Therapy} A cancer treatment in which a radioactive particle is attached to a molecule that binds a target on cancer cells. The molecule delivers radiation to cells carrying that target; use depends on the cancer, target expression, organ function, and treatment-specific safety monitoring.

\medterm{Radiologic Health} The safe use of medical imaging and radiation-based treatment to protect patients, workers, and the public from unnecessary ionizing radiation while preserving clinical benefit. Safety measures include appropriate justification, shielding, and dose optimization.

\medterm{Radiologic-Pathologic Correlation} Comparing scan findings with tissue examined under a microscope to determine whether the sample explains the abnormal area. If they do not match, another sample or further testing may be needed.

\medterm{Radiologic Response} The change seen on medical scans after treatment, such as shrinkage, disappearance, stability, growth, or new lesions. Standard measurement rules help classify the response, but scans do not always show the full biological effect.

\medterm{Radiologic Technologist} A trained health professional who performs selected medical-imaging examinations and helps protect patients from unnecessary radiation exposure.

\medterm{Radiologist} A doctor trained to interpret medical images and sometimes perform image-guided procedures. Radiologists evaluate X-rays, ultrasound, CT, MRI, and other scans to diagnose disease, guide treatment, or obtain tissue samples.

\medterm{Radiology} The medical field that uses images to find and monitor disease and guide procedures. Radiologists interpret X-rays, ultrasound, CT, MRI, and other scans and may perform image-guided biopsies, drainages, injections, or treatments.

\medterm{Radiology-Patient Identification} Verification of the correct patient, examination, body site, indication, pregnancy
status when relevant, allergies, renal function, and prior imaging before image
acquisition or intervention. It is a core diagnostic and procedural safety step.

\medterm{Radiology Report} A written explanation of an imaging study. It describes why the study was done, the important findings, and their meaning; urgent or unexpected findings should be communicated promptly to the treating team.

\medterm{Radiolucent} Describing a material that lets more X-rays pass through and therefore appears darker on an X-ray or CT image. Air is highly radiolucent; some stones, foreign bodies, and abnormal tissues can also have this appearance.

\medterm{Radiomimetic} Radiomimetic describes a chemical or treatment that produces cellular or DNA damage resembling the effects of ionizing radiation.

\medterm{Radionuclide} An unstable form of an atom that releases radiation as it changes into a more stable form. In medicine, radionuclides can be attached to substances for imaging or used to deliver treatment to selected tissues.

\medterm{Radionuclide Angiography} An imaging test in which a small amount of radioactive tracer is used to follow blood flow through the heart or blood vessels. A camera detects the tracer and produces information about circulation or heart pumping.

\medterm{Radionuclide Cystogram} A nuclear medicine examination that tracks radiolabeled urine to detect vesicoureteral reflux or evaluate bladder emptying.

\medterm{Radionuclide Generator} A device that produces a short-lived radioactive substance from a longer-lasting radioactive substance. It supplies certain tracers for nuclear scans or treatment when needed and must be operated under strict safety and quality controls.

\medterm{Radionuclide Heart Scan} A nuclear medicine study that uses a radioactive tracer to evaluate myocardial blood flow, heart muscle viability, or ventricular function.

\medterm{Radionuclide Imaging} A scan that uses a small amount of radioactive tracer to show the structure or activity of an organ. A camera detects the tracer, helping evaluate blood flow, bone, thyroid, kidney, heart, or cancer-related changes.

\medterm{Radionuclide Imaging Or Scan} A scan that detects radiation from a small amount of radioactive tracer given to a person. The tracer collects in selected organs or tissues, allowing clinicians to examine blood flow, structure, or function and identify some diseases.

\synonyms
Scintigraphy; Nuclear Medicine Scan; Gamma Scan

\medterm{Radionuclides} Radioactive forms of atoms used in nuclear medicine. Attached to a suitable substance, they can show how an organ works on a scan or deliver radiation to selected diseased tissue during treatment.

\medterm{Radionuclide Therapy} Treatment that delivers radiation from inside the body using a radioactive substance directed toward selected tissue. It can treat some thyroid diseases, cancers, or painful bone lesions. The medicine, dose, precautions, and side effects depend on the target and the person’s kidney, blood, and overall health.

\medterm{Radionuclidic Purity} The proportion of radioactivity in a preparation that comes from the intended radioactive atom rather than unwanted radioactive contaminants. High purity helps ensure the expected scan or treatment and limits unnecessary radiation to other tissues.

\medterm{Radio-Opaque} Blocking X-rays and appearing relatively white on a radiograph; radiopaque is the modern spelling.

\medterm{Radiopaque} Describing a material that attenuates X-rays strongly and therefore appears relatively
white on a radiograph or CT image. Bone, metal, and some contrast agents are radiopaque.

\medterm{Radiopharmaceutical} A medicine containing a small amount of radioactive material. It can show how an organ or tissue works on a scan or deliver radiation to selected diseased cells; its distribution determines the information or treatment it provides.

\medterm{Radiopharmaceutical Extravasation} Leakage of a radioactive medicine from the vein into nearby tissue during an injection. It can reduce the accuracy of a scan and, with treatment doses, injure local tissue; staff assess and document the event.

\medterm{Radiopharmaceutical Labeling Kit} A prepared set of sterile ingredients used to attach a radioactive atom to a substance that travels to a particular organ or follows a body process. The finished product is tested before it is given for imaging.

\medterm{Radiopharmaceutical Quality Control} Testing performed to verify identity, radionuclidic purity, radiochemical purity, pH,
sterility, endotoxin limits, and other release specifications before clinical
administration. Requirements depend on the preparation and jurisdiction.

\medterm{Radiopharmacy} The preparation, testing, storage, and dispensing of medicines containing radioactive materials. Radiopharmacy requires accurate activity measurement, cleanliness, sterility checks, quality control, safe handling, and protection of workers and patients.

\medterm{Radioprotective Agent} A medicine or substance intended to reduce injury to healthy tissue from ionizing radiation. It may limit chemical damage or support cell repair, but its benefit and risks depend on the type, dose, and treatment setting.

\medterm{Radio-Radial Delay} A palpable asymmetry or latency between the right and left radial arterial pulses when examined simultaneously. It indicates proximal arterial obstruction affecting one subclavian artery, as seen in aortic dissection, subclavian steal syndrome, or thoracic aortic aneurysm.

\medterm{Radiosensitization} Making cells more vulnerable to radiation damage. Some medicines or treatments are combined with radiotherapy to increase damage to cancer cells, while treatment planning aims to limit harm to nearby healthy tissue.

\medterm{Radiosurgery} Radiosurgery delivers highly focused radiation to a small target, usually in the brain or body, without making a conventional surgical incision.

\medterm{Radiotherapy} Treatment using high-energy radiation to damage and control abnormal cells, especially cancer. It may be given from outside or inside the body and can cause temporary or lasting effects in nearby healthy tissue.

\medterm{Radiotherapy Bolus} A layer of material placed on the skin during selected radiation treatments so more radiation reaches the skin surface. It is positioned carefully to deliver the planned dose to the target area.

\medterm{Radiotherapy Dosage} Radiotherapy dosage is the amount of ionizing radiation prescribed for a treatment target, usually divided into carefully planned fractions.

\medterm{Radiotracer} A small amount of radioactive material attached to a substance that travels through a specific body process. A scanner detects where it goes, helping show blood flow, organ function, bone activity, or selected disease-related changes.

\medterm{Radioulnar Joint} A joint between the radius and ulna bones of the forearm. It allows the hand and forearm to turn palm-up and palm-down.

\medterm{Radium} A naturally radioactive chemical element that releases penetrating radiation as it breaks down. Internal exposure can damage bone and other tissues and increase cancer risk, so handling requires strict safety controls.

\medterm{Radium-223} A radioactive medicine that behaves like calcium and collects in areas of increased bone activity. It is used for selected prostate cancers that have spread to bone and may cause low blood counts or other side effects.

\medterm{Radius} The forearm bone on the thumb side. It joins the elbow and wrist and rotates around the ulna, allowing the palm to turn up or down.

\medterm{Radius Fracture} A break in the radius, one of the two bones of the forearm, often occurring near the wrist after a fall onto an outstretched hand.

\medterm{Radon} A naturally occurring radioactive gas that can accumulate indoors and increase lung-cancer risk when inhaled over time.

\medterm{Ragged Red Fibers} An appearance seen in a muscle biopsy when abnormal mitochondria build up around muscle fibers. It supports the diagnosis of some mitochondrial muscle diseases but is not specific to one condition. The biopsy result must be interpreted with symptoms, blood tests, genetic testing, and other findings.

\synonyms
RRF; Ragged-Red Muscle Fibers

\medterm{Raised Hemidiaphragm} A raised hemidiaphragm is elevation of one side of the diaphragm, caused by abdominal pressure, lung volume loss, phrenic-nerve dysfunction, or diaphragmatic disease.

\medterm{Rale} A rale, or crackle, is an abnormal discontinuous breath sound produced by fluid, collapsed airways, or opening of small airways during breathing.

\medterm{Rales} Discontinuous, adventitious non-musical respiratory lung sounds heard on auscultation, generated by the sudden explosive reopening of collapsed small airways and alveoli containing fluid or exudate; categorized into fine crackles (high-pitched, dry) and coarse crackles (low-pitched, wet).

\synonyms
Crackles; Crepitations

\medterm{Raloxifene} A medicine that acts like estrogen in bones but blocks some estrogen effects elsewhere. It may reduce spine fractures and selected breast-cancer risk after menopause, but can increase blood-clot risk.

\medterm{Raloxifene Hydrochloride} Raloxifene hydrochloride is a selective estrogen-receptor modulator used to prevent or treat osteoporosis and reduce breast-cancer risk in selected postmenopausal women.

\medterm{Raltegravir Potassium} Raltegravir potassium is an HIV integrase inhibitor used with other antiretroviral medicines to suppress HIV replication.

\medterm{Raltitrexed} A chemotherapy medicine used for some colorectal cancers. It blocks a cell process needed to make DNA, slowing the growth of cancer cells; possible effects include diarrhea, mouth sores, low blood counts, tiredness, and liver problems.

\medterm{Ramipril} Ramipril is an angiotensin-converting-enzyme inhibitor used for hypertension, heart failure, and cardiovascular risk reduction; cough, hyperkalaemia, renal-function changes, angioedema, and fetal toxicity are important risks.

\medterm{Rampant Caries} Severe tooth decay affecting many teeth or surfaces at once. It may be linked with frequent sugar, dry mouth, vomiting, poor oral hygiene, limited dental care, medicines, or other health problems.

\medterm{Ramsay Hunt Syndrome} Reactivation of the shingles virus near the facial nerve, causing severe ear pain, a blistering rash in or around the ear or mouth, and weakness on one side of the face. Hearing loss or dizziness may occur; early treatment can improve recovery.

\medterm{Ramsay Hunt Syndrome Type I} A rare neurological disorder characterized by progressive myoclonus epilepsy, intentional cerebellar tremor, progressive ataxia, and mild cognitive decline. Historically termed dyssynergia cerebellaris myoclonica; distinct from Ramsay Hunt syndrome type II (herpes zoster oticus).

\synonyms
Dyssynergia Cerebellaris Myoclonica; Dentatorubral Tremor Syndrome

\medterm{Ramsay Hunt Syndrome Type II} Reactivation of latent varicella-zoster virus within the geniculate ganglion of the facial nerve (cranial nerve VII), presenting with the clinical triad of acute ipsilateral peripheral facial paralysis, otalgia, and a painful vesicular eruption on the external auditory canal, auricle, or soft palate.

\synonyms
Herpes Zoster Oticus

\medterm{Ramucirumab} A monoclonal antibody that blocks vascular endothelial growth factor receptor 2 and is used in selected advanced cancers, including gastric, colorectal, lung, and liver cancers.

\medterm{Ramus} A ramus is a branch-like extension of a bone, nerve, artery, or other structure, such as the mandibular ramus.

\medterm{Random Allocation} Assignment of participants to study groups by chance to reduce selection bias and balance known and unknown factors.

\medterm{Randomised Controlled Trial} A study in which participants are assigned by chance to interventions or controls and outcomes are compared.

\medterm{Randomization Validation} Checking that a clinical trial's assignment process truly places participants into study groups by chance and is implemented as planned. This helps reduce selection bias and protects the reliability of comparisons between treatments.

\medterm{Randomized Controlled Trial} An experimental study in which participants are assigned by chance to different interventions or a control condition and followed for specified outcomes.

\medterm{Range Of Motion} The amount a joint can move in different directions. It may be measured by a healthcare professional or by the person moving alone; pain, swelling, stiffness, weakness, injury, or a shortened muscle can reduce it.

\medterm{Ranibizumab} An intravitreal monoclonal antibody fragment that inhibits vascular endothelial growth factor and treats neovascular macular degeneration, diabetic macular edema, retinal vein occlusion, and related retinal disease.

\medterm{Ranke Complex} The late, healed, fibrocalcified sequela of primary pulmonary tuberculosis, consisting of a calcified subpleural Ghon parenchymal focus and calcified ipsilateral draining hilar or mediastinal lymph nodes visible on chest radiography.
\synonyms
Calcified Primary Tuberculous Complex; Healed Ghon Complex

\medterm{Ranolazine} Ranolazine is an antianginal drug that inhibits the late sodium current in cardiac cells and reduces chronic angina without substantially lowering heart rate or blood pressure; QT prolongation and drug interactions require caution.

\medterm{Ranson Criteria} An older scoring system that estimates the severity and risk of death in acute pancreatitis using clinical and blood-test findings at admission and again after 48 hours. It is less commonly used now because newer scores and organ-failure assessment may be more practical.

\medterm{RANTES} RANTES, now called CCL5, is a chemokine that attracts immune cells and participates in inflammation, antiviral defense, and immune signaling.

\medterm{Ranula} A mucocele (mucus extravasation cyst) of the sublingual salivary gland, presenting as a
blue-translucent, fluctuant swelling in the floor of the mouth, typically lateral to the
lingual frenulum. A ranula results from trauma or obstruction of a sublingual gland
duct, causing mucus accumulation in the sublingual space.

\medterm{Raoultella} Raoultella is a genus of gram-negative bacteria that can cause opportunistic urinary, respiratory, wound, or bloodstream infections.

\medterm{Rape} Nonconsensual sexual penetration or sexual violence; it is a serious crime and a medical, safeguarding, and trauma-care concern.

\medterm{Raphe} A seam-like line or ridge where structures join, such as the median raphe or raphe nuclei of the brainstem.

\medterm{Raphe Nuclei} The raphe nuclei are brainstem groups of neurons that produce serotonin and influence mood, pain, sleep, arousal, and autonomic functions.

\medterm{Rapid Eye Movement Sleep} A sleep stage with rapid eye movements, vivid dreaming, and active brain patterns. Most skeletal muscles are temporarily relaxed during it, while breathing and heart rate can become more variable.
\synonyms
REM Sleep; Paradoxical Sleep

\medterm{Rapid Eye Movement Sleep Behavior Disorder} Alternate name; see \textit{REM sleep behavior disorder}.

\medterm{Rapid Gastric Emptying} Abnormally rapid movement of food from the stomach into the small intestine, sometimes called dumping syndrome.

\medterm{Rapid Heartbeat} Alternate name; see \textit{Tachycardia}.

\medterm{Rapidly Progressive Glomerulonephritis} Severe inflammation of the kidney’s filtering units that causes kidney function to worsen over days or weeks. It may cause blood or protein in urine, swelling, high blood pressure, or reduced urine.

\medterm{Rapid Plasma Reagin Test} A blood test that detects antibodies associated with syphilis. It is useful for screening and monitoring treatment, but false-positive and false-negative results occur, so reactive results need a confirmatory treponemal test.

\medterm{Rapid Sequence Induction} A method of giving anesthetic and muscle-relaxing medicines to place a breathing tube quickly while reducing the chance of vomit entering the lungs. It is used when urgent airway control is needed and requires preparation for difficult intubation.

\medterm{Rapid Sequence Intubation} An emergency airway management technique involving the near-simultaneous administration of a rapid-acting sedative induction agent (such as etomidate or propofol) and a neuromuscular blocking paralytic agent (succinylcholine or rocuronium) without intervening bag-mask ventilation, designed to minimize aspiration risk in an unfasted patient.

\synonyms
RSI

\medterm{Rapid-Sequence Intubation} An emergency method for placing a breathing tube quickly after giving medicines that induce unconsciousness and temporarily relax the muscles. It is used when airway protection is urgent, with oxygen preparation, backup equipment, blood-pressure monitoring, and confirmation that the tube is correctly positioned.

\medterm{Rapid Shallow Breathing} Breathing that is faster than normal but moves only small amounts of air each time. It may occur with lung disease, pain, anxiety, metabolic illness, or tiring breathing muscles and can be a warning sign of respiratory failure.

\medterm{Rapid Strep Test} A test performed on a throat swab to detect Group A Streptococcus. It provides a rapid
result, but a negative result may require confirmatory testing depending on the patient's age, symptoms, and local
guidance.

\medterm{Rare Disease} A condition affecting relatively few people in a defined population; many rare diseases are genetic, require specialized diagnosis, and may have limited evidence or treatment options.

\medterm{Rarefaction} A reduction in density or pressure, such as the low-pressure part of a sound wave. In medicine, understanding rarefaction helps explain ultrasound imaging, sound transmission, and pressure changes in tissues.

\medterm{Rasburicase} A recombinant urate-oxidase enzyme that converts uric acid to the more soluble allantoin. It is used
in selected patients at risk of or experiencing tumor-lysis-related hyperuricemia. It can
cause serious hypersensitivity or hemolysis in people with glucose-6-phosphate
dehydrogenase deficiency and requires appropriate screening and monitoring.

\medterm{Rash} A visible change in skin color, texture, or surface, such as spots, bumps, blisters, or scaling. Causes include infection, allergy, medicines, inflammation, and other disease; rapid or severe rashes need assessment.
\synonyms
Exanthem; Skin Eruption; Dermatosis

\medterm{Rasmussen Encephalitis} A rare, progressive, chronic immune-mediated neuro-inflammatory disorder occurring predominantly in children, characterized by unilateral hemispheric inflammation, drug-resistant focal motor seizures (frequently epilepsia partialis continua), progressive contralateral hemiparesis, cognitive deterioration, and progressive unilateral cortical atrophy on neuroimaging.

\synonyms
Chronic Focal Encephalitis; Rasmussen Syndrome

\medterm{Rasp} A roughened surgical instrument used to file, shape, or smooth bone or other hard tissue.

\medterm{Ras Peptide} A small protein fragment related to Ras proteins, which help control cell growth and division. Changes in Ras signaling can contribute to cancer.

\medterm{RAST} An older name for a blood test that measures IgE antibodies directed against specific allergens. Modern laboratories usually use newer specific-IgE immunoassays; results show sensitization and must be interpreted with symptoms and exposure history.

\synonyms
Radioallergosorbent Test; Specific IgE Test

\medterm{Rastelli Procedure} Open-heart surgery for certain complex birth defects involving the heart’s major arteries, a hole between the lower chambers, and a narrowed route to the lungs. A patch directs blood from the left ventricle to the aorta, and a tube with a valve connects the right ventricle to the lung arteries. The tube may need later replacement.

\synonyms
Rastelli Operation; Rastelli Repair

\medterm{Rat Bite Fever} A systemic bacterial infection caused by Streptobacillus moniliformis or Spirillum
minus, transmitted through rat bites or contact with rat-contaminated food or water. It
presents with relapsing fever, rash, polyarthritis, and lymphadenopathy. The
characteristic rash is petechial or maculopapular on the extremities.

\medterm{Rat-Bite Fever} Alternate name; see \textit{Rat Bite Fever}.

\medterm{Rathke Cleft Cyst} A benign, non-neoplastic, mucin-filled sellar or suprasellar epithelial cyst derived from remnants of the embryonic Rathke pouch, which may compress the optic chiasm (causing bitemporal hemianopia) or normal pituitary gland (producing hypopituitarism).
\synonyms
RCC; Pituitary Cleft Cyst

\medterm{Rathke Pouch} An embryonic ectodermal diverticulum arising from the roof of the stomodeum (primitive oral cavity) that ascends toward the developing brain to form the anterior pituitary gland (adenohypophysis: pars distalis, pars intermedia, pars tuberalis); incomplete obliteration gives rise to craniopharyngiomas or Rathke cleft cysts.

\medterm{Ratio} A ratio compares two quantities by division, such as a measurement of risk, concentration, or body proportions.

\medterm{Rationalization} A psychological defense in which a person creates plausible explanations for behavior or feelings while avoiding the underlying motive or conflict.

\medterm{Rat Lungworm Disease} Infection caused by the rat lungworm parasite Angiostrongylus cantonensis, usually after eating contaminated raw or undercooked produce or animals. It can cause eosinophilic meningitis with headache, neck stiffness, or neurologic symptoms.

\medterm{RATS} Alternate name; see \textit{Robotic-Assisted Thoracoscopic Surgery}.

\medterm{Ravuconazole} Ravuconazole is a triazole antifungal medicine studied and used in some settings for invasive fungal infections.

\medterm{Raw Egg} Raw eggs may contain Salmonella and are safest when properly refrigerated, handled hygienically, and cooked unless pasteurized products are used.

\medterm{Raynaud Phenomenon} Episodes in which small blood vessels in the fingers or toes suddenly narrow after cold or stress. The digits may turn white, blue, then red and become numb or painful; severe or new symptoms may indicate another disease.

\medterm{Raynaud's Disease} The primary form of Raynaud phenomenon, in which small blood vessels in the fingers or toes narrow episodically without an associated underlying disease. Cold or emotional stress may trigger attacks that cause numbness, tingling, or pain and a characteristic change in color from white to blue and then red as blood flow returns.

\medterm{Raynaud's Phenomenon} Episodic constriction of small arteries in the fingers or toes, causing color changes, numbness,
or pain after cold exposure or emotional stress.

\synonyms
Raynaud Disease; Raynaud Syndrome; Vasospastic Attack

\medterm{Raynaud’s Syndrome} Episodes in which fingers or toes change color and become cold, numb, or painful after cold exposure or stress because small blood vessels narrow. It may occur alone or with autoimmune disease.

\medterm{Raynaud Syndrome} Episodes in which small blood vessels in the fingers or toes narrow after cold or stress, causing numbness, pain, and color changes that may progress from white to blue to red. It may occur alone or result from another disease; severe attacks can damage tissue.

\medterm{Razor Bumps} Inflamed bumps caused by hairs curling back into the skin after shaving or other hair removal. They are more common with coarse or curly hair and may cause itching, pain, dark marks, or scarring; avoiding very close shaving and using proper shaving technique can reduce recurrence.

\medterm{RBC Count} A measurement of the number of red blood cells in a volume of blood, interpreted with hemoglobin, hematocrit, and red-cell indices to evaluate anemia or erythrocytosis.

\medterm{RBC Indices} Measurements describing the size, hemoglobin content, and variation of red blood cells. They help classify anemia and interpret a complete blood count alongside hemoglobin and other results.

\medterm{RBC Nuclear Scan} An imaging test using a small amount of the person’s blood labeled with a radioactive tracer. It can help locate ongoing bleeding in the digestive tract or assess blood flow through the heart.

\medterm{RBC Urine Test} A urine examination detecting or counting red blood cells, used to evaluate hematuria from urinary-tract infection, stones, glomerular disease, tumors, trauma, or other causes.

\medterm{RDA} Recommended dietary allowance, the daily nutrient intake sufficient for nearly all healthy people in a defined group.

\synonyms
Recommended Dietary Allowance; Recommended Daily Allowance

\medterm{Reabsorption} The transport of a substance from a body cavity or tubular fluid back into tissue or the bloodstream, such as renal tubular recovery of water and filtered solutes.

\medterm{Reaction} A response of the body or mind to a stimulus, illness, medicine, treatment, or event. Examples include an allergic reaction, an immune reaction, a treatment response, and a psychological reaction; the cause and severity determine its importance.

\medterm{Reactivation Tuberculosis} Active tuberculosis developing when a previously contained infection becomes active again, often years later. It commonly affects the upper lungs and may cause a long-lasting cough, coughing blood, fever, night sweats, weight loss, and spread to others through the air.

\medterm{Reactive Arthritis} Painful inflammation of joints that develops during or after certain intestinal or genital infections. It may also cause eye inflammation, skin changes, or urinary symptoms, and the triggering infection may already have cleared.

\synonyms
Reiter Syndrome; Post-Infectious Arthritis

\medterm{Reactive Attachment Disorder} A rare childhood condition in which a child persistently avoids or barely responds to comfort from caregivers. It is linked to severe neglect or unstable care early in life and may involve sadness, irritability, limited trust, and difficulty forming relationships.

\medterm{Reactive Gastropathy Pathology} Gastric mucosal injury characterized by foveolar hyperplasia, edema, smooth-muscle fibers in the lamina propria, vascular congestion, and relatively little active inflammation. Chemical injury, bile reflux, and medicines are common associations.

\medterm{Reactive Hypoglycemia} A fall in blood sugar occurring within a few hours after eating. It may cause shakiness, sweating, hunger, weakness, dizziness, or confusion; repeated episodes need assessment to confirm low sugar and identify the cause.

\medterm{Reactive Oxygen Species} Unstable oxygen-containing molecules produced during normal cell activity, infection, or inflammation. Small amounts can help cells signal and defend against germs, but excess amounts can damage cell membranes, proteins, and DNA.

\medterm{Reactive Perforating Collagenosis} A skin condition in which damaged collagen from deeper skin is pushed through the surface. It causes small, firm, very itchy bumps with a central plug, often on the arms, legs, or hands, and may be associated with diabetes or kidney disease.

\medterm{Reality Testing} The ability to recognize whether perceptions, thoughts, and beliefs match shared reality. Clinicians assess it during a mental-status examination; impaired reality testing can occur with psychosis, severe mood episodes, delirium, or some brain disorders.

\medterm{Rebound Effect} Worsening or return of symptoms after stopping or reducing a treatment that had been suppressing them.

\medterm{Rebound Headache} Alternate name; see \textit{Medication overuse headaches}.

\medterm{Rebound Insomnia} Return or worsening of insomnia after stopping a sleep medicine.

\medterm{Rebound Tenderness} A sharp increase in stomach pain when pressure on the belly is suddenly released, which signals that the inner lining of the abdomen is inflamed.

\medterm{Recalcitrant} Resistant to treatment, control, healing, or change; the term often describes a persistent disease or infection.

\medterm{Receiver Operating Characteristic} A graph showing how a diagnostic test’s ability to detect disease changes as its cutoff is changed. It compares true-positive and false-positive results; the area under the curve summarizes overall discrimination, not clinical usefulness by itself.

\medterm{Recent Memory} Recent memory is the ability to recall information or events from the near past and can be impaired by aging, delirium, brain disease, depression, or medicines.

\medterm{Receptor} A protein on or inside a cell that recognizes a specific signal, such as a hormone, nerve chemical, or medicine, and starts or changes a cell response.

\medterm{Receptor Binding} The attachment of a signaling molecule, medicine, or other ligand to a receptor on or inside a cell. Binding can activate, block, or change the receptor and its downstream cellular response.

\medterm{Receptor Desensitization} A reduced cell response after a receptor is stimulated continuously or repeatedly. The change can develop quickly and helps prevent excessive signaling, but it may also reduce the effect of a medicine or natural body signal.

\medterm{Receptor Downregulation} A decrease in the number or sensitivity of cell receptors after a signal stimulates them for a long time. The cell may remove receptors from its surface or break them down, becoming less responsive to the signal or medicine.

\medterm{Receptors} Proteins or structures on or inside cells that recognize specific signals and start responses such as growth, movement, or hormone action.

\medterm{Receptor Tyrosine Kinases} Cell-surface proteins that receive signals from hormones and growth factors. When switched on, they activate internal pathways controlling growth, survival, movement, and repair; abnormal activity can contribute to cancer and other disease.

\medterm{Receptor Upregulation} An increase in the number or sensitivity of cell receptors after prolonged blocking of a signal or reduced stimulation. This adjustment can make a cell respond more strongly when the signal or medicine is later present.

\medterm{Recession Of Medial Rectus Muscle} An eye operation that moves the medial rectus muscle backward to weaken its pull and correct certain inward eye deviations.

\medterm{Recessive} Describing a gene change whose effect usually appears only when a person has two altered copies, one inherited from each parent. A person with one copy may carry it without having the condition.

\medterm{RECIST} Standard rules for judging how a solid tumor changes on scans during treatment. Using measured lesions, they classify disease as completely gone, partly reduced, stable, or worsened; the categories support research and clinical decisions but do not replace overall assessment.

\medterm{Recognizing Medical Emergencies} Recognizing a medical emergency includes identifying severe breathing difficulty, chest pain, stroke signs, major bleeding, seizure, collapse, poisoning, or altered consciousness and calling emergency services.

\medterm{Recognizing Teen Depression} Recognizing depression in adolescents involves looking for persistent low mood or irritability, loss of interest, sleep or appetite changes, poor concentration, hopelessness, or self-harm thoughts.

\medterm{Recombinant} Produced by combining genetic material from different sources or by using engineered DNA in a cell or organism.

\medterm{Recombinant DNA} DNA made by joining genetic material from different sources. It is used in biotechnology to produce medicines such as insulin and growth hormone, develop some vaccines, and create vehicles for gene therapy.

\medterm{Recombinant Protein} A protein produced by putting its gene into a laboratory cell and growing the cell so it makes the protein. This method supplies some hormones, vaccines, antibodies, enzymes, and research materials.

\medterm{Recombinant Tissue Plasminogen Activator} A laboratory-made form of tissue plasminogen activator that helps convert plasminogen to plasmin and dissolve blood clots. It is used in selected emergencies, such as some ischemic strokes, heart attacks, or pulmonary emboli, when benefits outweigh bleeding risk.

\synonyms
rtPA; Alteplase; Recombinant tPA

\medterm{Recombinase} An enzyme that cuts, joins, or rearranges DNA segments. It helps cells exchange genetic material, repair some DNA damage, and assemble the gene parts needed to make a varied immune response.

\medterm{Recombinational Repair} A way cells repair a serious DNA break by using a matching, undamaged DNA copy as a guide. It restores genetic information accurately; inherited problems with this repair can increase cancer risk and treatment sensitivity.

\medterm{Recommended Daily Allowance} The Recommended Dietary Allowance is the average daily nutrient intake sufficient for nearly all healthy people in a specified age and sex group.

\medterm{Recommended Dietary Allowance} The average daily intake expected to meet the nutrient needs of nearly all healthy people in a defined group.

\synonyms
RDA

\medterm{Record} A health record is a documented account of a person’s history, findings, tests, diagnoses, treatments, decisions, and outcomes.

\medterm{Recovery Position} A side-lying position used for an unconscious but normally breathing person to help keep the airway open and reduce aspiration risk.

\medterm{Recreation Therapy} Recreation therapy uses adapted games, arts, exercise, and leisure activities to support physical, emotional, cognitive, or social functioning.

\medterm{Recrudescence} Recurrence of symptoms or infection after apparent improvement, often from persistence of the original organism or disease process.

\medterm{Rectal} Rectal means relating to the rectum, the final section of the large intestine before the anal canal.

\medterm{Rectal Biopsy} Removal of a small piece of tissue from the rectum for examination under a microscope. It can help investigate missing nerve cells, inflammation, infection, abnormal growth, or other rectal disease.

\medterm{Rectal Bleeding} Passage of bright-red or dark blood from the rectum, caused by hemorrhoids, fissures,
polyps, inflammatory bowel disease, diverticular disease, or cancer. Persistent, heavy,
or unexplained bleeding requires evaluation.

\synonyms
Hematochezia; Lower Gastrointestinal Bleeding

\medterm{Rectal Cancer} Cancer developing in the rectum, the final part of the large intestine. Symptoms may include bleeding, altered bowel habits, pain, or weight loss; treatment may combine surgery, radiation, chemotherapy, immunotherapy, or targeted therapy.

\synonyms
Rectal Carcinoma; Adenocarcinoma of the Rectum

\medterm{Rectal Carcinoma} Rectal carcinoma is cancer arising in the rectum and may cause bleeding, altered bowel habits, pain, or anemia.

\medterm{Rectal Culture} Laboratory culture of a rectal specimen to detect selected bacterial pathogens or colonization, depending on the clinical indication.

\medterm{Rectal Disease} A disorder affecting the rectum, including inflammation, infection, bleeding, obstruction, prolapse, functional disease, or neoplasm.

\medterm{Rectal Fistula} A rectal fistula is an abnormal tract connecting the rectum with skin or another organ, often causing drainage, infection, or passage of stool through an abnormal route.

\medterm{Rectal Hemorrhage} Bleeding from the rectum, which may result from hemorrhoids, fissure, inflammation, ischemia, polyps, cancer, trauma, or other gastrointestinal disease.

\medterm{Rectal Inflammation} Alternate name; see \textit{Proctitis}.

\medterm{Rectal Itching} Alternate name; see \textit{Anal itching}.

\medterm{Rectal Mucositis} Rectal mucositis is inflammation and injury of the rectal lining, often after radiation, chemotherapy, infection, inflammatory disease, or trauma.

\medterm{Rectal Necrosis} Rectal necrosis is death of rectal tissue from inadequate blood flow, severe inflammation, infection, pressure, or another destructive process.

\medterm{Rectal Neoplasm} An abnormal growth in the rectum that may be noncancerous, precancerous, or cancerous. It can cause bleeding, altered bowel habits, pain, anemia, or no symptoms and may require examination and biopsy.

\medterm{Rectal Obstruction} Impaired passage through the rectum caused by fecal impaction, stricture, tumor, foreign body, inflammation, or functional outlet disorder.

\medterm{Rectal Perforation} A tear or hole through the rectal wall that can release bowel contents and cause peritonitis, sepsis, or severe bleeding.

\medterm{Rectal Polyposis} Rectal polyposis is the presence of multiple polyps in the rectum, which may be sporadic or part of an inherited syndrome and can increase cancer risk.

\medterm{Rectal Polyps} Alternate name; see \textit{Colon Polyps}.

\medterm{Rectal Prolapse} The rectum slipping downward and partly or completely coming through the anus. It may cause a visible bulge, bleeding, mucus, discomfort, or difficulty controlling or passing stool and can result from weakened pelvic support, straining, or nerve problems.

\medterm{Rectal Prolapse Repair} A surgical or procedural treatment that restores a prolapsed rectum to its normal position and secures or removes affected tissue according to the prolapse and patient factors.

\medterm{Rectal Route} The rectal route administers a medicine or fluid through the anus into the rectum, where it may act locally or be absorbed systemically.

\medterm{Rectal Stenosis} Abnormal narrowing of the rectal lumen caused by inflammation, scarring, surgery, radiation, ischemia, trauma, or neoplasm.

\medterm{Rectal Tenesmus} Rectal tenesmus is a persistent or urgent feeling of needing to pass stool despite an empty or nearly empty rectum, caused by inflammation, infection, mass, or motility disorder.

\medterm{Rectal Ulcer} Alternate name; see \textit{Solitary rectal ulcer syndrome}.

\medterm{Rectocele} Bulging of the rectum into the back wall of the vagina because supporting tissues have weakened. It may cause pelvic pressure, difficulty emptying the bowel, or a vaginal bulge.

\medterm{Rectosigmoid Junction} The point where the sigmoid colon joins the rectum.

\medterm{Rectosigmoid Neoplasm} A rectosigmoid neoplasm is an abnormal growth where the sigmoid colon joins the rectum and may be benign or malignant.

\medterm{Rectouterine Pouch} The rectouterine pouch, or pouch of Douglas, is the peritoneal recess between the uterus and rectum where fluid or disease may collect.

\medterm{Rectovaginal Fistula} Abnormal connection between the rectum and vagina causing passage of stool or gas
through the vagina. Causes include obstetric injury particularly third and fourth degree
perineal tears, Crohn disease, radiation therapy, and surgical complications. Presents
with vaginal discharge of stool, flatus, or recurrent vaginal infections.

\medterm{Rectum} The final straight section of the large intestine, approximately 12-15 cm long,
extending from the sigmoid colon to the anal canal. The rectum stores feces before
defecation.

\medterm{Rectus Abdominis} The rectus abdominis is a paired abdominal muscle that flexes the trunk, stabilizes the pelvis, and helps compress abdominal contents.

\medterm{Rectus Abdominis Muscle} The paired vertical muscles at the front of the abdomen, commonly called the “six-pack.” They bend the trunk forward, stabilize the pelvis, and help increase pressure inside the abdomen; stretching can cause separation during pregnancy or after major weight change.

\medterm{Rectus Abdominis Sheath} A fibrous aponeurotic envelope formed by the abdominal wall muscles and surrounding the
rectus abdominis muscle. It contains the superior and inferior epigastric vessels and
related nerves.

\medterm{Rectus Sheath Hematoma} A collection of blood within the sheath surrounding the rectus abdominis muscle, often caused by muscle strain, coughing, trauma, anticoagulation, or vessel injury. It may cause localized abdominal pain and a tender mass.

\medterm{Recur} To recur means to return or happen again, as symptoms, disease, infection, or a tumor after a period of improvement or treatment.

\medterm{Recurrence} The return of a disease after it improved or seemed to disappear. Cancer may return at the original site, nearby lymph nodes, or another part of the body; recurrence does not always mean that treatment is no longer possible.

\medterm{Recurrence-Free Survival} The length of time after treatment during which a disease does not return. In cancer studies, it usually ends when cancer is found again or when a person dies, depending on the study rules. It measures control of disease, not a guarantee that treatment has cured it.

\medterm{Recurrence Risk} The chance that a disease or condition will return after treatment or occur again in a future pregnancy or family member. The estimate depends on the condition, previous history, test results, inheritance pattern, and other risk factors.

\medterm{Recurrent} Happening again after a period of improvement, remission, or apparent absence. A recurrent condition may return at the original site or elsewhere; examples include repeated urinary infections, repeated seizures, or cancer returning after treatment.

\medterm{Recurrent Breast Cancer} Breast cancer that returns after a period when tests found no detectable disease. It may recur near the original site or elsewhere and is assessed with examination, imaging, biopsy, and other tests.

\medterm{Recurrent Caries} Tooth decay that develops beneath or beside an existing filling, crown, or other restoration. It may result from leaking edges, plaque, frequent sugar, or dry mouth and can cause sensitivity, pain, or further tooth damage.

\medterm{Recurrent Disease} Return of a disease after a period of improvement, remission, or apparent absence, either at the original site or elsewhere.

\medterm{Recurrent Early Pregnancy Loss} Repeated miscarriages occurring early in pregnancy. Evaluation may consider uterine anatomy, chromosomal factors, antiphospholipid syndrome, endocrine or metabolic disease, parental genetics, and other causes, although no cause is found in many cases.

\medterm{Recurrent Fractures} Repeated bone fractures, suggesting persistent trauma, osteoporosis, bone fragility, metabolic disease, tumor, or an inherited skeletal disorder.

\medterm{Recurrent Laryngeal Nerve} A branch of the vagus nerve supplying most intrinsic muscles of the larynx; injury can cause hoarseness.

\medterm{Recurrent Laryngeal Nerve Palsy} A peripheral cranial neuropathy resulting from surgical injury (e.g., during thyroidectomy), trauma, or mediastinal tumor compression affecting the recurrent laryngeal branch of the vagus nerve (CN X), causing ipsilateral vocal cord immobility, hoarseness, and loss of airway protective reflexes.
\synonyms
RLN Palsy; Vocal Cord Paralysis

\medterm{Recurrent Oral Ulcers} Alternate name; see \textit{Canker sore}.

\medterm{Recurrent Pregnancy Loss} Two or more miscarriages. Possible causes include chromosome changes, an unusual shape of the uterus, hormone or thyroid problems, antiphospholipid syndrome, or other medical conditions, although no cause is found in many people. Evaluation is individualized and may identify treatment that improves the chance of a future healthy pregnancy.

\synonyms
Recurrent Miscarriage; Habitual Abortion; RPL

\medterm{Recurrent Prostate Cancer} Alternate name; see \textit{Prostate cancer recurrence}.

\medterm{Recurrent Respiratory Papillomatosis} A disease in which wart-like growths caused by human papillomavirus develop in the larynx or other breathing passages. The growths can cause hoarseness, noisy breathing, airway narrowing, or repeated recurrence.

\medterm{Recurrent Urinary Tract Infection} Repeated urinary tract infections separated by periods without infection. They may affect the bladder or kidneys and can be related to bacteria that persist, incomplete bladder emptying, stones, urinary blockage, or other urinary-tract problems.

\medterm{Red Birthmarks} Red birthmarks are vascular skin marks present at or soon after birth, ranging from harmless capillary marks to lesions requiring monitoring or specialist treatment.

\medterm{Red Blood Cell} A blood cell that carries oxygen from the lungs to the body's tissues. It contains hemoglobin, has a flattened, indented shape, and lacks a nucleus, which helps it bend through small blood vessels. Red blood cells also carry some carbon dioxide back to the lungs and help keep the blood's acid level balanced. They are made mainly in the bone marrow and normally live for about 120 days.

\medterm{Red Blood Cell Casts} Tube-shaped groups of red blood cells seen in urine. They form inside the kidney's small tubes and strongly suggest bleeding or inflammation in the kidney's filtering units, such as glomerulonephritis or vasculitis.

\synonyms
RBC Casts

\medterm{Red Blood Cell Labeling} A nuclear-medicine technique in which a person’s red blood cells are marked with a small amount of radioactive material. Scans can then show ongoing digestive-tract bleeding or measure blood flow and heart function.

\medterm{Red Bone Marrow} Bone marrow that makes red blood cells, white blood cells, and platelets. It is found mainly in certain bones and in the spongy parts of many bones.

\medterm{Red Cell Distribution Width} A standard automated complete blood count parameter that measures the degree of variation in red blood cell volume (anisocytosis); an elevated RDW aids in distinguishing iron deficiency anemia (high RDW) from uncomplicated thalassemia trait (normal RDW).
\synonyms
RDW; Erythrocyte Anisocytosis Index

\medterm{Red Cell Indices} Standardized automated complete blood count parameters derived from hemoglobin, hematocrit, and erythrocyte count that quantify red blood cell physical properties: mean corpuscular volume (MCV), mean corpuscular hemoglobin (MCH), and mean corpuscular hemoglobin concentration (MCHC).

\medterm{Red Eye} Redness of the white part of the eye caused by irritation, dryness, allergy, infection, injury, inflammation, or increased surface blood flow.

\medterm{Red Flag} A symptom, sign, or test result suggesting a potentially serious condition that requires
urgent assessment or escalation of care. Red flags are context-specific and should not
be used as a substitute for clinical judgment.

\medterm{Red Hepatization} The second pathological stage of acute lobar pneumonia (occurring 2 to 4 days after infection), in which the affected lung lobe becomes consolidated, heavy, and firm ('liver-like') due to massive filling of alveolar spaces with packed erythrocytes, neutrophils, and fibrin strands.
\synonyms
Second Stage of Lobar Pneumonia; Consolidation Stage

\medterm{Red Infarction} Tissue death caused by loss of blood flow in which blood also leaks into the damaged area, making it dark red. It is more likely in organs with a second blood supply or when a vein is blocked.

\synonyms
Hemorrhagic Infarction

\medterm{Red Man Syndrome} An adverse, non-IgE-mediated pseudoallergic hypersensitivity reaction caused by rapid intravenous infusion of vancomycin, in which direct degranulation of mast cells and basophils triggers massive histamine release, producing flushing, erythema, and pruritus of the face, neck, and upper torso.

\synonyms
Vancomycin Flushing Reaction

\medterm{Red-Man Syndrome} An acute, non-immunologic (anaphylactoid) rate-dependent adverse drug reaction caused by rapid intravenous infusion of vancomycin, mediated by direct mast cell and basophil degranulation with histamine release, presenting with erythematous flushing and pruritus of the face, neck, and upper torso.
\synonyms
Vancomycin Flushing Syndrome; Histaminoid Infusion Reaction

\medterm{Red Nucleus} A group of nerve cells in the midbrain that helps coordinate movement through connections with the cerebellum and spinal cord. Damage to nearby pathways can cause weakness, tremor, or abnormal movement.

\medterm{Reductase Deficiency} A disorder or laboratory finding caused by inadequate activity of a reductase enzyme, impairing a specific metabolic or hormone pathway.

\medterm{Reduction} Restoration of a displaced fracture or joint to better alignment, or decrease in amount, size, or activity.

\medterm{Reduction Division} The first division of meiosis, when a reproductive cell's chromosome number is reduced by half. It produces cells with one chromosome from each pair and helps create genetic variation in eggs and sperm.

\medterm{Reduction Mammaplasty} Surgery that removes breast tissue, fat, and excess skin to reduce breast size and relieve symptoms such as neck, shoulder, or back discomfort; it may also reshape the breasts.

\medterm{Redundant Mitral Valve} Excess leaflet tissue of the mitral valve, most commonly the posterior leaflet, which
may prolapse into the left atrium during systole. Often associated with myxomatous
degeneration and mitral valve prolapse. May be asymptomatic or cause mild mitral
regurgitation detected incidentally on echocardiography.

\medterm{Reduviidae} A family of insects including assassin bugs; some species transmit pathogens such as the parasite causing Chagas disease.

\medterm{Reed-Sternberg Cells} Large abnormal immune cells derived from B lymphocytes and commonly found in classic Hodgkin lymphoma. Their appearance in a tissue sample helps pathologists confirm that diagnosis.

\medterm{Reexpansion Pulmonary Edema} A sudden buildup of fluid in a lung after it rapidly expands from a collapsed state, often after drainage of a large or long-standing pneumothorax. It can cause severe breathlessness and low oxygen and needs urgent treatment.

\medterm{Refeeding Syndrome} Dangerous changes in body salts and fluid balance when food is restarted after severe or prolonged undernutrition. Low phosphate, potassium, or magnesium can cause weakness, swelling, abnormal heart rhythms, breathing failure, or death and requires careful medical monitoring.

\medterm{Reference Interval} The range of results usually found in a selected healthy reference group, often covering the middle 95 percent. A result outside it does not by itself prove disease and must be interpreted with symptoms, history, and the laboratory’s own method.

\medterm{Referral} A request for assessment or treatment by another healthcare professional or service. Referrals help patients access specialized expertise, tests, procedures, rehabilitation, or coordinated care for a particular concern.

\medterm{Referred Otalgia} Ear pain felt in an anatomically normal ear caused by sensory nerve sharing between the ear and distant head and neck structures via cranial nerves V3 (auriculotemporal nerve), VII (facial), IX (glossopharyngeal - Jacobson nerve), X (vagus - Arnold nerve), and upper cervical nerves C2-C3; mandates careful examination of the pharynx and larynx for occult malignancy.
\medterm{Referred Pain} Pain felt in an area away from the tissue causing it because nearby or connected nerves share pathways to the brain. Examples include shoulder pain from irritation below the diaphragm or arm and jaw pain during a heart attack.

\medterm{Reflex} A rapid, automatic response to a stimulus, such as pulling the hand away from heat. Reflex testing helps assess pathways involving the nerves, spinal cord, and brain.

\medterm{Reflex Action} A rapid, automatic response to a stimulus that occurs through a nerve pathway called a reflex arc. Reflexes protect the body, maintain posture, and provide information about nerve and spinal-cord function.

\medterm{Reflex Anoxic Seizure} A brief fainting episode that can look like a seizure after sudden pain, fright, or emotional stress. The heart slows briefly, the person becomes pale and unresponsive, and jerking may occur before rapid recovery.

\medterm{Reflex Arc} The nerve pathway that produces a rapid, automatic response. It includes a sensor, a sensory nerve, a processing center in the spinal cord or brain, a motor nerve, and the muscle or gland that responds.

\medterm{Reflex Epilepsy} Reflex epilepsy is epilepsy in which specific stimuli such as flashing light, sound, reading, or touch reliably trigger seizures.

\medterm{Reflexes} Automatic, rapid responses to stimuli carried through nerve pathways called reflex arcs. Testing them helps assess sensory nerves, motor nerves, the spinal cord, and parts of the brain.

\medterm{Reflux} The backward flow of a body fluid, especially stomach contents into the esophagus or urine from the bladder toward the kidneys. Symptoms and risks depend on the type, frequency, and underlying cause.

\medterm{Reflux Esophagitis Pathology} Mucosal injury associated with gastroesophageal reflux, often showing basal-zone hyperplasia, elongation of lamina-propria papillae, intercellular edema, and variable eosinophils or neutrophils.

\medterm{Reflux Laryngitis} Irritation or inflammation of the larynx attributed to reflux of stomach contents into the throat or upper airway. Hoarseness, throat clearing, cough, globus sensation, or throat discomfort may occur, although other causes of these symptoms should also be assessed.
\medterm{Reflux Nephropathy} Renal scarring and chronic kidney damage caused by vesicoureteral reflux, recurrent
urinary infection, or both. It may lead to hypertension, proteinuria, impaired growth,
and chronic kidney disease.

\medterm{Refraction} The bending of light as it passes through the cornea, aqueous humour, lens, and vitreous humour. In an eye examination, refraction determines the optical power needed to focus a clear image on the retina; refractive power is measured in dioptres.

\medterm{Refractive Disorder} A vision problem caused by the eye not focusing light correctly on the retina. Nearsightedness, farsightedness, astigmatism, and presbyopia are common refractive disorders and may be corrected with glasses, lenses, or surgery.

\medterm{Refractive Error} A refractive error occurs when the eye does not focus light correctly on the retina, causing nearsightedness, farsightedness, astigmatism, or presbyopia.

\medterm{Refractive Errors And Astigmatism} Refractive errors are vision disorders caused by the eye's inability to properly focus
light onto the retina, resulting in blurred vision.

\medterm{Refractive Surgery} Eye surgery that changes the shape of the cornea or another part of the eye to reduce short-sightedness, long-sightedness, or astigmatism. It may reduce dependence on glasses but can cause dry eyes, glare, infection, or imperfect vision.

\medterm{Refractometer} A refractometer measures how much light bends through a substance and is used for fluid concentration, optical properties, or eye refraction.

\medterm{Refractometry} Measurement of the eye’s refractive power to determine the optical correction needed for myopia, hyperopia, astigmatism, or presbyopia.

\medterm{Refractory} Resistant to a treatment or not responding adequately to it. The term may refer to cancer or other diseases.

\medterm{Refractory Anemia} An older term for anemia that does not improve with usual treatment. It may refer to a myelodysplastic syndrome or another cause, so evaluation looks for blood-cell disorders, bleeding, nutritional problems, kidney disease, and other conditions.

\medterm{Refractory Cancer} Refractory cancer is cancer that does not respond adequately to a treatment or begins progressing during treatment.

\medterm{Refractory Disease} Disease that does not improve adequately with appropriate standard treatment or that continues to worsen despite therapy. The term does not mean that every possible treatment has failed.

\medterm{Refractory Period} The brief recovery time after a nerve or muscle cell sends an electrical signal. During this time it cannot, or can only partly, send another signal, which helps control the speed and direction of nerve and heart activity.

\medterm{Refrigerant Poisoning} Poisoning from inhaling or contacting a refrigerant gas or liquid. Exposure can cause dizziness, abnormal heart rhythms, suffocation, frostbite, chemical burns, or lung injury.

\medterm{Refsum Disease} A rare inherited disorder in which a fatty substance builds up in the body. It may cause night blindness, loss of smell, nerve damage, poor coordination, hearing loss, and scaly skin.

\medterm{Regadenoson Pharmacological Stress} A medicine used during a heart scan to widen the heart’s arteries when a person cannot exercise adequately. It can cause flushing, headache, chest discomfort, shortness of breath, or abnormal heart rhythms.

\medterm{Regeneration} Growth or replacement of damaged or lost cells and tissues, helping restore structure or function after injury.

\medterm{Regimen} A planned schedule for medicines, treatments, diet, exercise, or other healthcare steps. It states what to do, how much or how often, and for how long, and may be changed according to response or side effects.

\medterm{Regional} Relating to a particular body area, geographic area, or localized distribution rather than the whole body.

\medterm{Regional Anesthesia} Numbing a larger body area by injecting local anesthetic near the nerves supplying it. The person is usually awake, although sedation may be added; the block can provide numbness during surgery and pain relief afterward.

\synonyms
Conduction Anesthesia; Nerve Block Anesthesia

\medterm{Regional Chemotherapy} Treatment that delivers anticancer medicine directly to a body region or organ to produce a high local concentration while reducing, but not eliminating, exposure elsewhere. It is suitable only for selected cancers and carries procedure-related and medicine-related risks.

\medterm{Regional Enteritis} An older name for Crohn disease, a chronic inflammatory bowel disease that can affect any part of the digestive tract.

\medterm{Regional Ileitis} An older name for Crohn disease affecting mainly the terminal ileum. It causes long-term inflammation that may lead to abdominal pain, diarrhea, weight loss, narrowing, fistulas, or nutritional problems.

\medterm{Regulation Of Health Professions} Legal and professional systems that govern healthcare workers' education, licensing, permitted duties, conduct, and accountability. Regulation aims to protect patients, maintain competence, support ethical care, and address unsafe practice.

\medterm{Regulatory Approval} Authorization by a competent regulatory authority for a medicine, device, or procedure to be marketed or used for specified purposes after review of evidence.

\medterm{Regulatory Authority} A regulatory authority is a government or official body that sets and enforces standards for medicines, devices, facilities, professionals, or public health.

\medterm{Regulatory Gene} A gene that controls when, where, or how strongly another gene is used. Changes in regulatory genes or their control regions can alter development, cell growth, metabolism, or disease risk.

\medterm{Regulatory T Cells} Immune cells that limit excessive immune reactions and help prevent attacks against the body’s own tissues. They develop mainly in the thymus or from other T cells in tissues. Problems with their function can contribute to autoimmune disease.

\medterm{Regurgitation} Backward flow of blood through a heart valve that does not close fully. Mild leakage may cause no symptoms; more severe leakage can enlarge or weaken the heart and cause breathlessness, tiredness, or swelling.

\medterm{Rehab} A structured rehabilitation program that helps restore function, independence, participation, or quality of life after illness or injury.

\medterm{Rehabilitation} A coordinated process that helps a person restore or adapt physical, cognitive, psychological, and social
function after illness, injury, surgery, or disability.

\medterm{Rehabilitation Clinic} An outpatient facility where a team helps people regain or adapt movement, communication, thinking, daily activities, or independence after illness, injury, surgery, or disability. Services may include physical, occupational, speech, or other rehabilitation therapies.

\medterm{Rehabilitation Engineering} The use of engineering and technology to improve independence after illness or disability. It includes artificial limbs, braces, wheelchairs, communication aids, home adaptations, and devices that help control the environment.

\medterm{Rehabilitation Motivation} A person’s willingness and ability to take part in rehabilitation and work toward functional goals. Pain, mood, fatigue, cognition, support, confidence, and the expected benefit can all influence participation.

\medterm{Rehabilitation Potential} The expected capacity to improve function, independence, participation, or quality of
life through rehabilitation. It depends on diagnosis, prognosis, cognition, motivation,
comorbidity, social support, and environmental resources.

\medterm{Rehabilitation Unit} A hospital unit for people who need coordinated medical care and intensive therapy after a serious illness or injury, such as stroke, brain or spinal-cord injury, or major surgery. The team works to restore safe movement, daily skills, communication, and independence.

\medterm{Rehydration} Replacement of water and electrolytes lost through dehydration, commonly using oral fluids or intravenous treatment when necessary.

\medterm{Reid Index} A microscope measurement comparing the thickness of mucus-making glands with the wall of a large airway. A higher value may be seen in chronic bronchitis, but it is not used alone to diagnose the condition.
\synonyms
Bronchial Gland-to-Wall Ratio; Reid Histopathologic Index

\medterm{Reifenstein Syndrome} An inherited disorder of androgen action, usually involving partial androgen insensitivity in an individual with a 46,XY chromosome pattern.

\medterm{Reiki} Reiki is a complementary practice involving light touch or hands near the body. Evidence that it treats disease is limited.

\medterm{Reinfection} A new episode of infection by the same or a related pathogen after the original
infection has resolved. Reinfection is distinct from relapse, in which the original
infection persists or returns from an incompletely eradicated focus.

\medterm{Reinke Edema} Chronic polypoid degeneration and diffuse, gelatinous mucoid edema within the superficial lamina propria (Reinke's space) of both true vocal cords, strongly associated with chronic cigarette smoking and vocal strain, presenting with progressive hoarseness and a deepened voice.
\synonyms
Polypoid Corditis; Diffuse Vocal Fold Edema

\medterm{Reinnervation} Restoration of nerve supply to a muscle or tissue after nerve repair, regeneration, or collateral nerve growth.

\medterm{Rejection} An immune response in which the body recognizes transplanted tissue or an organ as foreign and attacks it. Rejection can damage or destroy the transplant, may be sudden or gradual, and is monitored and treated with immune-suppressing medicines.
\synonyms
Allograft Rejection; Graft Rejection

\medterm{Relapse} Return of a disease or its signs and symptoms after a period of improvement or remission. The term is used for cancer, infections, mental-health and substance-use disorders, and other conditions; its precise meaning depends on the disease.

\medterm{Relapse-Free Survival} The time from completion of treatment to recurrence of the cancer or death from any
cause.

\medterm{Relapsing Fever} An acute febrile illness caused by Borrelia species, characterized by recurrent episodes
of fever separated by afebrile intervals.

\medterm{Relapsing Polychondritis} A rare immune disease in which episodes of inflammation damage cartilage. It commonly affects the ears, nose, eyes, joints, or windpipe and may cause pain, swelling, hearing or breathing problems; airway involvement can be life-threatening.

\medterm{Relapsing-Remitting Multiple Sclerosis} A form of multiple sclerosis in which neurological symptoms appear in attacks and then improve partly or completely for periods of time. Symptoms may include vision changes, numbness, weakness, balance problems, or fatigue. Treatment can reduce relapses and new nerve damage, although the course varies between people.

\synonyms
RRMS; Relapsing-Remitting MS

\medterm{Relative Afferent Pupillary Defect} Unequal response of the pupils because one eye or optic nerve sends less light information to the brain than the other. During a swinging-light test, both pupils may appear to dilate when the light moves to the affected eye; it can indicate retinal or optic-nerve disease.

\synonyms
RAPD; Marcus Gunn Pupil

\medterm{Relative Biological Effectiveness} A comparison of how strongly two types of radiation produce the same biological effect at the same absorbed dose. It depends on the radiation type, dose, tissue, and biological outcome.

\medterm{Relative Bradycardia} A physical finding in which the heart rate fails to increase appropriately in response to fever (sphygmothermic dissociation), where pulse rises by less than 10 beats per minute per degree Celsius of fever, characteristic of Salmonella typhoid fever, Legionnaires' disease, and yellow fever (Faget sign).
\synonyms
Sphygmothermic Dissociation; Inappropriate Bradycardia in Fever

\medterm{Relative Risk} A comparison of the chance of an outcome in an exposed group with its chance in an unexposed group. A value above 1 suggests higher risk, a value below 1 suggests lower risk, and 1 suggests no difference.

\medterm{Relaxation} A reduction in muscle tension, physiologic arousal, or psychological stress; techniques may include breathing, biofeedback, or progressive muscle relaxation.

\medterm{Relaxation Response} The body's calming changes after activities such as slow breathing, meditation, or progressive muscle relaxation. Heart rate, breathing, muscle tension, and stress-related alertness may decrease; these practices can support health but do not replace needed medical treatment.

\medterm{Relaxation Techniques For Stress} Relaxation techniques for stress include slow breathing, progressive muscle relaxation, meditation, guided imagery, and other methods that reduce physical arousal.

\medterm{Relaxation Therapy} Techniques such as breathing, muscle relaxation, or imagery used to reduce stress, arousal, or symptom burden.

\medterm{Relaxin} Relaxin is a reproductive hormone that helps remodel connective tissue and may contribute to cervical, pelvic, and vascular changes during pregnancy.

\medterm{Religion} A system of spiritual beliefs and practices that may influence health decisions, coping, diet, treatment preferences, and end-of-life care.

\medterm{Relocation Test} A shoulder examination often performed after an apprehension test. The examiner gently supports the front of the shoulder while the arm is in a vulnerable position. Relief of fear or pain may support shoulder instability, but the result must be interpreted with the rest of the examination.

\synonyms
Fowler Test; Jobe Relocation Test

\medterm{REM} Rapid eye movement, a sleep stage associated with vivid dreams and intense brain activity.
\synonyms
Rapid Eye Movement; REM Sleep

\medterm{REM Behavior Disorder} A sleep disorder where a person physically acts out dreams because normal muscle relaxation during REM sleep is lost.

\synonyms
RBD; REM Sleep Behavior Disorder

\medterm{Remedy} A treatment or measure intended to relieve a symptom or illness; safety and effectiveness depend on the specific remedy and condition.

\medterm{Remifentanil} A very short-acting opioid pain medicine given by infusion during anesthesia or intensive procedures. Its effects wear off quickly after the infusion stops, but it can slow breathing, lower blood pressure, and cause opioid-related complications.

\medterm{Remission} A period when the signs and symptoms of a disease decrease or disappear. Complete remission means no evidence of the disease can be found with available tests; partial remission means it has improved but remains detectable.

\medterm{Remittent Fever} Fever that rises and falls during the day but remains above normal temperature, rather than returning completely to normal.

\medterm{Remodeling} The body's ongoing process of removing old tissue and building new tissue to adapt to use, injury, growth, or disease. Bone remodeling responds to stress, hormones, nutrition, and medicines throughout life.

\medterm{Remote Surgery} Surgery performed when the surgeon is physically separated from the patient and controls robotic instruments through a telecommunications system. It may also include remote surgical assistance or guidance, with a qualified local team available to manage the patient and respond to emergencies.

\synonyms
Telesurgery
Telepresence Surgery

\medterm{Removal Technique} The method used to take away a foreign object, diseased tissue, stone, clot, or other material from the body. The appropriate method depends on its location, size, cause, and risk of bleeding, infection, or injury to nearby structures.

\medterm{REM Rebound} An increase in rapid-eye-movement sleep after sleep loss, REM deprivation, or stopping a medicine that suppresses REM sleep. It may cause vivid dreams, nightmares, or temporary changes in sleep patterns.

\medterm{REM Sleep} A sleep stage with rapid eye movements, active brain activity, and temporary relaxation of most skeletal muscles. Vivid dreams are common, and REM sleep supports memory, learning, emotional processing, and healthy development.

\medterm{REM Sleep Behavior Disorder} A sleep disorder in which the normal muscle paralysis of REM sleep is reduced or absent, allowing a person to move or vocalize while dreaming.

\synonyms
Rapid Eye Movement Sleep Behavior Disorder

\medterm{Renal} Renal means relating to the kidneys, which filter blood, regulate fluid and electrolytes, produce hormones, and excrete urine.

\medterm{Renal Abscess} A localized collection of pus within or around the kidney, usually caused by infection. Diagnosis and treatment may require imaging, antibiotics, and drainage.

\medterm{Renal Adenoma} A renal adenoma is a usually benign kidney tumor, often discovered incidentally; its size, imaging features, and growth determine whether follow-up is needed.

\medterm{Renal Agenesis} A condition present at birth in which one or both kidneys fail to form. One absent kidney may cause no symptoms if the other works well, but absence of both kidneys is usually fatal before or shortly after birth.

\medterm{Renal and Ureteral Injury} Mechanical, penetrating, blunt, or iatrogenic trauma affecting the kidneys or ureters. Clinical manifestations include gross or microscopic hematuria, flank pain, retroperitoneal hemorrhage, urinoma formation, or obstructive hydronephrosis requiring contrast-enhanced CT staging and potential surgical or endoscopic stenting.

\synonyms
Upper Urinary Tract Trauma; Renal Laceration; Ureteral Disruption

\medterm{Renal And Urological Diseases} Renal and urological diseases affect the kidneys, urinary tract, bladder, urethra, or male reproductive urinary organs.

\medterm{Renal Anemia} Anemia associated with chronic kidney disease. Diseased kidneys make less erythropoietin, a hormone that stimulates red-cell production, and iron deficiency or inflammation may add to the problem, causing tiredness and breathlessness.

\medterm{Renal Arteriography} Renal arteriography uses contrast X-rays to visualize kidney arteries and may identify narrowing, bleeding, aneurysm, or vascular anatomy during an intervention.

\medterm{Renal Artery} One of the two main arteries carrying blood from the abdominal aorta to the kidneys. Each divides into smaller branches inside the kidney so blood reaches the filtering tissue.

\medterm{Renal Artery Aneurysm} A focal enlargement of a renal artery, often asymptomatic but capable of causing hypertension,
thromboembolism, or rupture. Management depends on size, symptoms, anatomy, pregnancy, and
growth, and may involve surveillance, endovascular treatment, or surgery.

\medterm{Renal Artery Disease} Alternate name; see \textit{Renal artery stenosis}.

\medterm{Renal Artery Doppler} Doppler ultrasound that measures blood flow in the renal arteries. It may help assess narrowing or blockage that can contribute to renovascular hypertension.

\medterm{Renal Artery Embolism} Sudden blockage of an artery supplying the kidney, often by a clot from the heart in atrial fibrillation or another heart condition. It may cause abrupt side pain, blood in urine, nausea, or kidney injury.

\medterm{Renal Artery Stenosis} Narrowing of an artery supplying a kidney. It can reduce kidney blood flow and cause difficult-to-control high blood pressure or worsening kidney function. Causes include fatty buildup and fibromuscular dysplasia.

\synonyms
RAS; Renal Artery Narrowing; Renovascular Disease

\medterm{Renal Biopsy} Alternate name; see \textit{Kidney biopsy}. It is used selectively for unexplained acute kidney injury, nephrotic syndrome, glomerular hematuria, or transplant dysfunction; bleeding is the main complication.

\medterm{Renal Calculi} Alternate medical term; see \textit{Nephrolithiasis}.

\medterm{Renal Calculus} A kidney stone, a hard mass of mineral crystals formed in the urinary tract. It may pass without symptoms or cause severe wave-like pain, blood in the urine, blockage, vomiting, or urinary infection.

\medterm{Renal Cancer} Alternate name; see \textit{Kidney cancer}.

\medterm{Renal Carcinoma} Renal carcinoma is cancer arising from kidney tissue, most commonly renal cell carcinoma. Cancer arising from kidney tissue, most often from cells lining the kidney's small filtering tubes.

\medterm{Renal Cell Carcinoma} A cancer beginning in the kidney’s small filtering tubes, most often the clear-cell type. It may cause blood in urine, pain, weight loss, fever, or no symptoms until found on imaging.

\medterm{Renal Cell Carcinoma Pathology} Pathologic classification of renal cell carcinoma by growth pattern, cytology, immunophenotype, necrosis, sarcomatoid or rhabdoid change, and invasion. Clear cell, papillary, chromophobe, and other recognized subtypes have different features and behavior.

\medterm{Renal Changes In Pregnancy} Pregnancy increases renal blood flow and glomerular filtration. These changes lower
usual serum creatinine values and may increase urinary loss of glucose and amino acids.

\medterm{Renal Circulation} Blood flow through the kidneys, which supports filtration, oxygen delivery, hormone production, and regulation of fluid and electrolytes.

\medterm{Renal Clearance} A measure of how quickly the kidneys remove a substance from the blood. It is usually calculated from blood and urine levels and helps assess kidney function or adjust some medicine doses.

\medterm{Renal Colic} Severe, wave-like pain caused when a kidney stone or other blockage stretches the urine tube. Pain usually starts in the side and may move toward the groin, with nausea, vomiting, or blood in urine; fever with blockage is an emergency.

\medterm{Renal Corpuscle} The first part of a kidney filtering unit, made of a tiny ball of blood vessels inside a cup-shaped capsule. Water and small dissolved substances pass out of the blood here to begin urine formation.

\medterm{Renal Cortex} The outer layer of the kidney. It contains many of the structures that filter blood and process the first fluid that will become urine, while extensions of the cortex separate the inner kidney pyramids.

\medterm{Renal Cryptococcosis} Infection of the kidney by Cryptococcus species, usually as part of disseminated cryptococcosis in
an immunocompromised person.

\medterm{Renal Cyst} Alternate name; see \textit{Kidney Cysts}.

\medterm{Renal Dialysis} Alternate term for dialysis: kidney replacement therapy using hemodialysis or peritoneal dialysis to remove excess fluid and dissolved waste products when kidney function is inadequate.

\medterm{Renal Diet} An eating plan adjusted for kidney disease. Depending on kidney function and treatment, it may limit sodium, potassium, phosphorus, or fluid while providing enough calories and protein; needs differ between people and during dialysis.

\medterm{Renal Disease} Any condition affecting the kidneys, including infection, inflammation, stones, obstruction, inherited disease, or chronic loss of filtering function. Kidney disease can alter fluid balance, blood pressure, and body chemistry.

\medterm{Renal Diseases} Disorders affecting kidney tissue or function, including glomerular, tubular, vascular, infectious, obstructive, and neoplastic disease.

\medterm{Renal Dose Adjustment} Modification of a medication dose, dosing interval, or both to account for reduced
kidney clearance. The adjustment should use an appropriate estimate of renal function,
drug-specific pharmacokinetics, treatment indication, and therapeutic monitoring when
available.

\medterm{Renal Dysplasia} Abnormal development of kidney tissue before birth. It may affect one or both kidneys and can cause reduced kidney function, high blood pressure, urine reflux, or other urinary-tract problems.

\medterm{Renal Dysplasia And Cystic Disease} A group of developmental or inherited kidney disorders in which the kidneys have abnormal tissue, cysts, or both. They may affect one or both kidneys and can reduce kidney function or cause high blood pressure.

\medterm{Renal Failure} Severe loss of kidney function, causing waste, extra fluid, and abnormal body salts to build up. It may develop suddenly and sometimes improve, or progress gradually as chronic kidney disease; dialysis or transplant may be needed.

\medterm{Renal Function} The kidneys’ ability to filter blood, remove waste, regulate fluids and electrolytes, maintain acid-base balance, and produce or activate hormones.

\medterm{Renal Glomerulus} A tiny ball of blood vessels inside each kidney filtering unit. Water and small dissolved substances pass from the blood through its walls to begin urine formation, while blood cells and most proteins remain in the bloodstream.

\medterm{Renal Hamartoma} A renal hamartoma is a benign disorganized overgrowth of mature tissues in the kidney, sometimes used to describe an angiomyolipoma-like lesion.

\medterm{Renal Hypertension} High blood pressure caused by kidney disease or narrowing of the renal arteries. High blood pressure caused by kidney disease or narrowed renal arteries, often requiring treatment of both causes.

\medterm{Renal Hypertrophy} Renal hypertrophy is enlargement of kidney tissue, often compensating for loss or reduced function of the other kidney but sometimes caused by disease.

\medterm{Renal Hypoperfusion} Reduced blood flow to the kidneys, causing decreased glomerular filtration and potentially acute kidney injury. Causes include dehydration, blood loss, shock, heart failure, renal artery disease, and severe hypotension. Management requires treating the cause and restoring adequate circulation while avoiding further kidney injury.

\medterm{Renal Infarction} Death of part of the kidney caused by a blocked blood vessel. It may cause sudden flank or abdominal pain, nausea, blood in the urine, or reduced kidney function.

\medterm{Renal Insufficiency} Renal insufficiency is reduced kidney function that impairs waste removal, fluid and electrolyte balance, acid handling, or hormone production.

\medterm{Renal Leukemia Infiltration} Spread of leukemia cells into kidney tissue. It may enlarge both kidneys and reduce kidney function, sometimes causing sudden kidney injury, blood or protein in urine, or high blood pressure; treatment focuses on the underlying leukemia.

\medterm{Renal Lithiasis} Alternate name; see \textit{Kidney stones}.

\medterm{Renal Manifestations of Sickle Cell Disease} Kidney complications of sickling disorders, including excessive urination, impaired concentration, hematuria, proteinuria, papillary necrosis, acute kidney injury, and progressive chronic kidney disease. Blood pressure and urine abnormalities require ongoing monitoring.

\medterm{Renal Mass} An abnormal area seen in or on a kidney during imaging. It may be a harmless cyst, infection, scar, benign growth, or cancer; its appearance, size, growth, and sometimes biopsy determine the next step.

\medterm{Renal Medulla} The inner part of the kidney containing tubes and blood vessels that help concentrate urine. It receives the fluid processed by the outer kidney and directs the finished urine toward the kidney’s central collecting area.

\medterm{Renal Necrosis} Renal necrosis is death of kidney tissue caused by severe ischemia, toxins, infection, obstruction, or other injury and can cause acute kidney failure.

\medterm{Renal Oncocytoma} A renal oncocytoma is a usually benign kidney tumor arising from oncocytic cells; imaging may not always distinguish it reliably from renal cancer.

\medterm{Renal Osteodystrophy} Bone and mineral changes caused by advanced chronic kidney disease. Abnormal calcium, phosphate, vitamin D, and parathyroid hormone levels can weaken or deform bones, cause bone pain, and increase fracture risk; treatment corrects these imbalances.

\medterm{Renal Papillary Necrosis} Renal papillary necrosis is death of the kidney papillae from ischemia, infection, obstruction, diabetes, sickle-cell disease, or analgesic injury and may cause blood or obstruction.

\medterm{Renal Parenchyma} The working tissue of the kidney, containing the tiny filtering units and supporting structures that remove waste from blood, balance body fluids, and produce urine.

\medterm{Renal Pelvis} The funnel-shaped expansion of the ureter within the kidney, formed by the union of
major calyces. It collects urine from the renal papillae and directs it into the ureter.

\medterm{Renal Pelvis Or Ureter Cancer} Cancer of the renal pelvis or ureter usually arises from urothelial cells lining the urinary tract and may cause blood in the urine.

\medterm{Renal Perfusion Scintiscan} A nuclear medicine study using a radiotracer to evaluate renal blood flow and relative perfusion of the kidneys.

\medterm{Renal Pyramids} Cone-shaped structures in the renal medulla that contain tubules and collecting ducts and drain urine toward the renal papillae.

\medterm{Renal Reabsorption} The return of water and useful substances, such as salt, sugar, and bicarbonate, from kidney tubes into the blood. This prevents their unnecessary loss and helps control body fluids and chemistry.

\medterm{Renal Replacement Therapy} Treatment that takes over some functions of severely damaged kidneys, including removing waste and extra fluid and correcting dangerous chemical imbalances. Options include hemodialysis, peritoneal dialysis, continuous treatment in intensive care, or kidney transplantation.

\medterm{Renal Replacement Therapy Indications} Reasons to consider urgent dialysis include dangerously high potassium, severe acid buildup, certain poisonings, fluid overload causing breathing difficulty, or symptoms from waste buildup. The decision also considers kidney recovery, overall health, and the person’s goals.

\medterm{Renal Salt-Wasting Syndrome} Excessive loss of sodium and other salts through the kidneys, causing low blood sodium, dehydration, low blood pressure, or muscle symptoms. Causes include inherited tubular disorders, medicines, kidney disease, and hormonal problems.

\medterm{Renal Sarcoma} Renal sarcoma is a rare malignant tumor arising from connective or supporting tissues of the kidney.

\medterm{Renal Scan} A nuclear medicine test using a small radioactive tracer to assess kidney blood flow, drainage, or separate function. The tracer and scan method are chosen according to the clinical question.

\medterm{Renal Scintigraphy} Nuclear imaging that follows a radioactive tracer through the kidneys. It can compare kidney function, assess blood flow and urine drainage, detect obstruction, or monitor a transplanted kidney.

\medterm{Renal Stone} A renal stone is a solid crystal aggregate formed in the kidney that can pass through the urinary tract or obstruct urine flow.

\medterm{Renal Transitional Cell Carcinoma} A urothelial carcinoma arising in the renal pelvis or calyces rather than the kidney’s filtering tissue. It may cause blood in the urine, flank pain, or obstruction and is evaluated with imaging, urine testing, endoscopy, and tissue diagnosis.

\medterm{Renal Transplantation} Surgical placement of a functioning donor kidney into a person with advanced kidney failure. It can restore
kidney function but requires suitability assessment and long-term immunosuppression.

\synonyms
Kidney Transplant; Kidney Transplantation

\medterm{Renal Trauma} Kidney injury caused by a blow, fall, crash, sports injury, or penetrating wound. Assessment may include urine testing and CT imaging; treatment ranges from observation to procedures or surgery depending on severity.

\medterm{Renal Tubular Acidosis} A group of disorders in which kidney tubules do not remove enough acid or return enough bicarbonate to the blood. This makes the blood too acidic and may cause weakness, kidney stones, bone disease, or high potassium.

\medterm{Renal Tubular Acidosis Type 1} A metabolic acidosis caused by an impaired ability of distal intercalated renal tubular cells to secrete hydrogen ions into the collecting duct lumen. Characterized by normal anion gap metabolic acidosis, inappropriately high urinary pH (greater than 5.5 in the setting of systemic acidemia), hypokalemia, hypocitraturia, and recurrent nephrolithiasis or nephrocalcinosis.

\synonyms
Distal RTA; Type 1 RTA

\medterm{Renal Tubular Acidosis Type 2} A normal anion gap hyperchloremic metabolic acidosis caused by defective proximal tubular reabsorption of filtered bicarbonate, resulting in bicarbonate wasting into the urine. Urinary pH is typically elevated initially but drops below 5.5 when systemic acidosis becomes severe; frequently associated with generalized proximal tubular dysfunction (Fanconi syndrome) including glucosuria, aminoaciduria, phosphaturia, and osteomalacia.

\synonyms
Proximal RTA; Type 2 RTA

\medterm{Renal Tubular Acidosis Type 4} A normal anion gap metabolic acidosis characterized by persistent hyperkalemia caused by absolute aldosterone deficiency or renal tubular resistance to aldosterone (hypoaldosteronism). Impaired distal sodium reabsorption and potassium/hydrogen excretion reduce renal ammonium generation, yielding an acidic urinary pH (typically less than 5.5); common in diabetic nephropathy and chronic tubulointerstitial nephritis.

\synonyms
Hyperkalemic RTA; Hypoaldosteronemic RTA

\medterm{Renal Tubular Disorders} Diseases affecting the kidney tubules, which reabsorb water and useful substances and remove selected wastes into urine. They can disturb acid-base balance, electrolytes, urine concentration, or blood pressure.

\medterm{Renal Tubule} A tiny nephron tube that reabsorbs needed water and substances and removes others into urine.

\medterm{Renal Tubule Acidosis} A group of disorders in which the kidney tubules do not remove enough acid or conserve enough bicarbonate. The blood becomes too acidic, and complications may include low potassium, kidney stones, bone disease, weakness, or poor growth.

\medterm{Renal Ultrasound} Noninvasive imaging used to assess kidney size, structure, blood flow, and urinary obstruction. It is
particularly useful when obstruction, cystic disease, or structural abnormality is
suspected and is often used in the evaluation of acute or chronic kidney disease.

\medterm{Renal Vascular Disease} Alternate name; see \textit{Renal artery stenosis}.

\medterm{Renal Vein} One of two large veins that carry blood away from the kidneys to the body’s main abdominal vein. A clot in a renal vein can impair drainage and damage kidney function.

\medterm{Renal Vein Thrombosis} Renal vein thrombosis is a clot in a kidney vein that can cause flank pain, blood or protein in urine, kidney dysfunction, swelling, or pulmonary embolism.

\medterm{Renal Venogram} Contrast imaging of the renal veins used selectively to evaluate venous anatomy, thrombosis, obstruction, or vascular abnormalities.

\medterm{Renaut Body} A Renaut body is a specialized cluster of cells in the peripheral nerve endoneurium, associated with pressure sensitivity and sometimes nerve pathology.

\medterm{Renin} An enzyme from kidney cells that starts a hormone chain controlling blood pressure, salt balance, and fluid levels.

\medterm{Renin-Angiotensin-Aldosterone System} A hormone system that raises blood pressure and preserves body fluid when circulation or kidney blood flow falls. It narrows blood vessels and signals the adrenal glands to retain salt and water; overactivity contributes to hypertension and heart failure.

\medterm{Renin-Angiotensin-Aldosterone System Physiology} When blood pressure or kidney blood flow falls, the kidneys release renin. A chain of chemical signals then narrows blood vessels and prompts the adrenal glands to retain salt and water, raising blood pressure and restoring circulation.

\medterm{Renin-Angiotensin System} A hormone and enzyme system that regulates blood pressure, fluid balance, and vascular tone. Renin initiates formation of angiotensin II, which promotes vasoconstriction and aldosterone release; excessive activation contributes to hypertension and heart failure.

\medterm{Renin Blood Test} A blood test measuring renin concentration or activity, interpreted with aldosterone, posture, sodium status, medications, and blood pressure to evaluate hypertension or mineralocorticoid disorders.

\medterm{Rennin} An older spelling that may refer to renin or to rennet enzyme; in humans, renin regulates blood pressure and fluid balance.

\medterm{Renography} A kidney scan that follows a small amount of radioactive tracer as it enters and leaves the kidneys. It can compare kidney function and help identify reduced blood flow or blocked urine drainage.

\medterm{Renovascular Disease} Disease of the blood vessels supplying or draining the kidneys. It includes narrowing, blockage, aneurysm, or other vessel abnormalities and may reduce kidney blood flow or contribute to high blood pressure.

\medterm{Renovascular Hypertension} High blood pressure caused by reduced blood flow through one or both kidney arteries, usually from narrowing. It may be suspected when hypertension begins suddenly or is difficult to control; treatment addresses blood pressure and, in selected cases, the narrowed artery.

\medterm{Renshaw Cell} A nerve cell in the spinal cord that helps limit and regulate the activity of motor neurons.

\medterm{Reoviridae Infection} An infection caused by a reovirus; illness varies by virus and host and is often mild or asymptomatic in humans.

\medterm{Rep} An abbreviation most often meaning one repetition of an exercise or movement. In medical records it can also mean representative or another phrase, so the full wording should be used when the meaning is important.

\medterm{Repair Of Webbed Fingers Or Toes} Surgery to separate joined fingers or toes and reconstruct the skin between them. It may improve movement, growth, function, or appearance; wound problems, scarring, and recurrence of joining are possible.

\medterm{Repeat Surgery} Repeat surgery is a second or subsequent operation on the same or related area to treat persistent disease, recurrence, complications, or failure of prior treatment.

\medterm{Reperfusion} Restoration of blood flow to tissue after it has been reduced or interrupted. Reperfusion can save injured tissue but may also cause additional inflammation or oxidative injury.

\medterm{Reperfusion Injury} Damage that occurs when blood flow returns to tissue after a period of inadequate flow. The returning oxygen and inflammation can worsen injury, causing swelling, bleeding, or loss of function in organs such as the heart, brain, kidneys, or limbs.

\medterm{Reperfusion Therapy} Treatment that restores blood flow to tissue deprived of oxygen, such as clot-dissolving medicine, angioplasty, stenting, or surgery. Benefit depends on the tissue, duration of blockage, and risk of bleeding.

\medterm{Repetitive Motion Disorders} Conditions caused or worsened by repeated movement, force, vibration, or awkward posture. They may affect muscles, tendons, nerves, or joints and can cause pain, numbness, weakness, or reduced function.

\medterm{Repetitive Nerve Stimulation} A test that gives a nerve several small electrical impulses while recording the muscle’s response. A changing response may support a disorder in communication between nerves and muscles, such as myasthenia gravis or Lambert-Eaton syndrome. The result is interpreted with the examination and other tests.

\synonyms
RNS; Repetitive Stimulation Test

\medterm{Repetitive Strain Injury} Pain or dysfunction caused by repeated forceful movements, sustained positions, or overuse of muscles, tendons, nerves, or other tissues.

\medterm{Replantation Of An Avulsed Tooth} Putting a completely knocked-out permanent tooth back into its socket after an injury. The tooth should be handled by its crown, kept moist, and seen urgently; later splinting, root treatment, and follow-up may be needed to prevent loss.

\medterm{Replantation Of Digits} Microsurgical replacement of a completely amputated finger or toe, reconnecting bone, tendons, nerves, arteries, and veins when feasible.

\medterm{Replication Initiation} The time and process when copying of DNA or another genetic molecule begins. Specific proteins recognize a starting region, separate the strands, and recruit enzymes that build matching new genetic material.

\medterm{Replicon} A unit of genetic material that can replicate from one origin, such as a chromosome segment, plasmid, or viral genome.

\medterm{Repolarization} The return of the membrane potential toward the resting value after depolarization. In
neurons and skeletal/cardiac muscle, repolarization is primarily caused by the efflux of
potassium through voltage-gated potassium channels and inactivation of sodium channels.

\medterm{Reportable Diseases} Reportable diseases are illnesses that must be notified to public-health authorities under applicable laws to support surveillance and outbreak control.

\medterm{Reporting Standards} Agreed rules for recording and communicating clinical or imaging findings. Structured reports improve completeness and consistency and can make important results, follow-up needs, and comparisons easier for clinicians to understand.

\medterm{Repressed Memory Therapy} A controversial approach aimed at recovering allegedly repressed memories; suggestive techniques can create false memories and require caution.

\medterm{Repression} The reduction or stopping of gene activity, cell activity, behavior, or a body process. In genetics, regulatory proteins can silence genes; in psychology, the term may describe unconsciously keeping distressing thoughts outside awareness.

\medterm{Reproduction} The biological process by which living organisms produce new offspring, either sexually or without combining reproductive cells. In humans, it involves reproductive organs, hormones, eggs, sperm, pregnancy, and childbirth.

\medterm{Reproductive Cells} Cells used for reproduction, mainly sperm in males and eggs in females. Each normally carries half the usual number of chromosomes.

\medterm{Reproductive History} Reproductive history records pregnancies, births, miscarriages, abortions, menstrual function, fertility, contraception, and related conditions relevant to health assessment.

\medterm{Reproductive Hormones} Hormones that regulate sexual development, reproductive organs, ovulation, menstruation, sperm production, pregnancy, and lactation. Examples include estrogen, progesterone, testosterone, luteinizing hormone, and follicle-stimulating hormone.

\medterm{Reproductive Organs} The organs involved in reproduction. Female: ovaries (produce eggs and hormones),
fallopian tubes (transport eggs, site of fertilization), uterus (nurtures the developing
fetus), cervix (lower part of uterus, dilates during labor), vagina (passage for sperm
and baby), and vulva (external genitalia).

\medterm{Reproductive System} The organs and tissues that produce gametes and reproductive hormones and support sexual development and reproduction. Depending on the anatomy present, it may enable sperm or egg production, fertilization, pregnancy, childbirth, and lactation.

\medterm{Reproductive Technique} A reproductive technique is a medical or laboratory method used to assist conception, preserve fertility, diagnose reproductive disease, or manage pregnancy.

\medterm{Reproductive Toxicity} Harmful effects of a substance on fertility, pregnancy, the developing baby, birth outcomes, or later growth. Risks depend on the substance, dose, timing, route, and individual health and require careful exposure assessment.

\medterm{Reproductive Years} The part of life when a person may be able to become pregnant or produce sperm. Fertility varies widely and is affected by age, health, reproductive organs, hormones, and medicines; it should not be assumed from age alone.

\medterm{Reprogram} To reprogram a cell means to alter its gene-expression state so it acquires a different developmental potential or function.

\medterm{Resectable} Describing a tumor or lesion that can be removed completely by surgery with acceptable risk and an appropriate margin, based on its extent and relationship to vital structures.

\medterm{Resection} Surgical removal of part or all of an organ, tissue, or abnormal growth. The removed material may be examined to confirm diagnosis, determine treatment, or reduce disease while preserving useful function.

\synonyms
Surgical Excision; Surgical Removal

\medterm{Resectoscope} An endoscopic surgical instrument with a viewing system and working loop used to remove or cauterize tissue, especially during prostate or bladder procedures.

\medterm{Reserve Volume} The additional amount of air that can be inhaled or exhaled beyond a normal tidal breath.

\medterm{Reservoir} A human, animal, or environmental habitat in which an infectious agent normally lives,
multiplies, and persists. A reservoir may serve as the source from which the agent is
transmitted to a susceptible host.

\medterm{Reservoir Of Infection} A reservoir of infection is a person, animal, environment, or material in which an infectious organism normally lives and from which transmission can occur.

\medterm{Resident} A resident is a physician or other trainee receiving supervised specialty education after initial professional qualification; the term can also mean a person living in a place.

\medterm{Residential Mobility} Residential mobility is movement between homes or communities and can affect continuity of care, social support, exposure, and health access.

\medterm{Residual} Residual means remaining after a process, treatment, measurement, or event, such as residual urine, disease, risk, or volume.

\medterm{Residual Limb} The portion of a limb remaining after amputation. Care includes edema control, skin
inspection, shaping, strengthening, desensitization, contracture prevention, and
prosthetic preparation.

\medterm{Residual Limb Pain} Pain in the remaining portion of a limb after amputation.

\synonyms
Stump Pain

\medterm{Residual Neoplasm} A residual neoplasm is tumor tissue that remains after treatment intended to remove or destroy the original tumor.

\medterm{Residual Risk} Residual risk is the chance of harm that remains after preventive measures, treatment, screening, or controls have been applied.

\medterm{Residual Tumor} A residual tumor is cancer or other tumor tissue that remains after surgery, radiation, chemotherapy, or another treatment.

\medterm{Residual Volume} The amount of air remaining in the lungs after breathing out as forcefully as possible. It keeps small airways open and cannot be measured accurately by ordinary spirometry alone.

\medterm{Resistance} The ability of a microorganism to survive a medicine that previously inhibited or killed it. Resistance can make infections harder to treat and may spread between organisms through genetic material.
\synonyms
Antimicrobial Resistance; Drug Resistance

\medterm{Resistance Training} Exercise in which muscles work against weights, bands, body weight, or other resistance. Regular training can improve strength, endurance, balance, bone health, function, and independence when progressed safely.

\medterm{Resistant Hypertension} High blood pressure remaining above goal despite appropriate use of three or more medicines, including a water pill.

\synonyms
Refractory Hypertension; Treatment-Resistant Hypertension

\medterm{Resistant Starch} Starch that resists digestion in the small intestine and is fermented by gut bacteria in the colon.

\medterm{Resistin} A signaling protein associated with adipose tissue and inflammation; its human clinical significance remains under study.

\medterm{Resolution} The ability of an imaging system to distinguish separate structures, or the improvement or disappearance of a disease or symptom.

\medterm{Resonance} Amplification or characteristic vibration produced when a system is driven at a suitable frequency; used in sound and MRI contexts.

\medterm{Resorcinol} A phenolic compound used in selected topical dermatologic preparations; concentrated exposure can irritate or poison.

\medterm{Resorption} The body's breakdown and removal of tissue or material, especially old bone. The released components may be absorbed, reused, or removed from the body.

\medterm{Resorption Atelectasis} Atelectasis caused by complete obstruction of an airway, leading to absorption of
trapped air by the pulmonary capillaries. Common causes include mucus plugging, foreign
body, tumor, and post-surgical shallow breathing. The affected area appears as increased
density on CXR with volume loss.

\medterm{Respiration} The processes of breathing, gas exchange, and cellular energy production involving oxygen and carbon dioxide.

\medterm{Respiratory} Respiratory means relating to breathing, the airways, lungs, gas exchange, or the organs and muscles that move air.

\medterm{Respiratory Acidosis} A condition in which breathing does not remove enough carbon dioxide, making the blood too acidic. It can result from lung disease, blocked airways, weak breathing muscles, brain disorders, or medicines that slow breathing.

\medterm{Respiratory Alkalosis} A condition in which breathing removes too much carbon dioxide, making the blood less acidic. It can occur with anxiety, pain, fever, low oxygen, pregnancy, severe infection, or other conditions that cause rapid breathing.

\medterm{Respiratory Arrest} Stopping of effective breathing. It is an immediate emergency that may require airway support, rescue breathing or ventilation, and treatment of the cause.

\medterm{Respiratory Burst} A rapid release of chemically reactive oxygen molecules by certain immune cells after they engulf germs. This helps kill the germs inside the cell, although excessive activity can also damage nearby tissue.

\medterm{Respiratory Changes In Pregnancy} Pregnancy changes breathing as the growing uterus raises the diaphragm and hormones increase the drive to breathe. Breaths may become deeper, overall ventilation increases, and the amount of air left after a normal breath out decreases.

\medterm{Respiratory Compliance} The physical measure of the distensibility and elastic stretchability of the lungs and chest wall, expressed mathematically as the change in lung volume per unit change in distending pressure ($\Delta V / \Delta P$); reduced in acute respiratory distress syndrome, pulmonary edema, and pulmonary fibrosis.

\synonyms
Lung Compliance

\medterm{Respiratory Depression} Breathing that is abnormally slow or shallow, causing inadequate ventilation and possible
carbon-dioxide retention or low oxygen. Medicines, brain disorders, and metabolic illness
can cause it.

\medterm{Respiratory Disease} Any disease affecting the airways, lungs, breathing muscles, or gas exchange. Examples include asthma, chronic obstructive pulmonary disease, pneumonia, tuberculosis, and lung cancer.

\medterm{Respiratory Disorder} A disease affecting the airways, lungs, pleura, respiratory muscles, or control of breathing and causing impaired ventilation, oxygenation, or gas exchange.

\medterm{Respiratory Distress} Difficult or labored breathing caused by problems such as asthma, infection, a blood clot, a collapsed lung, an allergic reaction, or heart failure. Fast breathing, struggling to speak, blue lips, or severe exhaustion may occur.

\medterm{Respiratory Distress Syndrome} A serious breathing problem in premature newborns whose lungs do not make enough surfactant, a substance that keeps the tiny air sacs open. Babies may breathe rapidly, grunt, or need oxygen and breathing support. Antenatal steroids, noninvasive breathing support, and surfactant treatment improve survival.

\medterm{Respiratory Drive} The brain’s signal to the breathing muscles that controls how often and how deeply a person breathes. It responds mainly to carbon-dioxide and oxygen levels and can be reduced by brain disease, severe illness, or sedating medicines.

\medterm{Respiratory Failure} A life-threatening inability of the lungs to provide enough oxygen or remove enough carbon dioxide. It may cause severe breathlessness, confusion, blue or gray lips, extreme sleepiness, or collapse. Causes include lung disease, infection, heart problems, weak breathing muscles, brain injury, and drug overdose.

\medterm{Respiratory Failure Types 1 and 2} The clinical classification of acute respiratory failure: Type 1 (hypoxemic, $P_aO_2 < 60\text{ mmHg}$ with normal or low $P_aCO_2$) resulting from ventilation-perfusion mismatch or shunt (e.g., pneumonia, pulmonary edema); and Type 2 (hypercapnic, $P_aCO_2 > 45-50\text{ mmHg}$) resulting from alveolar hypoventilation (e.g., COPD, neuromuscular weakness, sedative overdose).

\medterm{Respiratory Gating} A radiation-treatment method that tracks breathing and delivers the beam only during a selected part of the breathing cycle. It helps account for movement of tumors in the chest or upper abdomen and may reduce radiation to nearby healthy tissue.

\medterm{Respiratory Health} The proper functioning of the lungs and airways, enabling efficient gas exchange and
oxygen delivery to the body. Maintaining respiratory health involves avoiding tobacco
smoke, minimizing exposure to air pollutants and occupational hazards, and receiving
vaccinations against respiratory infections like influenza and pneumococcal disease.

\medterm{Respiratory Illness} A short-term or long-term illness affecting breathing or the respiratory tract. Symptoms may include cough, fever, congestion, wheezing, shortness of breath, or chest discomfort, depending on the cause.

\medterm{Respiratory Infection} An infection of the upper or lower respiratory tract caused by viruses, bacteria, fungi, or other organisms.

\medterm{Respiratory Infections} Illnesses caused by pathogens including viruses, bacteria, and fungi that affect the airways and lungs.

\medterm{Respiratory Insufficiency} Inadequate breathing or gas exchange that causes low oxygen, high carbon dioxide, or both. It can result from lung or airway disease, weak breathing muscles, or problems in the brain’s breathing control and may progress to respiratory failure.

\medterm{Respiratory Muscle} A muscle used for breathing, especially the diaphragm and muscles between the ribs. These muscles expand and contract the chest to move air into and out of the lungs.

\medterm{Respiratory Neoplasm} A respiratory neoplasm is an abnormal growth in the airways or lungs and may be benign or malignant.

\medterm{Respiratory Paralysis} Respiratory paralysis is failure of the muscles needed for breathing because of nerve, muscle, brainstem, toxin, or spinal-cord dysfunction and is an emergency.

\medterm{Respiratory Physiology} The study of how breathing works, including movement of air, exchange of oxygen and carbon dioxide in the lungs, transport of gases in blood, and control of breathing.

\medterm{Respiratory Protection} Measures and equipment used to reduce inhalation of harmful dusts, fumes, gases, aerosols, or infectious particles.

\medterm{Respiratory Quotient} The ratio of carbon dioxide produced to oxygen used by the body’s cells. The value changes according to whether the body is mainly using carbohydrate, fat, or protein and helps estimate metabolism and energy needs.

\medterm{Respiratory Rate} The number of breaths taken per minute. The expected rate varies with age, activity, temperature, illness, and other factors.

\medterm{Respiratory Rehabilitation} A program of exercise, education, breathing techniques, and self-management that improves
function and quality of life in people with chronic respiratory disease.

\medterm{Respiratory Syncytial Virus} A respiratory virus that commonly causes bronchiolitis and pneumonia in infants
and young children. Infection is usually mild in adults but may be severe in infants, older adults, and people with
chronic heart or lung disease.

\synonyms
RSV; Human Orthopneumovirus

\medterm{Respiratory Syncytial Virus Bronchiolitis} An acute lower respiratory tract infection occurring predominantly in infants under two years of age caused by respiratory syncytial virus (RSV), characterized by bronchiolar mucosal inflammation, edema, and epithelial sloughing leading to wheezing, tachypnea, cough, and intercostal retractions.

\medterm{Respiratory Syncytial Virus Infection} Infection caused by respiratory syncytial virus. It commonly causes a cold-like illness but may involve the lower airways, producing bronchiolitis or pneumonia. Severity varies with age, immune status, and underlying heart or lung disease.

\medterm{Respiratory System} The nose, mouth, pharynx, larynx, trachea, bronchi, bronchioles, lungs, pleura, and breathing muscles, including the diaphragm and chest-wall muscles. It moves air into and out of the lungs and exchanges oxygen and carbon dioxide between the alveoli and blood.

\medterm{Respiratory Therapy} Respiratory therapy includes treatments and breathing support that improve ventilation, oxygenation, airway clearance, or lung function.

\medterm{Respiratory Tract Infections} Infections affecting any part of the respiratory tract, classified by anatomical
location into upper respiratory tract infections (URTIs) and lower respiratory tract
infections (LRTIs).

\medterm{Respiratory Transport} Respiratory transport is movement of oxygen and carbon dioxide between lungs, blood, and tissues, largely through hemoglobin and dissolved gas.

\medterm{Respiratory Viruses} Viruses that cause respiratory tract infections. Major respiratory viruses include
influenza, respiratory syncytial virus (RSV), parainfluenza virus, adenovirus, and human
metapneumovirus. SARS-CoV-2 causes COVID-19. Symptoms range from mild upper respiratory
illness to severe pneumonia.

\medterm{Respiratory Viruses And Face Masks} Face masks can reduce spread of respiratory viruses by blocking droplets and some aerosols, especially in crowded or healthcare settings; fit and local guidance matter.

\medterm{Respirologist} A doctor specializing in diseases of the lungs and breathing system. They assess breathing problems, interpret tests, manage chronic and acute illness, and may evaluate people for lung transplantation.

\medterm{Respirovirus} A genus of paramyxoviruses that includes human parainfluenza viruses, which can cause croup and other respiratory infections.

\medterm{Respite Care} Temporary care for a person who needs support so that a usual family or unpaid caregiver can rest, attend appointments, or manage other responsibilities. It may be provided at home, in a facility, or through short-term services.

\medterm{Response Assessment} Evaluation of how a disease responds to treatment. Clinicians may use symptoms, examination, blood tests, scans, tumor measurements, or activity levels to decide whether treatment is helping, failing, or needs adjustment.

\medterm{Response Evaluation Criteria} Standardized rules for judging whether a disease, especially a tumor, has responded to treatment, remained stable, or progressed. RECIST mainly uses tumor measurements on imaging; other systems may use metabolic or blood-based findings.

\medterm{Responsible Drinking} Responsible drinking means limiting alcohol according to health guidance, avoiding driving or hazardous activities after drinking, and avoiding alcohol when it is unsafe or contraindicated.

\medterm{Rest} Rest is reduced physical or mental activity used to recover, though excessive inactivity can impair strength and function in some conditions.

\medterm{Restenosis} Re-narrowing of a blood vessel after it has been opened by angioplasty or another procedure.

\medterm{Resting Energy Expenditure} The energy the body uses while resting to keep essential functions going, such as breathing, circulation, temperature control, and cell activity. It is affected by body size, age, muscle, illness, hormones, and recent food or activity.

\synonyms
REE; Resting Metabolic Rate; Basal Metabolic Rate

\medterm{Resting Membrane Potential} The small electrical difference across a nerve or muscle cell when it is not sending a signal. Unequal amounts of charged particles inside and outside the cell create it and allow the cell to respond quickly to stimulation.

\medterm{Resting Phase} The resting phase is a period of relative inactivity, such as G0 in the cell cycle or a physiologic interval between active events.

\medterm{Resting Tremor} An involuntary, rhythmic oscillation (typically four to six hertz) occurring in a body part that is fully supported against gravity and relaxed, which diminishes or temporarily ceases with voluntary movement. It typically presents as a pill-rolling hand tremor in Parkinson's disease and other parkinsonian syndromes.

\medterm{Restless Anal Syndrome} An uncommon term for uncomfortable sensations around the anus with an urge to move, touch, or relieve the area. The cause and best treatment are uncertain, so other anal and bowel conditions should be assessed.

\medterm{Restless Legs} Alternate name; see \textit{Restless Legs Syndrome}.

\medterm{Restless Legs Syndrome} A neurological and sleep-related condition causing an irresistible urge to move the legs, often with crawling, pulling, aching, or tingling sensations. Symptoms begin or worsen during rest, improve temporarily with movement, and are worse in the evening or at night; they may disrupt sleep and daytime function. Iron deficiency, kidney disease, pregnancy, some medicines, and genetic factors may contribute.

\synonyms
RLS; Willis-Ekbom Disease

\medterm{Restless Leg Syndrome} Alternate name; see \textit{Restless Legs Syndrome}.

\medterm{Restlessness} Restlessness is an uncomfortable need to move or inability to remain still, caused by anxiety, akathisia, restless-legs syndrome, pain, stimulants, withdrawal, or medical illness.

\medterm{Restorative Dentistry} Specialty of restoring damaged teeth. Procedures: fillings, inlays, onlays, crowns,
bridges, implants. Goals: function, aesthetics, durability.

\medterm{Restraint} A device, medicine, or other measure used to limit a person’s movement. Restraints can cause injury, distress, falls, and loss of independence, so trained staff should use the least restrictive option only when necessary and monitor closely.

\medterm{Restraint-Related Harm} Physical or psychological injury associated with using a device, medication, or
environmental restriction to limit a person's movement. Older adults are at risk of
falls, pressure injury, delirium, aspiration, and functional decline; the least
restrictive alternatives should be prioritized.

\medterm{Restriction Enzyme} An enzyme made by bacteria that cuts DNA at particular short sequences. Scientists use these enzymes to analyze, join, or modify DNA in laboratory research and medical testing.

\medterm{Restrictive Cardiomyopathy} A disease in which stiff heart muscle cannot relax and fill normally. It may cause breathlessness, tiredness, swelling, or abnormal rhythms and can result from conditions such as amyloidosis, sarcoidosis, iron overload, or scarring.

\medterm{Restrictive Lung Disease} A pattern of reduced total lung capacity caused by stiff lung tissue, chest-wall limitation, respiratory-muscle weakness, obesity, or other factors. It may cause exertional breathlessness and low oxygen; lung volumes, imaging, muscle strength, and the underlying cause guide management.

\medterm{Resuscitation} Emergency measures used to restore or support breathing and circulation, including opening the airway, rescue breathing, chest compressions, defibrillation, medicines, and treatment of the cause.
\synonyms
Cardiopulmonary Resuscitation; CPR; Clinical Revival

\medterm{Resuscitative Thoracotomy} Emergency chest surgery used for selected trauma patients who are in cardiac arrest or close to it, especially after a penetrating chest injury. It may relieve blood around the heart, control bleeding, or allow direct heart compressions. It is a last-resort procedure with a high risk of death and complications.

\synonyms
Emergency Department Thoracotomy; EDT; Clamshell Thoracotomy; Trauma Thoracotomy

\medterm{Retained Placenta} Failure of the placenta to come out after childbirth, usually when it remains in the uterus for about 30 minutes or longer or bleeding is excessive. It can cause severe bleeding or infection.

\medterm{Retained Products Of Conception} Pregnancy tissue that remains in the uterus after childbirth, miscarriage, or abortion. It may cause continuing or heavy bleeding, cramping, fever, or infection and may be treated with medicine or a procedure to remove it.

\medterm{Retainer} A dental appliance worn after orthodontic treatment to help keep teeth in their corrected positions.

\medterm{Retching} Repeated involuntary effort to vomit without necessarily producing stomach contents.

\medterm{Retention Of Urine} Inability to empty the bladder completely or at all despite making urine. It may cause painful fullness, leaking, or kidney damage and can result from blockage, weak bladder muscles, medicines, or nerve disease.

\medterm{Retention Time} The time a substance takes to pass through a laboratory separation system from injection until it is detected. Comparing it with a tested reference can help identify a substance, but timing alone is not definitive proof.

\medterm{Rete Testis} A network of small spaces inside the testis that collects sperm from the seminiferous tubules and passes them to the small ducts leading to the epididymis.

\medterm{Reticular} Net-like or arranged in a network pattern, describing the appearance or distribution of a skin lesion, rash, or vascular pattern.

\medterm{Reticular Activating System} A network of brain structures that helps maintain wakefulness, alertness, attention, and transitions between sleep and waking. Damage or severe disturbance can cause unusual sleepiness, reduced awareness, or coma.

\medterm{Reticular Dermis} The deeper layer of the dermis containing dense connective tissue, collagen, elastin, blood vessels, nerves, glands, and hair follicles.

\medterm{Reticular Erythematous Mucinosis} A rare primary cutaneous mucinosis presenting with erythematous, reticulated macules and plaques with dermal mucin accumulation on the mid-chest and upper back in middle-aged individuals, exacerbated by ultraviolet light exposure.
\synonyms
REM Syndrome; Plaque-like Cutaneous Mucinosis

\medterm{Reticular Formation} A network of nerve cells in the brainstem that helps control alertness, sleep, attention, muscle tone, posture, breathing, heart function, and responses to pain.

\medterm{Reticulin} A network-forming type of collagen that supports bone marrow, liver, lymph nodes, and other organs.

\medterm{Reticulocyte} A young red blood cell recently released from bone marrow. The number of reticulocytes shows how actively the marrow is replacing red cells and helps distinguish reduced production from blood loss or red-cell breakdown.

\medterm{Reticulocyte Count} The number or percentage of immature red blood cells in blood. It reflects recent bone-
marrow red-cell production and helps assess anemia, especially when interpreted with the hemoglobin and clinical context.

\medterm{Reticulocyte Production Index} A calculated hematological parameter that corrects the raw reticulocyte percentage for both the patient's degree of anemia (hematocrit) and the premature release of \"shift\" reticulocytes from the bone marrow into circulating blood. Calculated as raw reticulocyte percent multiplied by (patient hematocrit / normal hematocrit) divided by maturation factor in days; an RPI of 2.0 or greater indicates an adequate marrow erythropoietic response to hemolysis or acute hemorrhage, whereas an RPI below 2.0 signifies an inadequate hypoproliferative marrow response.

\synonyms
RPI; Corrected Reticulocyte Count

\medterm{Reticulocytes} Immature red blood cells recently released from bone marrow; their proportion helps assess red-cell production.

\medterm{Reticulocytopenia} An abnormally low reticulocyte count indicating reduced effective red-cell production, which may occur with marrow failure, nutrient deficiency, inflammation, or kidney disease.

\medterm{Reticulocytosis} Reticulocytosis is an increased number of immature red blood cells in blood, usually indicating stimulated marrow production after blood loss or hemolysis.

\medterm{Reticulo-Endothelial System} An older name for the body’s network of macrophages and related immune cells in the blood, liver, spleen, lymph nodes, and tissues. These cells remove germs, particles, and worn-out blood cells.

\medterm{Reticuloendothelial System} A diffuse anatomical network of specialized phagocytic cells (primarily monocytes and tissue macrophages) situated in reticular connective tissue frameworks of the spleen, liver (Kupffer cells), lymph nodes, and bone marrow, responsible for filtering particulate foreign matter, microbial pathogens, and senescent erythrocytes.

\synonyms
Mononuclear Phagocyte System; MPS

\medterm{Reticulophagy} A cell-cleaning process that removes worn or damaged parts of the endoplasmic reticulum, the cell's internal membrane network. The material is broken down and recycled, helping cells control quality and respond to stress.

\medterm{Reticulum} A network-like structure; in anatomy or cell biology it may refer to a tissue network, endoplasmic reticulum, or ruminant stomach compartment.

\medterm{Retina} The light-sensitive neural tissue lining the back of the eye that converts light into nerve signals.

\medterm{Retina Disorder} A disease affecting the retina, including vascular, degenerative, inflammatory, infectious, traumatic, inherited, or neoplastic conditions that can impair vision.

\medterm{Retina Hemorrhage} Bleeding within or beneath the retina caused by vascular disease, diabetes, hypertension, trauma, retinal vein occlusion, inflammation, or other disorders.

\medterm{Retinal Artery Occlusion} Blockage of a retinal artery causing sudden, usually painless loss of vision. It may indicate a wider blood-vessel or clotting disorder.

\medterm{Retinal Cone} A light-sensing cell in the retina that provides color vision and sharp detail, especially in daylight. Damage to cone cells can reduce color perception and central vision.

\medterm{Retinal Damage} Injury or disease affecting the retina, potentially causing blurred vision, field loss, flashes, floaters, or permanent visual impairment.

\medterm{Retinal Degeneration} Progressive damage or loss of light-sensitive cells in the retina at the back of the eye. It may cause blurred or reduced central vision, night blindness, or loss of side vision and can result from inherited disease, aging, or other injury.

\medterm{Retinaldehyde} A form of vitamin A used in the eye’s light-sensing cycle. Light changes it inside photoreceptor cells, helping start the nerve signals that allow vision.

\medterm{Retinal Detachment} A separation of the retina, the light-sensitive tissue at the back of the eye, from the layer beneath it. It may follow a retinal tear, scar tissue, or fluid buildup. Sudden floaters, flashes, a curtain-like shadow, or reduced vision may occur, and permanent sight loss is possible.

\medterm{Retinal Detachment Repair} Surgical or pneumatic treatment to close retinal breaks and reattach the retina. Techniques include vitrectomy, scleral buckling, and pneumatic retinopexy, selected according to the type and location of detachment.

\medterm{Retinal Diseases} Disorders that damage the retina, the light-sensitive tissue at the back of the eye. They may cause flashes, floaters, distorted vision, blind spots, night-vision problems, or gradual or sudden sight loss.

\medterm{Retinal Disorders} Diseases affecting the light-sensing layer at the back of the eye. Causes include blood-vessel problems, inherited conditions, aging, inflammation, infection, injury, and cancer, potentially causing blurred vision or blindness.

\medterm{Retinal Dysplasia} Abnormal development and arrangement of the retina, the light-sensing layer at the back of the eye. It may impair vision and can occur with conditions present from birth, genetic disorders, or problems affecting development before birth.

\medterm{Retinal Dystrophy} Gradual loss or damage of retinal cells, often caused by an inherited condition. It may begin with night blindness, difficulty seeing in dim light, or loss of side vision and can progress to serious sight impairment.

\medterm{Retinal Fold} A retinal fold is a crease or ridge in the retina caused by traction, developmental abnormality, detachment, swelling, or another ocular process.

\medterm{Retinal Fundus} The inside back surface of the eye viewed during fundoscopy, including the retina, optic nerve, macula, and retinal blood vessels.

\medterm{Retinal Ganglion} A nerve cell in the retina that collects signals from other retinal cells and sends them through the optic nerve to the brain. It is essential for carrying visual information out of the eye.

\medterm{Retinal Hemorrhages} Bleeding in the light-sensitive tissue at the back of the eye. It may result from diabetes, high blood pressure, injury, bleeding disorders, or increased pressure in the skull; in infants and children, the pattern must be assessed carefully for its cause.

\medterm{Retinal Implant} An implanted visual prosthesis that detects or processes light and electrically stimulates remaining retinal or visual-pathway cells in selected severe retinal disorders. It provides limited visual information and is distinct from an intraocular lens.

\medterm{Retinal Ischemia} Too little blood and oxygen reaching the retina, often because a retinal artery is blocked or blood-vessel disease is severe. It may cause sudden, painless sight loss and is an emergency because prolonged injury can become permanent.

\medterm{Retinal Laser Photocoagulation} An ophthalmic procedure utilizing targeted laser energy to induce localized thermal denaturation of retinal tissue. Widely used to seal microvascular leakages in diabetic macular edema, ablate ischemic peripheral retina in proliferative diabetic retinopathy, or secure retinal tears to prevent rhegmatogenous retinal detachment.

\synonyms
Laser Photocoagulation; Retinal Laser Therapy; Panretinal Photocoagulation

\medterm{Retinal Layers} The organized layers of nerve cells, supporting cells, and blood vessels that make up the retina. Together they detect light, begin visual processing, and send signals through the optic nerve to the brain.

\medterm{Retinal Macula} The macula is the central region of the retina responsible for detailed straight-ahead vision, reading, and color discrimination.

\medterm{Retinal Neoplasm} A retinal neoplasm is an abnormal growth arising in or near the retina, such as retinoblastoma or a metastatic tumor.

\medterm{Retinal Perforation} A full-thickness hole or tear in the retina. Fluid can pass through the opening and separate the retina from the tissue beneath it, causing flashes, floaters, or loss of vision.

\medterm{Retinal Pigment} Retinal pigment refers to melanin-containing pigment in the retinal pigment epithelium, which supports photoreceptors and absorbs excess light.

\medterm{Retinal Pigment Epithelium} A layer of cells beneath the retina that supports the light-sensing cells, absorbs extra light, and helps remove worn-out cell material.

\medterm{Retinal Prosthesis} An implanted device designed to provide limited visual signals when the retina is severely damaged. A camera and electronic system stimulate remaining visual pathways; it may provide light or simple pattern perception, not normal sight.

\medterm{Retinal Tear} A break through the retina, often caused when the jelly inside the eye pulls on it. Fluid can pass through the tear and detach the retina; new flashes, floaters, or a curtain-like shadow may occur. Treatment may include laser or freezing treatment.

\medterm{Retinal Vasculitis} Inflammation of blood vessels in the retina. Leaking or blocked vessels can cause blurred vision, bleeding, swelling, or permanent sight loss; it may be linked to infection, autoimmune disease, or inflammation elsewhere in the body.

\medterm{Retinal Vein Occlusion} A blockage in a vein that drains the retina. It can cause sudden or gradual blurred vision from bleeding and swelling, especially in the macula, and may lead to abnormal blood-vessel growth or permanent sight loss.

\medterm{Retinitis} Inflammation or infection of the retina, the light-sensitive tissue at the back of the eye. It may cause blurred vision, floaters, flashes, or loss of sight and needs prompt eye assessment.

\medterm{Retinitis Pigmentosa} A group of inherited retinal diseases in which the retina’s light-sensing cells gradually stop working. Inheritance varies by gene and may be autosomal dominant, autosomal recessive, X-linked, or mitochondrial. Early symptoms often include poor night vision and loss of side vision, followed by narrowing of the visual field and possible severe sight loss.

\medterm{Retinoblastoma} The most common eye cancer in infants and young children, starting in the retina at the back of the eye. It is caused by changes in the \textit{RB1} gene. It can be inherited (often affecting both eyes) or occur by chance (usually affecting only one eye). The most common early signs are leukocoria (an abnormal white glow in the pupil, often noticed in flash photographs instead of the normal red reflex) and strabismus (crossed or misaligned eyes). Diagnosis is made through a specialized dilated eye examination and imaging scans (ultrasound or MRI) to check tumor size and optic nerve involvement. Treatment focuses on saving the child's life and preserving vision using chemotherapy (given into the bloodstream or directly into the eye's blood vessels), laser therapy, freezing (cryotherapy), or surgical removal of the eye (enucleation) in advanced cases.
\synonyms{Childhood Eye Cancer; Retinal Glioma; Malignant Retinoblastoma}

\medterm{Retinoblastoma Protein} A protein made from the RB1 gene that helps prevent cells from dividing when they should not. Loss of its protective function can contribute to retinoblastoma and other cancers.

\medterm{Retinoic Acid} An active form of vitamin A that helps control how cells develop and which genes they use. Related medicines are used for selected skin conditions and some cancers but can cause irritation or serious pregnancy risks.

\medterm{Retinoid} A vitamin-A-related substance used in selected skin treatments and other medical care. Retinoids can improve acne or abnormal cell growth but may irritate skin, interact with medicines, and harm an unborn baby.

\medterm{Retinoids} Vitamin-A-related compounds used in selected dermatologic, ophthalmic, and oncologic treatments.

\medterm{Retinol} A form of vitamin A needed for vision, healthy skin and body linings, growth, and immune function. It comes from animal foods and some supplements; too much can be harmful, especially during pregnancy.

\medterm{Retinopathy} Damage or disease of the retina, the light-sensing layer at the back of the eye. Diabetes, high blood pressure, prematurity, blocked vessels, aging, inflammation, and inherited conditions can cause it and may lead to sight loss.

\medterm{Retinopathy Of Prematurity} Abnormal growth of blood vessels in the retina of a premature or very low-birth-weight infant. Mild disease may resolve, but severe disease can cause retinal detachment and blindness, so eligible infants need scheduled eye examinations.

\medterm{Retinoschisis} Splitting between layers of the retina. It may result from aging or an inherited condition and can cause reduced sharpness or loss of part of the visual field; it differs from retinal detachment but may look similar on examination.

\medterm{Retinoscope} An instrument that shines light into the eye to help determine refractive error. A lighted instrument used to estimate focusing error by observing how light reflects from the retina.

\medterm{Retinoscopy} An eye examination that uses a retinoscope to estimate refractive error by observing how light reflects from the retina.

\medterm{Retirement} Retirement is withdrawal from regular employment or a career; changes in income, activity, social contact, and health coverage can affect wellbeing.

\medterm{Retracted Nipple} Alternate name; see \textit{Nipple Inversion Or Retraction}.

\medterm{Retractile Testicle} Alternate name; see \textit{Retractile testicle}.

\medterm{Retractile Testis} A testis that moves between the scrotum and groin because of a strong reflex but can be gently brought into the scrotum. It should be monitored because it may later become undescended.

\medterm{Retractor} A surgical instrument used to hold tissue or an organ aside so the area can be seen and reached during an examination or operation. Some retractors are held by hand; others stay in place by themselves.

\medterm{Retrobulbar Hemorrhage} Bleeding behind the eyeball within the orbit, often after trauma or surgery, that can rapidly increase orbital pressure and threaten vision.

\medterm{Retrobulbar Neuritis} Inflammation of the optic nerve behind the eyeball, causing reduced or blurred vision and pain, especially when the eye moves. It may be associated with multiple sclerosis or another immune condition and needs prompt assessment.

\medterm{Retroesophageal} Retroesophageal means located behind the esophagus, such as a mass, lymph node, vessel, or collection in the posterior mediastinum.

\medterm{Retrognathia} Retrognathia is backward positioning of the lower jaw relative to the upper jaw and may affect facial structure, bite, airway, or sleep breathing.

\medterm{Retrograde} Moving backward or against the normal direction, such as retrograde amnesia or retrograde blood flow.

\medterm{Retrograde Amnesia} Loss of memory for events that occurred before an injury, illness, or other triggering event, with variable preservation of more remote memories.

\medterm{Retrograde Cystography} Retrograde cystography uses contrast placed through the urethra to outline the bladder and detect leaks, injury, or structural abnormalities.

\medterm{Retrograde Degeneration} Breakdown of a nerve fiber that spreads back toward the nerve cell body after the fiber is injured. It can reduce movement, sensation, or organ function depending on which nerve is affected and how severe the injury is.

\medterm{Retrograde Ejaculation} A condition in which semen enters the bladder instead of leaving through the penis during orgasm. It may cause little or no semen to appear and can result from diabetes, surgery, nerve problems, or medicines. It may affect fertility.

\medterm{Retrograde Intrarenal Surgery} RIRS is a minimally invasive procedure using a flexible scope passed through the urinary tract to inspect and fragment or remove kidney stones.

\medterm{Retrograde Urethrogram} An X-ray study that outlines the urethra after contrast material is introduced through its opening. It can show narrowing, tears, blockage, or abnormal connections, especially after suspected urethral injury.

\medterm{Retroperitoneal Fibrosis} Formation of abnormal scar-like tissue in the space behind the abdominal organs. It often surrounds the tubes draining urine from the kidneys, causing back or abdominal pain, urine blockage, and reduced kidney function; medicines, immune disease, or cancer can cause it.

\medterm{Retroperitoneal Hematoma} A collection of extravasated blood within the retroperitoneal anatomical space, commonly provoked by pelvic fractures, penetrating or blunt abdominal trauma, rupture of an abdominal aortic aneurysm, or femoral vascular catheterization; classified into zones I (central), II (flank), and III (pelvic).
\medterm{Retroperitoneal Hemorrhage} Bleeding in the space behind the abdominal lining. It may follow injury, blood-thinning treatment, a procedure, or rupture of a blood vessel and can cause back or abdominal pain, weakness, low blood pressure, and life-threatening blood loss.

\medterm{Retroperitoneal Inflammation} Retroperitoneal inflammation is inflammatory disease in the space behind the peritoneum, caused by infection, pancreatitis, urinary disease, trauma, medicines, or fibrosis.

\medterm{Retroperitoneal Neoplasm} A retroperitoneal neoplasm is an abnormal growth in the tissue space behind the abdominal lining and may be benign or malignant.

\medterm{Retroperitoneal Space} The retroperitoneal space is the area behind the peritoneal lining containing the kidneys, adrenal glands, pancreas, parts of the bowel, vessels, nerves, and lymphatics.

\medterm{Retroperitoneum} The space behind the abdominal lining. It contains the kidneys, adrenal glands, ureters, major blood vessels, and parts of the bowel and pancreas; bleeding, infection, tumors, or inflammation there can cause back or abdominal symptoms.

\medterm{Retropharyngeal Abscess} A collection of pus in the tissues behind the throat, usually after infection or injury. Fever, severe throat or neck pain, difficulty swallowing, drooling, noisy breathing, or a stiff neck can occur. The abscess can threaten the airway.

\medterm{Retropharyngeal Space} A potential fascial space in the neck located between the buccopharyngeal fascia of the pharynx anteriorly and the alar layer of the prevertebral fascia posteriorly, extending from the skull base down into the superior mediastinum (providing an anatomical conduit for the spread of retropharyngeal abscesses into the mediastinum).

\medterm{Retrospective Studies} Retrospective studies examine existing records or past exposures and outcomes; they are efficient but vulnerable to missing data, bias, and confounding.

\medterm{Retrospective Study} A study that examines existing records or past exposures to investigate disease causes, treatments, or outcomes. It can be efficient, but missing information and differences between groups may limit reliable conclusions.

\medterm{Retrosternal Thyroid Surgery} An operation to remove part or all of a thyroid gland that extends behind the breastbone. It may relieve pressure on the airway or swallowing tube and requires planning around nearby nerves, blood vessels, and the parathyroid glands.

\medterm{Retrotracheal} Retrotracheal means located behind the trachea, such as a mass, lymph node, vessel, or collection in the posterior mediastinum.

\medterm{Retrotransposition} The copying and insertion of a genetic sequence into a new place in the genome using an RNA copy as an intermediate. New insertions can be harmless or disrupt a gene and cause disease.

\medterm{Retroversion} Backward tilting or turning of an organ, especially the uterus toward the spine. A backward-tilted uterus is often normal, but in some people it is linked with pain, endometriosis, or difficult examinations.

\medterm{Retroversion Of The Uterus} A uterine position in which the uterus tilts backward toward the spine; it is often a normal variant but can be associated with pain or endometriosis.

\medterm{Retroviridae} A family of RNA viruses that make a DNA copy of their genetic material and insert it into infected cells. HIV is a member; these viruses can cause persistent infection, immune disease, or cancer.

\medterm{Retroviridae Infection} An infection caused by a retrovirus, which may persist through genomic integration and cause immune, neurologic, or malignant disease.

\medterm{Retrovirus} A type of virus that uses RNA as its genetic material instead of DNA. It uses an enzyme called reverse transcriptase to make a DNA copy of its RNA. This copy can become part of the DNA of an infected cell, allowing the virus to make more copies. HIV is an example of a retrovirus.

\medterm{Rett Syndrome} A rare genetic disorder that affects brain development and occurs almost exclusively in girls. Development may appear typical for the first 6 to 18 months, followed by loss of previously learned language, purposeful hand use, coordination, or movement. Repeated hand movements such as hand-wringing or clapping are common. Slowed head growth, seizures, breathing changes, feeding problems, sleep problems, and curvature of the spine (scoliosis) may occur. Most cases are caused by a new change in the \textit{MECP2} gene; rare males with related gene changes can also be affected.

\medterm{Return Of Spontaneous Circulation} Restoration of a sustained heartbeat and blood flow after cardiac arrest, shown by a pulse, blood pressure, or other signs of circulation. Intensive post-arrest care then supports breathing, blood pressure, brain function, temperature, and treatment of the cause.

\medterm{Reuptake} The removal of a released nerve chemical from the space between nerve cells by returning it to the releasing cell or taking it into a nearby cell. This limits and ends the signal.

\medterm{Revalidation} Formal reassessment that a healthcare professional remains fit, competent, and up to date to practice.

\medterm{Revascularisation} Restoration of blood flow to an organ or limb by a procedure such as angioplasty, stenting, or surgical bypass.

\medterm{Revascularization} Restoration of blood flow to an area by bypass surgery, balloon treatment, stent placement, or another procedure. It can relieve symptoms, preserve tissue, and reduce damage caused by a blocked or narrowed vessel.

\medterm{Revascularization Surgery} Revascularization surgery restores blood flow around a blocked or narrowed vessel, such as in coronary artery bypass or limb bypass surgery.

\medterm{Reverberation Artifact} An ultrasound image error caused by repeated reflections, sometimes reduced by changing the angle or using specialized imaging.

\medterm{Reversal Learning} Reversal learning is the ability to change a learned response when the previously rewarded or correct choice is no longer so.

\medterm{Reverse Total Shoulder Replacement} Shoulder-replacement surgery that places the rounded joint surface on the shoulder blade and the socket on the upper-arm implant. This changes shoulder mechanics and can help when severe arthritis occurs with an irreparable rotator cuff.

\medterm{Reverse Transcriptase} An enzyme that synthesizes DNA using RNA as a template.

\medterm{Reverse Transcriptase Inhibitor} An antiviral medicine that blocks reverse transcriptase, used in combination treatment for HIV or selected infections.

\medterm{Reverse Transcriptase Inhibitors} A major class of antiretroviral medications utilized in the management of HIV-1 infection and chronic hepatitis B, categorized into nucleoside/nucleotide reverse transcriptase inhibitors (NRTIs, such as emtricitabine and tenofovir) and non-nucleoside reverse transcriptase inhibitors (NNRTIs, such as efavirenz and doravirine).

\synonyms
RTIs

\medterm{Reverse Transcription} The production of DNA using RNA as a template, carried out by reverse transcriptase. Some viruses use this process to copy their genetic material, and laboratories use it to study RNA.

\medterm{Reverse Trendelenburg Position} A supine position in which the head and upper body are elevated higher than the feet. The head-up tilt uses gravity to shift abdominal organs downward, improving surgical visualization during head, neck, upper gastrointestinal, and laparoscopic procedures while facilitating ventilation.

\synonyms
Head-Up Tilt Position; Anti-Trendelenburg Position

\medterm{Reversible Cell Injury} Temporary cell damage that can resolve when the harmful cause is removed. Cells may swell, store excess fat, or produce less energy, but their membranes and genetic material remain sufficiently intact to recover.

\medterm{Revision Arthroplasty} Surgical replacement, exchange, or repair of a failed or painful joint prosthesis.
Indications include infection, aseptic loosening, instability, periprosthetic fracture,
wear, and recurrent dislocation.

\medterm{Revision Total Hip Arthroplasty} Surgery to remove and replace some or all parts of a previous hip replacement that have loosened, worn out, become infected, dislocated repeatedly, or failed. It is more complex than first-time replacement and may require bone reconstruction.

\medterm{Revision Total Knee Arthroplasty} Surgery to remove and replace parts of a previous knee replacement because of loosening, wear, infection, instability, stiffness, or persistent pain. It may require specialized implants or bone grafting and has higher risks than the first operation.

\medterm{Reye-Johnson Syndrome} Alternate name; see \textit{Reye Syndrome}.

\medterm{Reye's Syndrome} A rare, potentially fatal illness causing acute brain dysfunction and liver injury, associated particularly with viral illness and aspirin exposure in children.

\medterm{Reye Syndrome} A rare, serious illness causing sudden brain swelling and liver dysfunction, usually in a child or teenager recovering from a viral infection. Repeated vomiting, confusion, unusual sleepiness, seizures, or loss of consciousness may occur; aspirin exposure is associated with increased risk.

\synonyms
Reye-Johnson Syndrome; Fatty Liver with Encephalopathy

\medterm{Reynolds Pentad} A combination of fever or chills, pain in the upper right abdomen, jaundice, low blood pressure, and confusion. Together these signs suggest severe infection and blockage of the bile ducts, which requires emergency treatment.

\medterm{Rezafungin} A long-acting echinocandin antifungal medicine that inhibits fungal cell-wall synthesis. It is used for selected invasive
candidiasis infections; its long duration permits an intermittent dosing schedule under
specialist supervision.

\medterm{R-Factor} A laboratory measure comparing the amount of an antibiotic needed to stop a bacterium with the expected effective amount. A higher value suggests reduced susceptibility and may help guide antibiotic choice.

\medterm{Rgs Proteins} Regulators of G-protein signaling are proteins that speed the switch-off of signals from G-protein-coupled receptors. They help control responses to hormones, neurotransmitters, and other signals and may influence disease when altered.

\medterm{Rhabdoid Tumor} A rare, fast-growing cancer that usually affects infants and young children. It may begin in the kidney, brain, or other tissues; changes in genes that control cell growth, especially SMARCB1, are common.

\medterm{Rhabdomyolysis} Rhabdomyolysis is rapid breakdown of skeletal muscle that releases muscle proteins and salts into the bloodstream. Causes include crush injury, extreme exertion, seizures, heat illness, medicines, and metabolic disease. It may cause muscle pain, weakness, dark urine, kidney injury, abnormal heart rhythms, or dangerous pressure in a limb. A high creatine kinase level supports the diagnosis.

\medterm{Rhabdomyolysis And AKI} Severe muscle breakdown releases muscle proteins and salts, potentially causing acute kidney injury, abnormal body salts, and abnormal heart rhythms. Findings may include muscle pain, weakness, dark urine, and a high creatine kinase level.

\medterm{Rhabdomyoma} A usually noncancerous tumor made of heart-muscle cells, seen mainly in infants and children. It is strongly associated with tuberous sclerosis and may shrink on its own, but large or multiple tumors can block blood flow or disturb the heart rhythm.

\medterm{Rhabdomyosarcoma} Rhabdomyosarcoma is a malignant tumor showing skeletal-muscle differentiation, occurring most often in children but also in adults.

\medterm{Rhabdomyosarcoma Pathology} A malignant mesenchymal tumor showing skeletal-muscle differentiation. Histologic patterns include embryonal, alveolar, spindle-cell, and pleomorphic forms, with primitive cells, rhabdomyoblasts, and variable myogenin or MyoD1 expression.

\medterm{Rhegmatogenous Retinal Detachment} The most common type of retinal detachment, caused by a hole or tear that lets the eye’s internal fluid pass beneath the retina. It can cause flashes, floaters, a curtain-like shadow, or sight loss.

\medterm{Rhesus Factor} The Rh blood-group antigen, especially the D antigen, whose presence or absence is important in transfusion and pregnancy.

\medterm{Rhesus Incompatibility} Rhesus incompatibility occurs when an RhD-negative pregnant person forms antibodies against RhD-positive fetal red cells, risking haemolytic disease in a later or affected pregnancy; anti-D immunoglobulin prevents sensitization.

\medterm{Rhesus System} The Rh blood-group system, a group of inherited red-cell antigens including D, C, c, E, and e. RhD status is commonly reported as positive or negative and is important in transfusion compatibility and prevention of hemolytic disease of the fetus and newborn.

\medterm{Rheumatic Chorea} Involuntary, irregular, fidgety movements caused by an immune reaction after group A streptococcal infection. It is a feature of acute rheumatic fever and may affect the face, arms, legs, speech, or handwriting.

\medterm{Rheumatic Disease} A broad group of disorders affecting joints, muscles, bones, connective tissue, or immune function, including arthritis and lupus.

\medterm{Rheumatic Fever} An inflammatory complication of untreated or inadequately treated group A streptococcal throat infection, affecting joints, heart, skin, or brain.

\medterm{Rheumatic Heart Disease} Permanent damage to heart valves after rheumatic fever, an immune reaction following group A streptococcal throat infection. It most often affects the mitral valve and can cause breathlessness, tiredness, swelling, abnormal rhythms, or heart failure.

\medterm{Rheumatism} Pain, stiffness, or discomfort in joints and surrounding soft tissues.

\medterm{Rheumatoid Arthritis} A long-term autoimmune disease in which the immune system mistakenly attacks the lining of the joints. This causes inflammation, pain, swelling, warmth, and stiffness, often worse after waking or resting. It commonly affects the same joints on both sides of the body, especially the hands, wrists, and feet. Over time, the inflammation can damage cartilage and bone and reduce joint function. It may also cause tiredness, low-grade fever, or loss of appetite and can affect the eyes, lungs, heart, blood vessels, skin, nerves, or blood cells.

\medterm{Rheumatoid Arthritis and Pregnancy} Rheumatoid arthritis often changes during pregnancy and may flare after delivery. Care balances disease control, joint function, medication safety, contraception or fertility goals, and fetal or newborn considerations, with coordinated rheumatology and obstetric review.

\medterm{Rheumatoid Arthritis Cause Gastrointestinal Issues} Rheumatoid arthritis can affect the gastrointestinal system indirectly through inflammation, associated autoimmune disease, reduced mobility, or medicines such as nonsteroidal anti-inflammatory drugs. Abdominal pain, bleeding, persistent diarrhea, or weight loss needs evaluation rather than attribution to arthritis alone.

\medterm{Rheumatoid Arthritis Inflammation Of The Brain} Direct brain inflammation is uncommon in rheumatoid arthritis, but vasculitis, medication effects, infection, stroke risk, or another autoimmune process can cause neurologic symptoms. New confusion, seizure, weakness, severe headache, or vision change requires urgent assessment.

\medterm{Rheumatoid Factor} An antibody that mistakenly recognizes part of another antibody. A blood test may support the diagnosis of rheumatoid arthritis, but it can also be positive in other diseases or in healthy people.

\medterm{Rheumatoid Lung Disease} Lung involvement associated with rheumatoid arthritis, including interstitial lung disease, pleural disease, nodules, or airway abnormalities.

\medterm{Rheumatoid Nodule} A rheumatoid nodule is a firm inflammatory lump under the skin or near a joint in rheumatoid arthritis, often over pressure points.

\medterm{Rheumatoid Nodules} Firm inflammatory lumps under the skin or near joints in people with rheumatoid arthritis, often over pressure points. They may cause no symptoms but can become tender or be affected by pressure.

\medterm{Rheumatoid Pneumoconiosis} Rheumatoid pneumoconiosis is lung disease involving rheumatoid arthritis and dust-related lung nodules, historically called Caplan syndrome.

\medterm{Rheumatoid Synovitis} Long-lasting immune inflammation of the tissue lining a joint in rheumatoid arthritis. It causes pain and swelling and can gradually damage cartilage, bone, and joint movement.

\medterm{Rheumatoid Vasculitis} Inflammation of small or medium-sized blood vessels associated with rheumatoid arthritis. It can reduce blood flow and affect the skin, nerves, eyes, or internal organs.

\medterm{Rheumatologist} A medical doctor who diagnoses and treats diseases affecting joints, muscles, bones, and connective tissues, especially autoimmune and inflammatory conditions such as rheumatoid arthritis, lupus, gout, and vasculitis.

\medterm{Rheumatology} The medical specialty concerned with inflammatory, autoimmune, and musculoskeletal disorders.

\medterm{Rh Factor} A system of antigens on the surface of red blood cells, especially the D antigen. A person
is Rh-positive when the D antigen is present and Rh-negative when it is absent. Rh incompatibility can affect pregnancy
and transfusion safety.

\synonyms
Rh D Antigen; Rhesus Antigen

\medterm{Rh Incompatibility} An immune reaction that can occur when an Rh-negative pregnant person carries an Rh-positive fetus and forms antibodies against fetal red blood cells. Anti-D immune globulin can prevent sensitization.

\medterm{Rhinencephalon} An older term for brain regions involved mainly in smell and closely connected with emotion and memory. Modern anatomy usually names the individual structures instead of using this broad term.

\medterm{Rhinitis} Inflammation or irritation of the lining of the nose. It can cause a blocked or runny nose, sneezing, itching, and reduced smell and may result from allergy, infection, smoke, strong smells, medicines, or other triggers.

\medterm{Rhinitis Medicamentosa} A condition of severe rebound nasal mucosal congestion and edema caused by prolonged, repetitive use of topical alpha-adrenergic nasal decongestant sprays (such as oxymetazoline) beyond 3 to 5 days.
\synonyms
Rebound Nasal Congestion; Topical Decongestant Rhinitis

\medterm{Rhinocephaly} A very rare birth defect involving severe abnormal development of the face and forebrain, with a nose-like structure that may lie unusually high or between the eyes. It is usually associated with major neurologic and developmental abnormalities.

\medterm{Rhinocerebral Mucormycosis} A rapidly progressive, fulminant, angioinvasive fungal infection caused by molds of the order Mucorales (such as \textit{Rhizopus} and \textit{Mucor}), occurring primarily in patients with diabetic ketoacidosis, severe neutropenia, or hematologic malignancies. Fungal hyphae invade blood vessels causing extensive thrombosis, tissue infarction, and black necrotic eschars of the nasal mucosa and hard palate, rapidly extending into the orbit and cavernous sinus.

\synonyms
Invasive Rhinocerebral Zygomycosis

\medterm{Rhinophyma} Progressive thickening and enlargement of the nose caused by long-standing changes often associated with rosacea. It may produce a rough, irregular appearance and can sometimes interfere with breathing or require surgical treatment.

\medterm{Rhinoplasty} Surgery to reshape, repair, or reconstruct the nose for breathing, injury, birth differences, or appearance. Recovery includes swelling and bruising, and results depend on anatomy, technique, healing, and patient goals.

\medterm{Rhinorrhea} Excessive watery or mucoid discharge from the nose, commonly caused by infection, allergy, irritants, or changes in temperature.

\medterm{Rhinorrhoea} British spelling of rhinorrhea, meaning drainage of watery or mucus-like fluid from the nose. Causes include allergy, infection, irritation, and rarely a cerebrospinal-fluid leak.

\medterm{Rhinoscleroma} A chronic, slowly progressive granulomatous infection of the nose and upper airway caused by Klebsiella rhinoscleromatis. It may progress from foul-smelling discharge and obstruction to scarring, deformity, and airway narrowing and usually needs prolonged antibiotics.

\medterm{Rhinoscopy} Examination of the nasal cavity and nasopharynx using an instrument or endoscope to assess the mucosa, septum, turbinates, masses, bleeding, or obstruction.

\medterm{Rhinosporidiosis} A chronic localized granulomatous disease of the mucous membranes caused by Rhinosporidium seeberi, an aquatic protistan parasite in the clade Mesomycetozoea, endemic primarily to rural wetland communities in India and Sri Lanka. It classically manifests in the anterior nasal cavity or nasopharynx as painless, friable, highly vascular, sessile or pedunculated polypoid lesions with a characteristic "strawberry-like" appearance dotted with visible white spherules (sporangia), causing nasal obstruction, epistaxis, and foreign-body sensation.

\synonyms
Nasal Rhinosporidiosis; Rhinosporidium Seeberi Infection

\medterm{Rhinovirus} A genus of non-enveloped, positive-sense RNA viruses in the Picornaviridae family. Rhinoviruses
are a common cause of the common cold and spread mainly through respiratory secretions and contaminated hands or
surfaces.

\medterm{Rhinovirus A} A species group of rhinoviruses that commonly causes the common cold and can worsen asthma or trigger wheezing.

\medterm{Rhinovirus C} A rhinovirus group associated with common-cold symptoms and, in some children, more severe wheezing or lower-respiratory illness.

\medterm{Rhinoviruses} A group of viruses that commonly cause colds. They spread through droplets, hands, and contaminated surfaces and may worsen asthma or cause more serious breathing illness in vulnerable people.

\medterm{Rh Isoimmunization} Maternal alloantibody production against fetal Rh D-positive red blood cells in an Rh
D-negative mother, occurring when fetal RBCs enter the maternal circulation
(fetomaternal hemorrhage).

\medterm{Rhizomucor} A genus of molds that can cause mucormycosis, especially in people with diabetes, neutropenia, or severe immune suppression.

\medterm{Rhizomucor Pusillus} A mold that can cause invasive mucormycosis involving the lungs, sinuses, skin, or other organs in severely immunocompromised people.

\medterm{Rhizopus} A genus of molds that includes important causes of mucormycosis, a rapidly invasive infection of sinuses, lungs, skin, or other tissues.

\medterm{Rhizopus Oryzae} A mold, also called Rhizopus arrhizus, that is a common cause of invasive mucormycosis.

\medterm{Rhizotomy} A procedure that cuts or heats selected nerve roots to reduce severe pain, muscle stiffness, or abnormal reflexes. It is used only in carefully selected cases because it can cause permanent numbness, weakness, or other nerve problems.

\medterm{Rhodium} A chemical element used in industrial and dental materials; its compounds can irritate tissues and may cause allergic reactions.

\medterm{Rhodococcus} A genus of gram-positive bacteria that includes Rhodococcus equi, an opportunistic pathogen causing pneumonia mainly in immunocompromised people.

\medterm{Rhodococcus Equi} A gram-positive bacterium that can cause cavitary pneumonia and disseminated infection, especially in people with advanced immune suppression.

\medterm{Rhodopsin} A light-sensitive protein in rod cells of the retina that enables vision in dim light.

\medterm{Rhodospirillaceae} A family of metabolically diverse bacteria found in soil and water; some species carry out photosynthesis or nitrogen fixation.

\medterm{Rhodospirillum} A genus of spiral-shaped aquatic bacteria capable of photosynthesis or nitrogen fixation; human infection is rare.

\medterm{RhoGAM} A brand name for Rh immune globulin, which contains anti-D antibodies. It is given to eligible RhD-negative pregnant people after specified pregnancy events and around delivery to reduce formation of antibodies that could harm an RhD-positive fetus or newborn.

\medterm{Rhombencephalon} The hindbrain region of an early embryo. It develops into the pons, medulla, cerebellum, and related structures that help control breathing, balance, movement, swallowing, and other essential functions.

\medterm{Rhonchi} Low-pitched breathing sounds caused by air moving through narrowed or mucus-filled larger airways, often changing after coughing.

\medterm{Rhonchus} An abnormal sound heard through a stethoscope during breathing, caused by air moving through fluid, mucus, narrowed airways, or abnormal lung tissue.

\synonyms
Crackles; Crepitations; Rales

\medterm{Rhubarb Leaves Poisoning} Poisoning after eating large amounts of rhubarb leaves, which contain oxalates and other irritants. It may cause burning, vomiting, abdominal pain, low calcium, kidney injury, weakness, or abnormal heart rhythms.

\medterm{Rhythm Method} A fertility-awareness method that estimates fertile days from menstrual-cycle timing to avoid or achieve pregnancy.

\medterm{Rib} One of 12 paired, curved bones forming the thoracic cage, protecting the heart, lungs,
and great vessels. True ribs (1--7) attach directly to the sternum via costal cartilage.
False ribs (8--10) attach indirectly to the sternum via the costal cartilage of rib 7.
Floating ribs (11--12) have no anterior attachment.

\synonyms
Costal Bone; Costa

\medterm{Rib Angle} The point where a rib bends most sharply as it curves from the spine toward the front of the chest.

\medterm{Ribavirin} An antiviral medicine used in selected viral infections, usually in combination and under specialist supervision.

\medterm{Rib Cage} The rib cage consists of ribs, sternum, thoracic vertebrae, muscles, and connective tissues that protect chest organs and assist breathing.

\medterm{Ribcage Pain} Ribcage pain may arise from muscles, ribs, joints, lungs, heart, or other organs; sudden severe pain or breathing difficulty may indicate a serious condition.

\medterm{Rib Fracture} A break in one or more ribs, usually causing localized chest pain worsened by breathing, coughing, or movement and sometimes associated with injury to underlying organs.

\medterm{Rib Head} The back end of a rib that joins with the bodies of one or two thoracic vertebrae.

\medterm{Rib Neck} The narrowed part of a rib between its head and tubercle.

\medterm{Riboflavin} A water-soluble vitamin that forms the coenzymes FAD and FMN. These cofactors
participate in oxidation-reduction reactions, including electron transfer, fatty-acid
oxidation, and energy metabolism.

\synonyms
Lactoflavin; Vitamin G

\medterm{Riboflavin Deficiency} Deficiency of vitamin B2, which can cause inflammation or fissuring of the mouth and lips, sore throat, dermatitis, anemia, and other mucosal or systemic findings.

\medterm{Ribonuclease} An enzyme that cuts RNA, a molecule involved in reading genetic instructions, into smaller pieces. Ribonucleases help control how long RNA lasts, recycle cell material, and protect cells from some unwanted genetic material.

\medterm{Ribonuclease H} An enzyme that removes the RNA strand when RNA and DNA are joined together. Cells use it during DNA copying, and some viruses depend on related activity when making DNA from their RNA.

\medterm{Ribonuclease III} An enzyme that cuts double-stranded RNA into smaller pieces. The products can help control gene activity, make proteins, or process genetic material in cells and some microorganisms.

\medterm{Ribonucleic Acid} RNA is a nucleic acid that carries genetic instructions, helps make proteins, and regulates cell activity; some viruses use RNA as their genetic material.

\medterm{Ribonucleoprotein} A structure made of RNA joined to one or more proteins. It helps make, process, transport, or control RNA; antibodies against certain ribonucleoproteins can occur in autoimmune diseases.

\medterm{Ribonucleoside} A molecule made of a ribose sugar joined to one of the bases used in RNA. Examples include adenosine, guanosine, cytidine, and uridine; adding phosphate groups produces RNA building blocks.

\medterm{Ribonucleotide Reductase} An enzyme that changes RNA building blocks into the related building blocks needed to make DNA. Cells require it to copy and repair DNA, so it is important for growth and is a target of some cancer medicines.

\medterm{Ribonucleotides} Nucleotides containing ribose sugar and RNA bases, serving as RNA building blocks and cellular energy or signaling molecules.

\medterm{Ribophagy} A cell-cleaning process that breaks down some ribosomes, the structures that make proteins, so their parts can be recycled. It helps cells adjust protein production when nutrients or growth conditions change.

\medterm{Ribose} A five-carbon sugar used to build RNA, ATP, and other molecules involved in storing genetic information and transferring energy. Cells make ribose from other nutrients and use it in many normal chemical processes.

\medterm{Ribosomal Protein} A protein that combines with ribosomal RNA to form a ribosome, the cell structure that builds proteins. Inherited changes in some ribosomal proteins can cause bone-marrow failure or developmental syndromes.

\medterm{Ribosomal RNA} RNA that forms the main structural and working part of ribosomes, the cell structures that make proteins. It helps position messenger RNA and joins amino acids together; its genes are also used in some laboratory studies to identify organisms.

\medterm{Ribosome} A cellular structure made of RNA and proteins that reads messenger RNA and links amino acids together to make proteins.

\medterm{Ribosome Subunit} One of the two RNA-and-protein parts of a ribosome. The small and large parts join around messenger RNA and work together to assemble amino acids into a protein.

\medterm{Rib Tubercle} A small projection near the back of a rib that joins with the transverse process of a thoracic vertebra.

\medterm{Richter's Syndrome} Transformation of chronic lymphocytic leukemia or small lymphocytic lymphoma into an aggressive lymphoma, most often diffuse large B-cell lymphoma.

\medterm{Richter Transformation} The rapid clinico-pathological transformation of chronic lymphocytic leukemia (CLL) or small lymphocytic lymphoma (SLL) into an aggressive, high-grade diffuse large B-cell lymphoma (DLBCL), presenting with rapidly enlarging lymphadenopathy, systemic B symptoms, and elevated serum LDH.
\synonyms
Richter's Syndrome; High-Grade Transformation of CLL

\medterm{Ricin} Ricin is a highly toxic plant protein that blocks protein synthesis; ingestion, inhalation, or injection can cause severe organ injury and death.

\medterm{Rickets} Defective mineralization of growing bone, usually due to vitamin D, calcium, or phosphate deficiency.

\synonyms
Nutritional Rickets; Hypocalcemic Rickets

\medterm{Rickettsia} A group of small intracellular bacteria transmitted mainly by arthropods and causing diseases such as typhus and spotted fevers.

\medterm{Rickettsiaceae} A family of bacteria that live inside cells and includes several organisms spread by ticks, fleas, lice, or mites. They can cause fever, rash, blood-vessel inflammation, and serious organ complications.

\medterm{Rickettsia Felis} A flea-associated bacterium that can cause a spotted-fever-like illness in people. Infection may produce fever, headache, rash, fatigue, or muscle aches and is assessed with exposure history and laboratory testing.

\medterm{Rickettsia Infections} Rickettsial infections are illnesses caused by intracellular bacteria spread mainly by ticks, fleas, lice, or mites and often causing fever and rash.

\medterm{Rickettsia Japonica} A bacterium spread by ticks and associated with Japanese spotted fever. Infection may cause fever, headache, rash, and a dark mark where the tick bit; treatment uses suitable antibiotics.

\medterm{Rickettsialpox} An infection spread by house mites and caused by Rickettsia akari. It usually begins with a small dark sore at the bite site, followed by fever and a blister-like rash.

\medterm{Rickettsia Prowazekii} A bacterium that causes epidemic typhus and is mainly spread through infected body-lice feces. Disease can cause high fever, headache, rash, confusion, and serious organ complications without prompt antibiotic treatment.

\medterm{Rickettsia Rickettsii} A bacterium spread by infected ticks that causes Rocky Mountain spotted fever. Fever, severe headache, rash, muscle aches, or confusion may develop; prompt antibiotic treatment is important because the illness can rapidly become life-threatening.

\medterm{Rickettsia Typhi} A bacterium spread mainly by infected fleas and responsible for murine typhus. The illness commonly causes fever, headache, chills, and sometimes a rash; prompt antibiotic treatment usually works, but untreated infection can become severe.

\medterm{Riddoch Phenomenon} A form of visual statuagnosia in which a patient with homonymous hemianopia or cortical blindness from an occipital lobe lesion is able to perceive moving visual targets within the blind field while remaining unable to visualize stationary, static stimuli.
\synonyms
Statuagnosia; Visual Motion Perception Preservation

\medterm{Ridged Cranial Sutures} Ridged cranial sutures are raised skull seams in an infant and may reflect normal molding or premature suture fusion, which requires assessment for craniosynostosis.

\medterm{Ridged Nail} A ridged nail has longitudinal or transverse grooves in its surface, which may result from aging, trauma, inflammation, nutritional problems, or systemic disease.

\medterm{Ridge Preservation} A dental procedure performed after a tooth is removed to reduce loss of the jawbone that held the tooth. Material is placed in the empty socket, sometimes with a covering, to help maintain shape for healing or a future implant.

\medterm{Riedel's Thyroiditis} A rare form of thyroid inflammation in which tough scar-like tissue makes the thyroid hard and may spread into nearby structures. It can cause a neck lump, difficulty swallowing or breathing, or reduced thyroid function.

\medterm{Riedel Thyroiditis} A rare, chronic fibroinflammatory disorder of the thyroid gland (an organ manifestation of IgG4-related disease) characterized by dense, invasive collagenous fibrous tissue replacing normal thyroid tissue and extending into adjacent strap muscles and trachea, presenting as a stony-hard, fixed ('woody') goiter.
\synonyms
Invasive Fibrous Thyroiditis; Woody Thyroiditis

\medterm{Rieger Syndrome} A genetic condition affecting development of the front of the eye and often the teeth. It may cause glaucoma, unusual teeth, or changes around the navel and can result from changes in genes such as PITX2 or FOXC1.

\medterm{Riehl Melanosis (Pigmented Contact Dermatitis)} An acquired gray-brown to slate-colored facial or neck hyperpigmentation associated with contact or photocontact exposure to fragrances, cosmetics, dyes, or other chemicals. Itching or redness may precede pigmentation; identifying and avoiding the trigger, sometimes with patch or photopatch testing, helps guide management.

\medterm{Rifabutin} A rifamycin antibiotic used mainly for selected mycobacterial infections and tuberculosis regimens; it has important drug interactions and can discolor body fluids.

\medterm{Rifampicin} A rifamycin antibiotic used in tuberculosis and selected other bacterial infections, usually in combination to prevent resistance.

\medterm{Rifampin} Rifampin is a rifamycin antibiotic that inhibits bacterial DNA-dependent RNA polymerase and is a key component of tuberculosis treatment; it causes orange body fluids and has many drug interactions and hepatotoxicity risk.

\medterm{Rifampin Toxicity} Adverse pharmacological effects of the first-line rifamycin antibiotic rifampin, including benign orange-red discoloration of body secretions (urine, tears, sweat), drug-induced hepatotoxicity, an immune-mediated flu-like syndrome with intermittent dosing, and potent induction of hepatic cytochrome P450 enzymes (CYP3A4).

\medterm{Rifamycin} Rifamycin is an antibiotic that inhibits bacterial RNA polymerase; related rifamycin medicines are used for tuberculosis and other infections.

\medterm{Rifaximin} Rifaximin is a minimally absorbed rifamycin antibiotic used for selected intestinal infections, prevention of recurrent hepatic encephalopathy, and some diarrhoeal syndromes; indication and resistance patterns guide use.

\medterm{Rift Valley Fever} A mosquito-borne viral disease affecting animals and humans, usually causing fever but sometimes severe eye, brain, or hemorrhagic disease.

\medterm{Right Atrium} The upper-right chamber of the heart. It receives oxygen-poor blood from the body and passes it through the tricuspid valve into the right ventricle, which sends it to the lungs.

\medterm{Right Bundle Branch} A group of heart fibers that carries the electrical signal to the right lower chamber so it contracts. If conduction is delayed or blocked, the ECG shows right bundle-branch block; this may be harmless or reflect heart or lung disease.

\medterm{Right Bundle Branch Block} Delay or interruption of electrical signals through the right bundle branch. It changes the ECG pattern and may occur without symptoms or with heart or lung disease.

\synonyms
RBBB

\medterm{Right Coronary Artery} A blood vessel that branches from the aorta and supplies much of the right side and lower part of the heart. A blockage can cause a heart attack or abnormal heart rhythm.

\medterm{Right Coronary Artery Occlusion} A blockage of the artery supplying much of the right and lower heart. It can cause a heart attack, slow or abnormal rhythms, low blood pressure, or failure of the right side of the heart.

\medterm{Right Heart Failure} Failure of the right side of the heart to pump blood effectively to the lungs. It may result from left-sided heart failure, high pressure in the lung arteries, a heart attack, lung clots, heart-muscle disease, or a stiff sac around the heart. It can cause swollen legs, neck veins, an enlarged liver, abdominal fluid, and weight gain.

\medterm{Right Heart Ventricular Angiography} An imaging test in which contrast dye and X-rays show the right lower heart chamber and how it pumps blood. It is used selectively to assess structure, movement, or blood flow during heart evaluation.

\medterm{Right Hypochondriac Region} The upper-right area of the abdomen, beneath the lower ribs. It contains much of the liver, the gallbladder, and part of the intestine. Pain there may arise from these organs.

\medterm{Right Iliac Region} The lower-right area of the abdomen, containing the appendix, beginning of the large intestine, and part of the small intestine. Pain there may occur with appendicitis, bowel disease, urinary problems, or a hernia and must be interpreted with other findings.

\medterm{Righting Reflex} The righting reflex is an automatic response that restores head and body orientation when balance or position changes.

\medterm{Right Lower Quadrant Abdominal Pain} Alternate name; see \textit{Appendicitis}.

\medterm{Right Lymphatic Duct} A short vessel that returns lymph from the right arm, right side of the head and neck, and right upper chest to the bloodstream near the right side of the neck. Damage or blockage can cause swelling in these areas.

\medterm{Right Sixth Cranial Nerve Palsy} Weakness of the nerve that moves the right eye outward. It can cause double vision, usually worse when looking to the right or at a distance, and may result from diabetes, high pressure in the skull, inflammation, stroke, or other disease.

\medterm{Right Ureterovesical Calculus} A stone lodged at the junction of the right ureter and bladder, causing flank pain,
hematuria, and urinary urgency or frequency. May require endoscopic or lithotripsy
intervention.

\medterm{Right Ventricle} The lower-right chamber of the heart. It receives oxygen-poor blood from the right atrium and pumps it through the pulmonary valve into the lung arteries to receive oxygen.

\medterm{Right Ventricular Strain} An electrocardiographic or echocardiographic pattern reflecting acute or chronic right ventricular pressure and volume overload, manifesting on ECG with T-wave inversions in leads V1 through V4, S1Q3T3 pattern, right bundle branch block, or rightward axis deviation (characteristic of acute massive pulmonary embolism).

\medterm{Rigid And Flexible Sigmoidoscopy} A sigmoidoscopy examines the rectum and lower colon with a flexible or rigid instrument. It can evaluate bleeding, pain, altered bowel habits, inflammation, polyps, or tumors; preparation, reach, comfort, and biopsy capability differ by technique.

\medterm{Rigidity} Abnormally increased resistance when a relaxed limb is moved by another person. It can occur with Parkinson disease, brain or spinal injury, severe infection, certain medicines, or other movement disorders.

\medterm{Rigler Sign} A classic abdominal radiographic sign of massive pneumoperitoneum (free intra-abdominal air) on plain supine radiographs, where both the inner (mucosal) and outer (serosal) surfaces of the bowel wall are clearly outlined by gas.
\synonyms
Double Wall Sign; Bas-Relief Sign

\medterm{Rigler Triad} Three imaging findings that strongly suggest gallstone ileus: blockage of the small intestine, air in the bile ducts, and a gallstone in the intestine. Not every person has all three findings.

\medterm{Rigor} Marked stiffness or rigidity; rigor mortis is postmortem muscle stiffening, while rigor may also mean a severe chill.

\medterm{Rigor Mortis} Stiffening of the body’s muscles after death. It develops as the muscles lose the energy needed to relax and later fades as decomposition advances; its timing is affected by temperature, activity, illness, and the environment.

\medterm{Rilpivirine} An HIV medicine that blocks an enzyme the virus needs to copy itself. It is taken with other antiretroviral medicines, and missed doses or drug interactions can allow resistance to develop.

\medterm{Riluzole} A medicine used to modestly slow disease progression in amyotrophic lateral sclerosis and to treat selected forms of spinal muscular atrophy.

\medterm{Rimmed Vacuoles} Small empty-looking spaces with a visible border seen inside muscle cells on biopsy. They indicate that muscle cells are breaking down or recycling parts of themselves and can occur in conditions such as inclusion-body myositis and GNE myopathy. The finding must be interpreted with symptoms and other tests.

\synonyms
Autophagic Rimmed Vacuoles

\medterm{Ring Block} Local anesthetic injected around the base of a finger or toe to numb the entire digit.

\medterm{Ring Chromosome} A chromosome whose two ends have joined to form a ring, often with loss or disruption of genetic material. It may cause developmental or other health problems, depending on the chromosome and changes involved.

\medterm{Ringing In The Ear} Alternate name; see \textit{Tinnitus}.

\medterm{Ringing In The Ears} Alternate name; see \textit{Tinnitus}.

\synonyms
Tinnitus; Aurium Acusticus

\medterm{Ring Vaccination} A public-health strategy in which close contacts and contacts of contacts around an identified case are vaccinated to create a protective ring and limit transmission.

\medterm{Ringworm} A contagious fungal infection of the skin, scalp, hair, or nails. It commonly causes an itchy, scaly rash with a ring-like edge; scalp infection may cause broken hairs or hair loss. The appearance and complications vary by the affected site.
\synonyms
Tinea Corporis; Dermatophytosis; Tinea

\medterm{Ringworm Of The Body} A contagious fungal infection of body skin causing an itchy, scaly, ring-shaped rash. It spreads through contact with infected people, animals, or objects and usually responds to antifungal treatment.

\medterm{Ringworm Of The Foot} Alternate name; see \textit{Athlete's foot}.

\medterm{Ringworm Of The Scalp} A fungal infection of the scalp and hair that can cause scaling, broken hairs, and hair loss.

\medterm{Rinne and Weber Tests} Complementary bedside tuning fork examinations (typically using a 512 Hz fork) to evaluate auditory hearing loss: Rinne compares air conduction to bone conduction in each ear, while Weber tests lateralization of sound when the vibrating fork is placed on the midline forehead.

\medterm{Rinnes Test} Rinne's test compares air and bone conduction with a tuning fork to help assess conductive hearing loss.

\medterm{Rinne Test} The Rinne test is a screening test that compares air conduction and bone conduction
hearing using a vibrating tuning fork (typically 512 Hz) to help differentiate between
conductive and sensorineural hearing loss.

\medterm{RIPE Therapy} Standard 4-drug regimen for active TB: Rifampin, Isoniazid, Pyrazinamide, Ethambutol.
Intensive phase 2 months, continuation phase 4 months (total 6 months). Pyridoxine
prevents isoniazid-induced neuropathy.

\medterm{Rippling Muscle Disease} A rare, benign myopathy characterized by mechanically-induced muscle contractions that
spread across the muscle as a visible wave (rippling), typically triggered by stretch,
percussion, or movement.

\medterm{Risedronate Sodium} The sodium salt of risedronate, a bisphosphonate used to prevent or treat osteoporosis and other disorders of excessive bone resorption.

\medterm{Risk} Risk is the probability of an unwanted event or outcome, considered together with its severity and the circumstances that influence it.

\medterm{Risk Adjustment} Risk adjustment accounts for differences in baseline patient characteristics so outcomes or payments can be compared more fairly across groups.

\medterm{Risk Factor} A characteristic, exposure, or behavior linked with a higher chance of disease or another outcome. It may be changeable, such as smoking, or unchangeable, such as age or inherited traits.

\medterm{Risk Management} Finding possible sources of harm, judging how likely and serious they are, and taking steps to reduce them. In healthcare, it includes safe systems, reporting problems, learning from errors, and checking whether safety measures work.

\medterm{Risk Register} A documented list of identified risks, their likelihood, impact, controls, owners, and review status.

\medterm{Risks Of Hip And Knee Replacement} Possible problems after joint replacement include infection, blood clots, bleeding, nerve or vessel injury, dislocation, stiffness, loosening, wear, persistent pain, and need for repeat surgery. Risk varies with the operation and the person's health.

\medterm{Risks Of Tobacco} Tobacco use increases risks of cancer, heart and blood-vessel disease, stroke, lung disease, pregnancy complications, addiction, and harm from secondhand smoke.

\medterm{Risks Of Underage Drinking} Alcohol use by minors can impair brain development, judgment, learning, safety, and mental health and increases risks of injury, dependence, and unsafe sexual or substance-related behavior.

\medterm{Risk-Taking} Risk-taking is behavior involving deliberate or inadvertent exposure to possible harm, influenced by perception, reward, development, mental health, and circumstances.

\medterm{Risky Drinking} Alcohol use that increases the chance of health, safety, work, relationship, or social problems. Risk depends on amount, frequency, pattern, personal health, medicines, pregnancy, age, and circumstances such as driving.

\synonyms
Hazardous Drinking; At-Risk Alcohol Consumption

\medterm{Risperidone} Risperidone is an atypical antipsychotic that blocks dopamine and serotonin receptors and treats schizophrenia, bipolar disorder, and selected severe behavioral symptoms; hyperprolactinaemia, metabolic effects, movement disorders, and sedation may occur.

\medterm{Ristocetin-Induced Platelet Aggregation} A blood test that checks how platelets clump when exposed to ristocetin, a substance that requires von Willebrand factor and a platelet surface protein. It helps investigate von Willebrand disease and some platelet disorders.

\medterm{Risus Sardonicus} A fixed, strained grin caused by painful facial-muscle spasms. It is classically associated with tetanus or strychnine poisoning and reflects severe muscle overactivity.

\medterm{Ritonavir} Ritonavir is an HIV protease inhibitor mainly used at low doses to boost levels of other antiviral medicines by slowing their breakdown.

\medterm{Ritual} A ritual is a repeated, structured behavior or ceremony; repetitive rituals may also occur as compulsions in obsessive-compulsive disorder.

\medterm{Rituximab} Rituximab is a monoclonal antibody against CD20 that depletes B cells and is used for B-cell malignancies and selected autoimmune diseases; infusion reactions, infections, hepatitis-B reactivation, and rare PML are important risks.

\medterm{Rivaroxaban} Rivaroxaban is an oral direct factor Xa inhibitor used to prevent or treat venous thromboembolism and reduce stroke risk in selected atrial fibrillation; bleeding and renal-function-dependent dosing are important considerations.

\medterm{Rivastigmine} A medicine that increases acetylcholine signaling and may temporarily improve thinking or daily function in Alzheimer disease or Parkinson disease dementia. Common effects include nausea, vomiting, diarrhea, appetite loss, dizziness, and slow heartbeat.

\medterm{River Blindness} Onchocerciasis, a parasitic infection caused by Onchocerca volvulus and transmitted by blackflies. It can
cause severe itching, skin disease, and visual impairment.

\medterm{RâLe} An older term for an abnormal crackling or bubbling respiratory sound, now generally described as a crackle.

\medterm{RLS} Restless legs syndrome, an urge to move the legs with uncomfortable sensations, especially at night.
\synonyms
Restless Legs Syndrome; Willis-Ekbom Disease

\medterm{RMDs} Abbreviation; see \textit{Repetitive Motion Disorders}.

\medterm{RNA} Ribonucleic acid, a molecule that helps cells read genetic instructions and make proteins. Some RNA carries messages from DNA, some helps build proteins, and some controls cell activity or forms part of viruses.

\medterm{RNA Binding} The attachment of a protein or other molecule to RNA, affecting RNA stability, transport, folding, processing, or translation.

\medterm{RNA Biosynthesis} The cellular production of RNA from a DNA template or, for some viruses, from an RNA template.

\medterm{RNA Degradation} The enzyme-controlled breakdown of RNA molecules. It controls how long genetic messages last, removes damaged or unnecessary messages, recycles their building blocks, and helps regulate protein production and cell behavior.

\medterm{RNA Folding} The process by which an RNA strand forms secondary and three-dimensional structures that influence its function.

\medterm{RNA Helicase} An enzyme that uses energy to separate, reshape, or move RNA structures and RNA-protein complexes. It helps cells copy, process, transport, and destroy RNA and is important for some viruses.

\medterm{RNA Interference} A gene-silencing process in which small RNAs guide proteins to degrade or block specific messenger RNA molecules.

\medterm{RNase P} An RNA-protein enzyme complex that processes precursor transfer RNA by removing its extra 5-prime sequence.

\medterm{RNA Splicing} Removal of introns and joining of exons from a precursor RNA to produce mature messenger RNA or other functional RNA.

\medterm{RNA Transport} Movement of RNA within a cell or between cells, often mediated by transport proteins and required for translation or signaling.

\medterm{RNA Virus} A virus whose genetic material is RNA; its replication often involves an RNA-dependent polymerase and may have a high mutation rate.

\medterm{Road Rage} Alternate name; see \textit{Intermittent explosive disorder}.

\medterm{Robinow Syndrome} A rare inherited skeletal-development disorder with distinctive facial features, short stature, vertebral or rib abnormalities, and genital differences.

\medterm{Robot-Assisted Rehabilitation} Robotic devices for repetitive task training. Upper extremity, lower extremity, gait
training. Objective measurement, high repetition.

\medterm{Robotic-Assisted Thoracoscopic Surgery} Chest surgery performed through several small openings using instruments controlled by a surgeon from a console. The camera provides a magnified three-dimensional view and the instruments can move precisely. It may reduce the size of the incisions, but bleeding, infection, air leakage, and conversion to open surgery remain possible.

\synonyms
RATS; Robotic Thoracic Surgery; Robotic Lobectomy

\medterm{Robotic Hysterectomy} Hysterectomy performed through minimally invasive abdominal access using a surgeon-controlled robotic platform. The extent of uterine, cervical, tubal, and ovarian removal depends on the indication and operative plan.

\medterm{Robotic Surgery} A surgical technique in which a surgeon controls a computer-assisted system that moves a camera and small instruments. It is usually performed through small incisions and may allow precise movements, but the system does not operate independently.

\synonyms
Robot-Assisted Surgery; Computer-Assisted Surgery; RAS

\medterm{Rocker-Bottom Foot} A congenital foot deformity characterized by a prominent, downward-convex plantar surface with a dorsiflexed midfoot and rigid equinovarus hindfoot (prominent calcaneus), classically caused by a congenital vertical talus and strongly associated with trisomy 18 (Edwards syndrome) and trisomy 13 (Patau syndrome).

\medterm{Rocky Mountain Spotted Fever} A severe, potentially fatal tick-borne rickettsial infection caused by \textit{Rickettsia rickettsii} transmitted by \textit{Dermacentor} ticks, presenting with high fever, severe headache, myalgias, and a petechial maculopapular rash that characteristically begins on the wrists and ankles and spreads centripetally to the trunk, palms, and soles.

\synonyms
RMSF

\medterm{Rod} A rod is a long cylindrical structure; in the retina, a rod photoreceptor supports dim-light and peripheral vision.

\medterm{Rodenticide Toxicity} A heterogeneous group of poisonings whose presentation depends on the active ingredient. Anticoagulant products cause delayed coagulopathy and bleeding; bromethalin causes neurologic toxicity, cholecalciferol causes hypercalcemia and kidney injury, and metal phosphides can cause shock and organ failure. Identify the product and provide agent-specific treatment.

\medterm{Rodent Ulcer} An older name for basal cell carcinoma, a common locally invasive skin cancer.

\medterm{Rod Photoreceptor} A light-sensitive cell in the retina that supports vision in dim light and the peripheral field. Rods do not provide color vision; their loss can cause night blindness and narrowing of the visual field.

\medterm{Rods} Retinal photoreceptor cells that provide vision in dim light and contribute to peripheral vision but do not detect color.

\medterm{Rods And Cones} Two types of light-sensing cells in the retina. Rods help with dim-light and side vision, while cones provide sharp and color vision in brighter light.

\medterm{Roentgen Rays} Alternate name; see \textit{X-Ray}.

\synonyms
Roentgen Radiation; Electromagnetic X-Radiation

\medterm{Rofecoxib} Rofecoxib is a selective COX-2 anti-inflammatory drug withdrawn in many countries because of increased cardiovascular risk.

\medterm{Romano-Ward Syndrome} The most common congenital long QT syndrome, inherited in an autosomal dominant pattern without associated neurosensory deafness (due to mutations in cardiac potassium or sodium channel genes \textit{KCNQ1}, \textit{KCNH2}, or \textit{SCN5A}), predisposing to syncope, polymorphous ventricular tachycardia (torsades de pointes), and sudden cardiac death.

\medterm{Romanus Lesion} An early inflammatory radiographic finding in ankylosing spondylitis characterized by enthesitis with focal erosion and surrounding reactive sclerosis at the anterior superior and inferior corners of vertebral bodies ('shiny corners'), preceding syndesmophyte formation.
\synonyms
Shiny Corner Sign; Spondylitis Anterior Enthesitis

\medterm{Rombergism} Unsteadiness or falling when standing with feet together and eyes closed, indicating impaired proprioception or vestibular function.

\medterm{Romberg Sign} Loss of balance when standing with the feet together and the eyes closed, suggesting impaired proprioception or vestibular function.

\medterm{Romberg's Sign} Increased unsteadiness when a person stands with the feet together and closes the eyes, suggesting impaired
proprioception or vestibular function.

\medterm{Romberg Test} The Romberg test is a clinical neurological test used to assess the integrity of the
dorsal column-medial lemniscus pathway (proprioception) and the cerebellum in
maintaining balance. The test is performed by asking the patient to stand with feet
together and eyes open, and then to close the eyes.

\medterm{Romidepsin} A histone deacetylase inhibitor used for selected T-cell lymphomas; it can cause cytopenias, infection risk, and ECG-related effects.

\medterm{Romiplostim} A thrombopoietin-receptor agonist that stimulates platelet production and is used mainly for chronic immune thrombocytopenia.

\medterm{Rongeur} A strong, jawed surgical instrument used to bite and remove small pieces of bone or other hard tissue.

\medterm{Ronnel} Ronnel is an organophosphate insecticide that can inhibit acetylcholinesterase and cause cholinergic toxicity after exposure.

\medterm{Root Canal} The channel inside a tooth root containing the pulp and its nerves and blood vessels.

\medterm{Root Canal Therapy} Dental treatment in which infected or inflamed pulp is removed from a tooth, the canal is cleaned and sealed, and the tooth is restored.

\medterm{Root Canal Treatment} Alternate term; see \textit{Root Canal Therapy}.

\medterm{Root Caries} Tooth decay affecting exposed root surfaces after gum recession or loss of protective covering. It is more common with aging, dry mouth, or limited cleaning and may be slowed with fluoride and dental care.

\medterm{Root Filling} Dental treatment sealing a cleaned tooth root canal to prevent bacteria entering again and causing reinfection.

\medterm{Root Fracture} A crack or break in the root of a tooth. A horizontal fracture may sometimes be stabilized and monitored, while a vertical fracture often allows bacteria to enter and may require removal of the tooth. Treatment depends on the fracture’s location and extent.

\medterm{Rooting Reflex} A primitive reflex present from birth to 3-4 months, triggered by stroking the infant's
cheek or corner of the mouth. The infant turns toward the stimulus and opens the mouth,
searching for the nipple. Facilitates breastfeeding. Absence may indicate neurologic
impairment. Disappears as voluntary feeding develops.

\medterm{Ropinirole} A medicine that stimulates dopamine receptors and is used for Parkinson disease and restless legs syndrome. It may cause nausea, dizziness, sleepiness, low blood pressure, hallucinations, swelling, or impulse-control problems such as gambling.

\medterm{Ropivacaine} A long-acting local anesthetic used for regional nerve blocks and epidural anesthesia; excessive blood levels can cause neurologic or cardiac toxicity.

\medterm{Ropivacaine Hydrochloride} The hydrochloride salt of ropivacaine, a long-acting local anesthetic used for regional and epidural anesthesia.

\medterm{Rorschach Test} A projective psychological test using inkblots; its diagnostic validity is limited and interpretation requires specialist training.

\medterm{Rosacea} A chronic inflammatory skin disorder that most often affects the central face. It may cause recurrent flushing, persistent redness, visible small blood vessels, burning or stinging, and acne-like bumps or pustules without being acne. Eye involvement may cause irritation or inflammation, and long-standing disease may produce thickening of the skin, especially of the nose (rhinophyma). Rosacea is not contagious, and its severity and features vary over time.

\medterm{Rosai-Dorfman Disease} A rare, benign, non-Langerhans cell histiocytic proliferative disorder (sinus histiocytosis with massive lymphadenopathy) characterized clinically by massive painless bilateral cervical lymphadenopathy and fever, and pathologically by sinusoidal infiltration of S100-positive histiocytes demonstrating emperipolesis.
\synonyms
Sinus Histiocytosis with Massive Lymphadenopathy; SHML

\medterm{Rose Bengal} Rose bengal is a dye used in eye-surface testing and laboratory or diagnostic procedures; it highlights damaged or poorly protected conjunctival cells.

\medterm{Roseola} Roseola infantum is a common viral illness, usually caused by human herpesvirus 6, with several days of high fever followed by a pink trunk rash as the fever resolves; febrile seizures may occur.

\synonyms
Exanthema Subitum; Sixth Disease; Roseola Infantum

\medterm{Roseola Infantum} A common childhood illness, usually caused by human herpesvirus 6 and sometimes HHV-7. Several days of high fever are followed by a pink or red trunk rash as the fever resolves; febrile seizures may occur, and illness is usually self-limited.

\medterm{Roseolovirus} A genus of herpesviruses that includes human herpesvirus 6 and 7, commonly associated with roseola and sometimes serious disease after reactivation.

\medterm{Rosiglitazone} A medicine for type 2 diabetes that helps the body respond better to insulin and lower blood sugar. It can cause weight gain, swelling, fractures, or worsening heart failure and is used selectively.

\medterm{Rosiglitazone Maleate} The maleate salt of rosiglitazone, a thiazolidinedione that improves insulin sensitivity in type 2 diabetes but may cause weight gain, edema, fractures, and heart-failure worsening.

\medterm{Ross Procedure} An operation for severe aortic-valve disease, especially in some children and young adults. The patient’s own pulmonary valve is moved to the aortic position, and a donor valve or tube replaces it. The moved valve can grow in children and usually avoids lifelong blood-thinning medicine, but either valve may later narrow or leak.

\synonyms
Pulmonary Autograft Procedure; Ross Operation

\medterm{Ross River Virus} A mosquito-borne alphavirus that causes Ross River fever, with fever, rash, and often prolonged joint pain or swelling.

\medterm{Rostral} Positioned toward the nose or front of the head. In descriptions of the brain, it usually means toward the forehead or front rather than toward the back of the head.

\medterm{Rosuvastatin Calcium} Rosuvastatin calcium is a statin medicine that lowers LDL cholesterol and reduces the risk of cardiovascular events.

\medterm{Rotation} Movement around an axis; in anatomy it describes turning a body part inward, outward, or around its longitudinal axis.

\medterm{Rotator Cuff} The group of four shoulder muscles and tendons that stabilize the glenohumeral joint and rotate the arm.

\synonyms
Musculotendinous Cuff; SITS Muscles

\medterm{Rotator Cuff Arthropathy} Shoulder-joint damage that develops after long-standing, irreparable tears of the rotator-cuff tendons. It can cause pain, weakness, grinding, and inability to raise the arm; severe cases may need specialized shoulder replacement.

\medterm{Rotator Cuff Exercises} Rotator cuff exercises strengthen and stretch shoulder muscles and are selected according to the injury, pain, range of motion, and rehabilitation plan.

\medterm{Rotator Cuff Injury} Damage to one or more shoulder tendons or muscles that stabilize and move the arm. It may cause pain, weakness, night discomfort, or limited movement after overuse, aging, or sudden injury.

\medterm{Rotator Cuff Problems} Rotator cuff problems include tendinopathy, bursitis, and tears that cause shoulder pain, weakness, painful motion, and difficulty with overhead activity.

\medterm{Rotator Cuff Repair} Surgery to reattach a torn shoulder tendon to the upper-arm bone, usually with small instruments or an open approach. Recovery requires structured rehabilitation and may take months; stiffness, weakness, recurrent tearing, or infection can occur.

\medterm{Rotator Cuff Tear} A partial or complete tear of one or more tendons stabilizing the shoulder, causing pain, weakness, and difficulty raising or rotating the arm.

\medterm{Rotator Cuff Tendinitis} Alternate name; see \textit{Rotator cuff syndrome}.

\medterm{Rotavirus} A very contagious virus that commonly causes vomiting and watery diarrhea in infants and young children. It can cause dangerous dehydration; vaccination greatly reduces severe illness, and testing is not always needed.

\medterm{Rotavirus Antigen Test} A rotavirus antigen test detects rotavirus proteins in stool and can help diagnose viral gastroenteritis, especially in infants and young children.

\medterm{Rotaviruses} A group of viruses causing vomiting and watery diarrhea, especially in young children; vaccination lowers severe disease risk.

\medterm{Rotavirus Infection} A highly contagious viral gastroenteritis causing vomiting and watery diarrhea, especially in infants and young children; vaccination prevents most severe disease.

\medterm{Rotavirus Vaccine} A live-attenuated oral viral vaccine administered to young infants in a 2- or 3-dose series to prevent severe dehydrating gastroenteritis and diarrhea caused by rotavirus infection, carrying a slight, transient risk of intussusception.

\medterm{Rotenone} Rotenone is a naturally occurring pesticide that inhibits mitochondrial complex I and can be toxic to people and animals after exposure.

\medterm{Rothmund-Thomson Syndrome} A rare inherited condition that causes a mottled, red, and thin skin rash beginning in infancy, short stature, sparse hair, skeletal or dental problems, and sometimes cataracts. It also increases the risk of certain cancers.

\medterm{Roth Spots} Small areas of bleeding in the retina with pale centers, seen during an eye examination. They may occur with infection of the heart valves, severe anemia, leukemia, diabetes, or other illnesses and are not specific to one diagnosis.

\medterm{Roth's Spots} Small retinal haemorrhages with pale centres, associated with conditions such as infective endocarditis, severe anaemia, leukaemia, and diabetes.

\medterm{Rotor Syndrome} A benign, autosomal recessive inherited hyperbilirubinemia caused by mutations in SLCO1B1 and SLCO1B3 genes, resulting in defective hepatic reuptake and storage of conjugated bilirubin, presenting with lifelong, asymptomatic conjugated hyperbilirubinemia with normal liver histology and elevated urinary coproporphyrin.
\synonyms
Rotor's Syndrome; Familial Conjugated Hyperbilirubinemia

\medterm{Rotterdam Criteria} A commonly used approach for diagnosing polycystic ovary syndrome in adults. After other causes are excluded, diagnosis generally requires at least two of: irregular or absent ovulation, signs or blood-test evidence of high androgen levels, or polycystic-appearing ovaries on ultrasound. Ultrasound alone does not establish the diagnosis.

\medterm{Rotter Nodes} Small interpectoral lymph nodes situated between the pectoralis major and pectoralis minor muscles that provide an alternate pathway for breast lymphatic drainage directly to Level II and III axillary lymph nodes.

\synonyms
Interpectoral Lymph Nodes; Rotter's Nodes

\medterm{Rouleau} A stack or column of red blood cells, commonly seen in blood with increased plasma
proteins.

\medterm{Rouleaux} A pattern where red blood cells stack like coins, often because blood plasma contains extra proteins.

\medterm{Rouleaux Formation} Linear, coin-like stacks or rolls of red blood cells visualized on a peripheral blood smear, occurring when elevated plasma concentrations of asymmetric acute-phase proteins (fibrinogen or monoclonal immunoglobulins) neutralize the normal negative repulsive zeta potential of erythrocytes, characteristically seen in multiple myeloma and severe inflammatory states.

\medterm{Round} Round means circular or approximately spherical; in anatomy or pathology it describes shape and must be interpreted with the structure being described.

\medterm{Round Back} Alternate name; see \textit{Kyphosis}.

\medterm{Round Ligament} A round ligament is a cord-like supporting structure, including the uterine ligaments that pass toward the groin and help maintain uterine position.

\medterm{Round Ligament of the Uterus} A fibrous and muscular cord originating from the superior-lateral aspect of the uterine cornu anterior to the fallopian tube, traversing through the inguinal canal to terminate in the subcutaneous tissue of the labium majus, functioning to hold the uterus in an anteverted and anteflexed position.

\medterm{Round Ligament Pain} Sharp or aching pain in the lower abdomen or groin caused by stretching of the round ligaments as the uterus enlarges during pregnancy. Persistent, severe, or bleeding-associated pain requires medical evaluation.

\medterm{Round Pneumonia} A lung infection that appears as a rounded shadow on a chest X-ray rather than the usual patchy pattern. It occurs especially in children and can resemble a tumor, so follow-up imaging may be needed after treatment.

\medterm{Roundworm} A parasitic worm, often referring to Ascaris lumbricoides, acquired by swallowing eggs from contaminated soil, food, or water. It may cause abdominal symptoms, poor nutrition, or intestinal blockage when many worms are present.

\medterm{Roundworms} Nematode parasites that may infect intestines or tissues, causing disease depending on their species and number.

\medterm{Roussy-Levy Syndrome} An autosomal dominant hereditary motor and sensory neuropathy (a phenotypic variant of Charcot-Marie-Tooth disease type 1A) characterized by early-onset sensory ataxia, prominent postural hand tremor, distal muscle weakness and wasting, pes cavus, and absent deep tendon reflexes.
\synonyms
Hereditary Areflexic Dystasia; CMT-Tremor Variant

\medterm{Route Of Administration} The way a medicine enters the body, such as by mouth, under the tongue, injection, inhalation, skin application, or through a vein. The route affects how quickly and how much of the medicine reaches the bloodstream.

\medterm{Roux-en-Y Gastric Bypass} A bariatric and metabolic surgical procedure in which a small (15 to 30 mL) upper gastric pouch is created and anastomosed to a Roux limb of proximal jejunum (gastrojejunostomy), bypassing the remaining stomach, duodenum, and proximal jejunum, connected to a downstream biliopancreatic limb (jejunojejunostomy).

\synonyms
RYGB

\medterm{Roux Stasis Syndrome} A chronic functional disorder of gastric and small bowel motility developing in up to 30 percent of patients following a Roux-en-Y gastrojejunostomy or gastrectomy. Characterized by delayed gastric emptying, severe postprandial fullness, nausea, vomiting of non-bilious food, and chronic epigastric pain, resulting from the loss of normal duodenal pacemaking coordination in the transected Roux jejunal limb.

\synonyms
Roux-en-Y Stasis Syndrome; Chronic Roux Limb Stasis

\medterm{Rovsing Sign} Pain in the lower right abdomen that occurs when the lower left abdomen is pressed. This examination finding can support suspicion of appendicitis, but it is not diagnostic by itself and must be interpreted with the symptoms and other tests.

\medterm{Roxithromycin} A macrolide antibiotic used in some countries for selected bacterial infections. It stops bacteria making proteins, but resistance, medicine interactions, stomach upset, liver problems, and heart-rhythm changes can occur.

\medterm{RPR Test} A blood screening test for syphilis that detects antibodies produced during tissue injury from the infection. A reactive result does not prove active syphilis and requires a second, more specific test and clinical assessment.

\medterm{RS3PE Syndrome} Remitting seronegative symmetrical synovitis with pitting edema, an inflammatory syndrome causing sudden swelling of the hands or feet, joint pain, and stiffness, usually in older adults. It may occur with another rheumatic disease, infection, or cancer.
\synonyms
Remitting Seronegative Symmetrical Synovitis with Pitting Edema

\medterm{RSI} Repetitive strain injury, a group of pain and dysfunction syndromes associated with repeated movement or sustained posture.

\medterm{RSV} A common respiratory virus that can cause mild cold symptoms or serious bronchiolitis and pneumonia, especially in infants and older adults.

\medterm{RSV Antibody Test} An RSV antibody test detects antibodies to respiratory syncytial virus and may help assess past exposure or immune response, but is not usually the main test for acute infection.

\medterm{RSV Infection} Alternate name; see \textit{Respiratory Syncytial Virus Infection}.

\medterm{RT-PCR} A laboratory test that converts RNA into DNA and then makes many copies of a selected genetic sequence so it can be detected. It can identify RNA viruses and measure selected genes in clinical samples.

\medterm{Rubber} Rubber is an elastic material used in gloves, devices, and equipment; natural-rubber latex can cause irritant or allergic reactions.

\medterm{Rubber Band Ligation} The most common and effective non-surgical outpatient procedure performed to treat symptomatic internal hemorrhoids. Through a small lighted viewing scope (an anoscope) inserted into the anal canal, a clinician uses a suction or mechanical ligator instrument to grasp the redundant hemorrhoidal tissue above the pain-sensitive dentate line and snap a tiny, tight elastic rubber band firmly around its base. The tight band strangles the hemorrhoid's blood supply; within 7 to 10 days, the ischemic hemorrhoidal tissue withers, dies, and painlessly detaches, sloughing off and passing out unnoticed during a bowel movement. The remaining small scar fixes the surrounding mucosa to the underlying tissue, preventing further prolapse or bleeding without requiring surgical incisions, general anesthesia, or hospital recovery.

\synonyms
Hemorrhoidal Banding; RBL; Band Ligation of Hemorrhoids; Elastic Band Ligation

\medterm{Rubber Cement Poisoning} Poisoning from swallowing or inhaling rubber cement, which contains volatile solvents. It may cause dizziness, drowsiness, confusion, breathing problems, or chemical pneumonitis if it enters the lungs.

\medterm{Rubber Dams} Sheets used in dentistry to isolate teeth from saliva and the rest of the mouth during procedures.

\medterm{Rubella} A contagious viral infection causing mild fever, swollen lymph nodes, and a fine rash. Infection during pregnancy can cause miscarriage, fetal growth problems, hearing loss, eye disease, congenital heart disease, or other congenital rubella syndrome; vaccination prevents most cases.

\synonyms
German Measles; Three-Day Measles

\medterm{Rubella Virus} The virus that causes rubella, a contagious illness with fever, swollen glands, and rash. Infection during pregnancy can seriously harm the developing baby, causing hearing loss, eye disease, heart problems, or other birth defects.

\medterm{Rubeola} Measles, a highly contagious viral illness caused by measles virus and characterized by fever, cough, conjunctivitis, and a widespread maculopapular rash.

\synonyms
Measles; Hard Measles; Morbilli

\medterm{Rubeosis Iridis} Growth of abnormal new blood vessels on the colored part of the eye. It usually results from poor blood flow in the retina, often caused by diabetes or a blocked retinal vein, and can lead to painful glaucoma and vision loss.

\medterm{Rubidium} A chemical element whose radioactive isotopes can be used in medical imaging; nonmedical rubidium compounds can affect potassium-related physiology.

\medterm{Rubin Maneuver} An emergency childbirth technique for shoulder dystocia, when a baby’s shoulder is stuck behind the birthing person’s pelvic bone. A clinician applies pressure and rotates the baby’s shoulder to reduce its width and help complete delivery safely.

\medterm{Rubinstein-Taybi Syndrome} A genetic developmental disorder associated with characteristic facial features,
short stature, broad thumbs and toes, and intellectual disability. It is commonly caused by pathogenic variants in
CREBBP or EP300.

\medterm{Rubivirus} A genus of viruses that includes rubella virus, which causes rubella and congenital rubella syndrome.

\medterm{Rubredoxins} Small proteins containing iron that help move electrons in some bacteria and other organisms. They take part in chemical reactions that release energy and protect cells from harmful oxygen-related chemicals.

\medterm{Rubricyte Count} A measurement of rubricytes, an older name for immature red blood cells, in blood or bone marrow. It may help assess how actively the marrow is making red cells, but modern tests usually report reticulocytes or other marrow findings.

\medterm{Rule Of Nines} A quick method for estimating the percentage of body surface burned by assigning simple proportions to body regions. The estimate helps guide fluids and emergency care, but children and unusual body shapes require more accurate methods.

\medterm{Rule Out} Rule out means evaluate a possibility and determine whether it is present; it does not mean the condition has already been excluded.

\medterm{Rumination} Repeated regurgitation and rechewing of food after eating, without the usual vomiting process. It is a behavior or symptom that may occur in rumination disorder or with other gastrointestinal, developmental, or behavioral conditions.

\medterm{Rumination Disorder} A feeding or eating disorder involving repeated, usually effortless return of recently eaten food to the mouth, followed by rechewing, reswallowing, or spitting. It differs from ordinary vomiting and causes clinically significant distress, impairment, or nutritional or medical consequences.

\medterm{Rumination Syndrome} Alternate name; see \textit{Rumination Disorder}.

\medterm{Ruminococcus Gnavus} An intestinal bacterium involved in mucin and carbohydrate metabolism and studied in relation to gut inflammation.

\medterm{Rumpel-Leede Sign} Small pinpoint bleeding spots that appear below a blood-pressure cuff after it has briefly been inflated. The sign may suggest fragile small blood vessels or a low platelet count, but it is not specific and is rarely used alone today.

\synonyms
Tourniquet Test; Capillary Fragility Test

\medterm{Rumpel-Leede Test} A clinical diagnostic test evaluating capillary wall fragility and thrombocytopenia in which a sphygmomanometer cuff is inflated around the arm midway between systolic and diastolic pressure for five minutes; the appearance of more than 10 to 20 petechiae in a designated circular area distal to the cuff indicates increased capillary fragility (as seen in dengue fever, thrombocytopenia, or scurvy).

\synonyms
Capillary Fragility Test; Tourniquet Test

\medterm{Runner's Knee} Alternate name; see \textit{Patellofemoral pain syndrome}.

\medterm{Runny Nose} Excess fluid or mucus draining from the nose, commonly caused by infection, allergy, irritants, or sinus disease.

\medterm{Rupture} A tear or complete break in an organ, vessel, muscle, or other tissue, often causing bleeding or loss of function.

\medterm{Ruptured Abdominal Aortic Aneurysm} A burst in a weakened section of the main artery in the abdomen, causing severe internal bleeding. Sudden abdominal or back pain, fainting, sweating, or collapse may occur.

\medterm{Ruptured Disk} A spinal disk whose outer layer has torn, allowing inner material to bulge outward.
\synonyms
Herniated Disc; Slipped Disc; Herniated Nucleus Pulposus

\medterm{Ruptured Eardrum} A ruptured eardrum is a tear in the tympanic membrane caused by infection, pressure, trauma, or foreign objects and may cause pain, drainage, hearing loss, or dizziness.

\medterm{Ruptured Perforated Eardrum} A perforated eardrum is a tear in the tympanic membrane from infection, pressure change, trauma, or an inserted object. Pain, drainage, hearing loss, or ringing may occur; keeping the ear dry and avoiding unprescribed drops helps while healing or repair is assessed.

\medterm{Russell-Silver Syndrome} A condition present from birth that causes poor growth before and after birth, a relatively large head, feeding difficulty, uneven growth of the limbs or body, and characteristic facial features. It is usually related to changes in gene regulation.

\medterm{Russell Viper Venom Time} A blood-clotting test used as part of testing for lupus anticoagulants. The venom activates clotting in the test tube; a prolonged result that shortens when extra phospholipid is added supports the presence of a lupus anticoagulant but does not establish the diagnosis by itself.

\synonyms
dRVVT

\medterm{Rx} A common abbreviation for prescription or prescribed treatment. It comes from the Latin word meaning “take” and is used on medication orders and in medical writing.

\medterm{Ryle's Tube} A soft tube passed through the nose into the stomach to remove stomach contents or give fluids, nutrition, or medicines. Position must be checked before use because a misplaced tube can enter the airway and cause serious harm.
